%============================================================
%  Bd. 05 - Funktionen
%============================================================

\documentclass[main.tex]{subfiles}

\ifSubfilesClassLoaded{
  \BandSetupStandalone{B05}
}{
}

\title{Bd. 05 - Funktionen}
\author{Martin Kunze}
\date{}
\setcounter{file}{5}

\hypersetup{
  pdftitle={Bd. 05 - Funktionen},
  pdfauthor={Martin Kunze}
}

\begin{document}

\title{Bd. 05 - Funktionen}
\author{Martin Kunze}
\date{}
\hypersetup{pageanchor=false}
\maketitle
\hypersetup{pageanchor=true}
\setcounter{tocdepth}{3}
\tableofcontents
\setcounter{file}{5}
\renewcommand{\FormulaBandID}{5}

\chapter{Einleitung}

Dieser Band entwickelt Funktionen als besondere totale Relationen auf der in
Band~03 aufgebauten Mengenlehre und der in Band~04 eingeführten Theorie
totaler Relationen. Er enthält die Begriffsbildung und alle Aussagen, die für
beliebige Funktionen gelten.

Die drei Eigenschaftstheorien sind bewusst ausgelagert: Band~06 behandelt
injektive, Band~07 surjektive und Band~08 bijektive Funktionen. Das
Auswahlprinzip und seine funktionalen Äquivalente bilden anschließend den
eigenständigen Band~09. Dadurch können spätere Beispiele genau dem Band
zugeordnet werden, dessen Abbildungseigenschaft sie besitzen.

In diesem Band verbleiben daher nur eigenschaftsneutrale Konstruktionen:
konstante Funktionen, Fall-, Faser-, Einschränkungs-, Kompositions-,
Potenzmengen-, Spur- und Erweiterungsabbildungen. Ihre Injektivität,
Surjektivität oder Bijektivität hängt jeweils von zusätzlichen Voraussetzungen
ab. Dagegen werden Inklusionsabbildungen erstmals in Band~06, Projektionen
erstmals in Band~07 sowie Identitäts- und Umkehrfunktionen erstmals in Band~08
eingeführt.

\chapter{Allgemeine Theorie der Funktionen}

\section{Begriff der Funktion}

\subsection{Axiom der funktionalen Eindeutigkeit}

\FormulaAxiomDelta[Funktionale Eindeutigkeit]
{(x,y)\in F\dsep (x,z)\in F \vdash y=z}
{
\DeltaRow{Mengen}{x\dsep y\dsep z\dsep F}
}

\subsection{Definition der Funktion}


\FormulaDefDeltaK[Begriff der Funktion]{F\colon A\to B}{Funktion}{
  \DeltaRow{Mengen}{A\dsep B\dsep F\dsep x\dsep y\dsep z}
  \DeltaRow{\textbf{Axiome}}{}
  %
  % — Existenz als Menge geordneter Paare —
  \DeltaRow{\makecell[l]{Menge geordneter\\Paare}}
           {F \subseteq A \times B}
           [\FormulaRefAuto{F \subseteq A \times B}]
  %
  % — Funktionale Eindeutigkeit —
  \DeltaRow{\makecell[l]{Funktionale \\Eindeutigkeit}}
           {(x,y)\in F\dsep (x,z)\in F \vdash y=z}
           [\FormulaRefAuto{(x,y)\in F\dsep (x,z)\in F \vdash y=z}]
  %
  % — Aussonderung (Schema) —
  \DeltaRow{Totalität}
           {x\in A \vdash \exists y\,(x,y)\in F}
           [\FormulaRefAuto{x\in A \vdash \exists y\,(x,y)\in F}]
%
    \DeltaRow{\textbf{Neue Symbole}}{}
    \DeltaPrem{Funktionen}{ }
}

\FormulaThmDeltaK[Der Funktionsgraph liegt im kartesischen Produkt]{%
  F\colon A\to B\vdash F\subseteq A\times B%
}{FunctionGraphSubsetProduct}{%
  \DeltaRow{Mengen}{A\dsep B\dsep F}
}
\begin{tabproof}
  \proofstep{1}{F\colon A\to B}{\rA}
  \proofstep{1}{F\subseteq A\times B}{%
    \FormulaRefAuto{F\subseteq A\times B}{1}}
\end{tabproof}

\section{Grundlegende Eigenschaften}

\FormulaThmDelta{F\colon A\to B \vdash \TotRel{F,A,B}}{
\DeltaRow{Mengen}{A\dsep B\dsep F}
}
\begin{tabproof}
    \proofstep{1}{F\colon A\to B}{\rA}
    \proofstep{1}{F \subseteq A \times B}{\FormulaRefAuto{R \subseteq A \times B}{1}}
    \proofstep{3}{x \in A}{\rA}
    \proofstep{1,3}{\exists y\,\bigl((x,y)\in F\bigr)}{%
      \FormulaRefAuto{x\in A \vdash \exists y\,(x,y)\in F}{1,3}}
    \proofstep{1}{\TotRel{F,A,B}}{\FormulaRefAuto{Totale Relation}{2,4}}
\end{tabproof}


\FormulaThmDelta{F\colon A\to B\dsep (x,y)\in F\vdash (x,y)\in A\times B}{
\DeltaRow{Mengen}{x\dsep y\dsep A\dsep B}
\DeltaRow{Funktionen}{F}
}
\begin{tabproof}
  \proofstep{1}{F\colon A\to B}{\rA}
  \proofstep{2}{(x,y)\in F}{\rA}
  \proofstep{1}{F\subseteq A\times B}{\FormulaRefAuto{F \subseteq A \times B}{1}}
  \proofstep{1,2}{(x,y)\in A\times B}{\FormulaRefAuto{A\subseteq B,\, x\in A \vdash x\in B}{3,2}}
\end{tabproof}

\FormulaThmDelta{F\colon A\to B\dsep (x,y)\in F\vdash x\in A}{
\DeltaRow{Mengen}{x\dsep y\dsep A\dsep B}
\DeltaRow{Funktionen}{F}
}
\begin{tabproof}
  \proofstep{1}{F\colon A\to B}{\rA}
  \proofstep{2}{(x,y)\in F}{\rA}
  \proofstep{1,2}{(x,y)\in A\times B}{\FormulaRefAuto{F\colon A\to B\dsep (x,y)\in F\vdash (x,y)\in A\times B}{1,2}}
  \proofstep{1,2}{x\in A}{\FormulaRefAuto{(a,b)\in A\times B\vdash a\in A}{3}}
\end{tabproof}

\FormulaThmDelta{F\colon A\to B\dsep (x,y)\in F\vdash y\in B}{
\DeltaRow{Mengen}{x\dsep y\dsep A\dsep B}
\DeltaRow{Funktionen}{F}
}
\begin{tabproof}
  \proofstep{1}{F\colon A\to B}{\rA}
  \proofstep{2}{(x,y)\in F}{\rA}
  \proofstep{1,2}{(x,y)\in A\times B}{\FormulaRefAuto{F\colon A\to B\dsep (x,y)\in F\vdash (x,y)\in A\times B}{1,2}}
  \proofstep{1,2}{y\in B}{\FormulaRefAuto{(a,b)\in A\times B\vdash b\in B}{3}}
\end{tabproof}

\subsection{Eindeutigkeit des Definitionsbereichs}

Der Zielbereich einer Funktion darf vergr\"o\ss ert werden, ohne ihren Graphen
zu ver\"andern. Ihr Definitionsbereich ist dagegen bereits durch den Graphen
festgelegt.

\FormulaThmDeltaK[Eindeutigkeit des Definitionsbereichs einer Funktion]{%
  F\colon A\to B\dsep F\colon C\to D\vdash A=C%
}{FunctionDomainUnique}{%
  \DeltaRow{Mengen}{A\dsep B\dsep C\dsep D\dsep F\dsep x\dsep y}
}
\begin{tabproof}
  \proofstep{1}{F\colon A\to B}{\rA}
  \proofstep{2}{F\colon C\to D}{\rA}
  \proofstep{3}{x\in A}{\rA}
  \proofstep{1,3}{\exists y\,((x,y)\in F)}{%
    \FormulaRefAuto{x\in A\vdash \exists y\,(x,y)\in F}{1,3}}
  \proofstep{5}{(x,y)\in F}{\rA}
  \proofstep{2,5}{x\in C}{%
    \FormulaRefAuto{F\colon A\to B\dsep (x,y)\in F\vdash x\in A}{2,5}}
  \proofstep{1,2,3}{x\in C}{\rEE{4,5,6}}
  \proofstep{1,2}{\forall x\,(x\in A\rightarrow x\in C)}{%
    \rUI{\rRI{3,7}}}
  \proofstep{1,2}{A\subseteq C}{%
    \FormulaRefAuto{A\subseteq B\coloneq
      \forall x\,(x\in A\rightarrow x\in B)}[def]{8}}
  \proofstep{10}{x\in C}{\rA}
  \proofstep{2,10}{\exists y\,((x,y)\in F)}{%
    \FormulaRefAuto{x\in A\vdash \exists y\,(x,y)\in F}{2,10}}
  \proofstep{12}{(x,y)\in F}{\rA}
  \proofstep{1,12}{x\in A}{%
    \FormulaRefAuto{F\colon A\to B\dsep (x,y)\in F\vdash x\in A}{1,12}}
  \proofstep{1,2,10}{x\in A}{\rEE{11,12,13}}
  \proofstep{1,2}{\forall x\,(x\in C\rightarrow x\in A)}{%
    \rUI{\rRI{10,14}}}
  \proofstep{1,2}{C\subseteq A}{%
    \FormulaRefAuto{A\subseteq B\coloneq
      \forall x\,(x\in A\rightarrow x\in B)}[def]{15}}
  \proofstep{1,2}{A=C}{%
    \FormulaRefAuto{A\subseteq B\dsep B\subseteq A\vdash A=B}[thm]{9,16}}
\end{tabproof}

\section{Funktionswert}

\FormulaThmDelta[Eindeutigkeit des Funktionswertes]{F\colon A\to B\dsep x\in A\vdash \exists! y\;(x,y)\in F}{
\DeltaRow{Mengen}{x\dsep A\dsep B}
\DeltaRow{Funktionen}{F}
}
\begin{tabproof}
  \proofstep{1}{F\colon A\to B}{\rA}
  \proofstep{2}{x\in A}{\rA}
  \proofstep{1,2}{\exists y(x,y)\in F}{\FormulaRefAuto{x\in A \vdash \exists y\,(x,y)\in F}{1,2}}
  \proofstep{4}{(x,y)\in F}{\rA}
  \proofstep{5}{(x,z)\in F}{\rA}
  \proofstep{4,5}{y=z}{\FormulaRefAuto{(x,y)\in F\dsep (x,z)\in F \vdash y=z}{4,5}}
  \proofstep{1,2}{\exists! y(x,y)\in F}{\UEI{3,4,5,6}}
\end{tabproof}
%%begin novalidate
\FormulaDefDelta[Funktionswert]{F\colon A\to B\dsep x\in A\vdash (F(x) \coloneqq \iota y\,\bigl((x,y)\in F\bigr)}{
\DeltaRow{Mengen}{x\dsep A\dsep B}
\DeltaRow{Funktionen}{F}
}
%%end novalidate

\FormulaThmDelta{F\colon A\to B\dsep x\in A\vdash (x,F(x))\in F}{\DeltaRow{Mengen}{x\dsep A\dsep B}
\DeltaRow{Funktionen}{F}
}
\begin{tabproof}
    \proofstep{1}{F\colon A\to B}{\rA}
    \proofstep{2}{x\in A}{\rA}
    \proofstep{1,2}{(x,F(x))\in F}{\FormulaRefAuto{F\colon A\to B\dsep x\in A\vdash (F(x) \coloneqq \iota y\,\bigl((x,y)\in F\bigr)}{1,2}}
\end{tabproof}

\subsection{Eigenschaften des Funktionswertes}

\FormulaThmDelta[Ersetzung im Funktionsargument]{%
  x = y \dsep a = F(y) \vdash a = F(x)
}{
  \DeltaRow{Mengen}{A \dsep B \dsep x \dsep y \dsep a}
  \DeltaRow{Funktionensymbol}{F}
}
\begin{tabproof}
  \proofstep{1}{x = y}{\rA}
  \proofstep{2}{a = F(y)}{\rA}
  \proofstep{1}{y = x}{\FormulaRefAuto{a = b \vdash b = a}{1}}
  \proofstep{1,2}{a = F(x)}{\rIE{3,2}}
\end{tabproof}



\FormulaThmDelta{F\colon A\to B\dsep x\in A\dsep y\in A\dsep x=y\vdash F(x)=F(y)}
{
\DeltaRow{Mengen}{x\dsep y\dsep A\dsep B}
\DeltaRow{Funktionen}{F}
}
\begin{tabproof}
    \proofstep{1}{F\colon A\to B}{\rA}
    \proofstep{2}{x\in A}{\rA}
    \proofstep{3}{y\in A}{\rA}
    \proofstep{4}{x=y}{\rA}
    \proofstep{1,2}{(x,F(x))\in F}{\FormulaRefAuto{F\colon A\to B\dsep x\in A\vdash (x,F(x))\in F}{1,2}}
    \proofstep{1,3}{(y,F(y))\in F}{\FormulaRefAuto{F\colon A\to B\dsep x\in A\vdash (x,F(x))\in F}{1,3}}
    \proofstep{1,2,4}{(y,F(x))\in F}{\rIE{4,5}}
    \proofstep{1,2,3,4}{F(x)=F(y)}{\FormulaRefAuto{(x,y)\in F\dsep (x,z)\in F \vdash y=z}{7,6}}
\end{tabproof}

\FormulaThmDelta{%
  F\colon A\to B\dsep x_1\in A \dsep x_2\in A \dsep
  F(x_1)\neq F(x_2) \vdash x_1\neq x_2%
}{
  \DeltaRow{Mengen}{A\dsep B\dsep x_1\dsep x_2}
  \DeltaRow{Funktionen}{F}
}
\begin{tabproof}
  \proofstep{1}{F\colon A\to B}{\rA}
  \proofstep{2}{x_1\in A}{\rA}
  \proofstep{3}{x_2\in A}{\rA}
  \proofstep{4}{F(x_1)\neq F(x_2)}{\rA}

  \proofstep{5}{x_1=x_2}{\rA}
  \proofstep{1,2,3,5}{F(x_1)=F(x_2)}{%
    \FormulaRefAuto{F\colon A\to B\dsep x\in A\dsep y\in A\dsep x=y\vdash F(x)=F(y)}{1,2,3,5}}
  \proofstep{1,2,3,4,5}{\bot}{\rBI{4,6}}
  \proofstep{1,2,3,4}{x_1\neq x_2}{\rCI{5,7}}
\end{tabproof}

\FormulaThmDelta{F\colon A\to B\dsep x\in A\dsep F(x)=y\vdash (x,y)\in F}
{
\DeltaRow{Mengen}{x\dsep y\dsep A\dsep B}
\DeltaRow{Funktionen}{F}
}
\begin{tabproof}
  \proofstep{1}{F\colon A\to B}{\rA}
  \proofstep{2}{x\in A}{\rA}
  \proofstep{3}{F(x)=y}{\rA}
  \proofstep{1,2}{(x,F(x))\in F}{\FormulaRefAuto{F\colon A\to B\dsep x\in A\vdash (x,F(x))\in F}{1,2}}
  \proofstep{1,2,3}{(x,y)\in F}{\rIE{3,4}}
\end{tabproof}

\FormulaThmDelta{F\colon A\to B\dsep x\in A\dsep y=F(x)\vdash (x,y)\in F}
{
\DeltaRow{Mengen}{x\dsep y\dsep A\dsep B}
\DeltaRow{Funktionen}{F}
}
\begin{tabproof}
  \proofstep{1}{F\colon A\to B}{\rA}
  \proofstep{2}{x\in A}{\rA}
  \proofstep{3}{y=F(x)}{\rA}
  \proofstep{3}{F(x)=y}{\FormulaRefAuto{a=b\vdash b=a}{3}}
  \proofstep{1,2,3}{(x,y)\in F}{\FormulaRefAuto{F\colon A\to B\dsep x\in A\dsep F(x)=y\vdash (x,y)\in F}{1,2,4}}
\end{tabproof}

\FormulaThmDelta{F\colon A\to B\dsep x\in A\vdash F(x)\in B}{
\DeltaRow{Mengen}{A\dsep B}
}
\begin{tabproof}
  \proofstep{1}{F\colon A\to B}{\rA}
  \proofstep{2}{x\in A}{\rA}
  \proofstep{1}{F\subseteq A\times B}{\FormulaRefAuto{F \subseteq A \times B}{1}}
  \proofstep{1,2}{(x,F(x))\in F}{\FormulaRefAuto{F\colon A\to B\dsep x\in A\vdash (F(x) \coloneqq \iota y\,\bigl((x,y)\in F\bigr)}{1,2}}
  \proofstep{1,2}{(x,F(x))\in A\times B}{\FormulaRefAuto{ A\subseteq B,\, x\in A \vdash x\in B }{3,4}}
  \proofstep{1,2}{F(x)\in B}{\FormulaRefAuto{(a,b)\in A\times B\vdash b\in B}{5}}
\end{tabproof}

\FormulaThmDelta{F\colon A\to B\dsep (x,y)\in F\vdash y=F(x)}{
\DeltaRow{Mengen}{x\dsep y\dsep A\dsep B}
\DeltaRow{Funktionen}{F}
}
\begin{tabproof}
  \proofstep{1}{F\colon A\to B}{\rA}
  \proofstep{2}{(x,y)\in F}{\rA}
  \proofstep{1,2}{x\in A}{\FormulaRefAuto{F\colon A\to B\dsep (x,y)\in F\vdash x\in A}{1,2}}
  \proofstep{1,2}{(x,F(x))\in F}{\FormulaRefAuto{F\colon A\to B\dsep x\in A\vdash (x,F(x))\in F}{1,3}}
  \proofstep{1,2}{y=F(x)}{\FormulaRefAuto{(x,y)\in F\dsep (x,z)\in F \vdash y=z}{2,4}}
\end{tabproof}

\FormulaThmDelta{F\colon A\to B\dsep (x,y)\in F\vdash F(x)=y}{%
  \DeltaRow{Mengen}{x\dsep y\dsep A\dsep B}
  \DeltaRow{Funktionen}{F}
}
\begin{tabproof}
  \proofstep{1}{F\colon A\to B}{\rA}
  \proofstep{2}{(x,y)\in F}{\rA}
  \proofstep{1,2}{y=F(x)}{%
    \FormulaRefAuto{F\colon A\to B\dsep (x,y)\in F\vdash y=F(x)}{1,2}}
  \proofstep{1,2}{F(x)=y}{%
    \FormulaRefAuto{a=b\vdash b=a}{3}}
\end{tabproof}

\FormulaThmDelta{%
  F\colon A\to B\dsep G\colon A\to B\dsep
  F(x)=G(x)\dsep (x,y)\in F\vdash (x,y)\in G%
}{
\DeltaRow{Mengen}{x\dsep A\dsep B}
\DeltaRow{Funktionen}{F\dsep G}
}
\begin{tabproof}
  \proofstep{1}{F\colon A\to B}{\rA}
  \proofstep{2}{G\colon A\to B}{\rA}
  \proofstep{3}{F(x)=G(x)}{\rA}
  \proofstep{4}{(x,y)\in F}{\rA}
  \proofstep{1,4}{y=F(x)}{%
    \FormulaRefAuto{F\colon A\to B\dsep (x,y)\in F\vdash y=F(x)}{1,4}}
  \proofstep{1,3,4}{y=G(x)}{\FormulaRefAuto{a=b, b=c \vdash a=c}{5,3}}
  \proofstep{1,4}{x\in A}{%
    \FormulaRefAuto{F\colon A\to B\dsep (x,y)\in F\vdash x\in A}{1,4}}
  \proofstep{1,2,3,4}{(x,y)\in G}{%
    \FormulaRefAuto{F\colon A\to B\dsep x\in A\dsep y=F(x)\vdash (x,y)\in F}{2,7,6}}
\end{tabproof}

\FormulaThmDelta{%
  \{x\in A\mid F(x)\in \{a\}\}=\{x\in A\mid F(x)=a\}%
}{
  \DeltaRow{Mengen}{A\dsep a\dsep x}
  \DeltaRow{Funktionensymbol}{F}
}
\begin{tabproofsplit}
  \proofpart{Teilmengenrichtung \(\{x\in A\mid F(x)\in \{a\}\}\subseteq \{x\in A\mid F(x)=a\}\)}
  \proofstep{1}{x\in \{x\in A\mid F(x)\in \{a\}\}}{\rA}
  \proofstep{1}{x\in A\land F(x)\in \{a\}}{%
    \FormulaRefAuto{x \in \{u \in A \mid P(u)\} \eqvdash x \in A \land P(x)}{1}}
  \proofstep{1}{x\in A}{\rAEa{2}}
  \proofstep{1}{F(x)\in \{a\}}{\rAEb{2}}
  \proofstep{1}{F(x)=a}{\FormulaRefAuto{x \in \{a\} \eqvdash x = a}{4}}
  \proofstep{1}{x\in \{x\in A\mid F(x)=a\}}{%
    \FormulaRefAuto{x \in A\dsep P(x)\vdash x \in \{x \in A \mid P(x)\}}{3,5}}
  \proofstep{}{\{x\in A\mid F(x)\in \{a\}\}\subseteq \{x\in A\mid F(x)=a\}}{%
    \FormulaRefAuto{A \subseteq B \coloneq \forall x\,(x\in A \rightarrow x\in B)}{\rUI{\rRI{1,6}}}}
  \closeproofpart

  \proofpart{Teilmengenrichtung \(\{x\in A\mid F(x)=a\}\subseteq \{x\in A\mid F(x)\in \{a\}\}\)}
  \setcounter{proofstepnr}{7}%
  \proofstep{8}{x\in \{x\in A\mid F(x)=a\}}{\rA}
  \proofstep{8}{x\in A\land F(x)=a}{%
    \FormulaRefAuto{x \in \{u \in A \mid P(u)\} \eqvdash x \in A \land P(x)}{8}}
  \proofstep{8}{x\in A}{\rAEa{9}}
  \proofstep{8}{F(x)=a}{\rAEb{9}}
  \proofstep{8}{F(x)\in \{a\}}{\FormulaRefAuto{x \in \{a\} \eqvdash x = a}{11}}
  \proofstep{8}{x\in \{x\in A\mid F(x)\in \{a\}\}}{%
    \FormulaRefAuto{x \in A\dsep P(x)\vdash x \in \{x \in A \mid P(x)\}}{10,12}}
  \proofstep{}{\{x\in A\mid F(x)=a\}\subseteq \{x\in A\mid F(x)\in \{a\}\}}{%
    \FormulaRefAuto{A \subseteq B \coloneq \forall x\,(x\in A \rightarrow x\in B)}{\rUI{\rRI{8,13}}}}

  \proofstep{}{\{x\in A\mid F(x)\in \{a\}\}=\{x\in A\mid F(x)=a\}}{%
    \FormulaRefAuto{A \subseteq B, B \subseteq A \vdash A = B}{7,14}}
  \closeproofpart
\end{tabproofsplit}


\subsection{Zur Gleichheit von Funktionen}

\FormulaThmDeltaR{%
  \begin{aligned}[t]
    &F\colon A\to B\dsep G\colon A\to B\dsep
      \forall x\in A\,(F(x)=G(x))\vdash{}\\[-2pt]
    &\forall (x,y)\,((x,y)\in F\rightarrow (x,y)\in G)
  \end{aligned}%
}{%
  F\colon A\to B\dsep G\colon A\to B\dsep
  \forall x\in A\,(F(x)=G(x))
  \vdash \forall (x,y)\,((x,y)\in F\rightarrow (x,y)\in G)%
}{
\DeltaRow{Mengen}{A\dsep B\dsep x\dsep y}
\DeltaRow{Funktionen}{F\dsep G}
}
\begin{tabproof}
  \proofstep{1}{F\colon A\to B}{\rA}
  \proofstep{2}{G\colon A\to B}{\rA}
  \proofstep{3}{\forall x\in A\,(F(x)=G(x))}{\rA}
  \proofstep{4}{(x,y)\in F}{\rA}
  \proofstep{1,4}{x\in A}{%
    \FormulaRefAuto{F\colon A\to B\dsep (x,y)\in F\vdash x\in A}{1,4}}
  \proofstep{1,3,4}{F(x)=G(x)}{\rRE{\rUE{3},5}}
  \proofstep{1,2,3,4}{(x,y)\in G}{%
    \FormulaRefAuto{%
      F\colon A\to B\dsep G\colon A\to B\dsep
      F(x)=G(x)\dsep (x,y)\in F\vdash (x,y)\in G%
    }{1,2,6,4}}
  \proofstep{1,2,3}{(x,y)\in F\rightarrow (x,y)\in G}{\rRI{4,7}}
  \proofstep{1,2,3}{\forall (x,y)\,((x,y)\in F\rightarrow (x,y)\in G)}{\rUI{8}}
\end{tabproof}



\FormulaThmDeltaR{%
  \begin{aligned}[t]
    &F\colon A\to B\dsep G\colon A\to B\dsep
      \forall x\in A\,(F(x)=G(x))\vdash{}\\[-2pt]
    &\forall (x,y)\,((x,y)\in G\rightarrow (x,y)\in F)
  \end{aligned}%
}{%
  F\colon A\to B\dsep G\colon A\to B\dsep
  \forall x\in A\,(F(x)=G(x))
  \vdash \forall (x,y)\,((x,y)\in G\rightarrow (x,y)\in F)%
}{
\DeltaRow{Mengen}{A\dsep B\dsep x\dsep y}
\DeltaRow{Funktionen}{F\dsep G}
}
\begin{tabproof}
  \proofstep{1}{F\colon A\to B}{\rA}
  \proofstep{2}{G\colon A\to B}{\rA}
  \proofstep{3}{\forall x\in A\,(F(x)=G(x))}{\rA}
  \proofstep{4}{x\in A}{\rA}
  \proofstep{3,4}{F(x)=G(x)}{\rRE{\rUE{3},4}}
  \proofstep{3,4}{G(x)=F(x)}{\FormulaRefAuto{a=b\vdash b=a}{5}}
  \proofstep{3}{\forall x\in A\,(G(x)=F(x))}{\rUI{\rRI{4,6}}}
  \proofstep{1,2,3}{\forall (x,y)\,((x,y)\in G\rightarrow (x,y)\in F)}{%
    \FormulaRefAuto{%
      F\colon A\to B\dsep G\colon A\to B\dsep
      \forall x\in A\,(F(x)=G(x))
      \vdash \forall (x,y)\,((x,y)\in F\rightarrow (x,y)\in G)%
    }{2,1,7}}
\end{tabproof}

\FormulaThmDeltaR{%
  \begin{aligned}[t]
    &F\colon A\to B\dsep G\colon A\to B\vdash{}\\[-2pt]
    &\Bigl(
      \forall x\in A\,(F(x)=G(x))\leftrightarrow{}\\[-2pt]
    &\qquad
      \forall (x,y)\,((x,y)\in F\leftrightarrow (x,y)\in G)
    \Bigr)
  \end{aligned}%
}{%
  F\colon A\to B\dsep G\colon A\to B\vdash
  \Bigl(
    \forall x\in A\,(F(x)=G(x))
    \leftrightarrow
    \forall (x,y)\,((x,y)\in F\leftrightarrow (x,y)\in G)
  \Bigr)%
}{
\DeltaRow{Mengen}{A\dsep B\dsep x\dsep y}
\DeltaRow{Funktionen}{F\dsep G}
}
\begin{tabproofsplit}
\proofpart{\(\vdash\)}
  \proofstep{1}{F\colon A\to B}{\rA}
  \proofstep{2}{G\colon A\to B}{\rA}
  \proofstep{3}{\forall x\in A\,(F(x)=G(x))}{\rA}
  \proofstep{1,2,3}{\forall (x,y)\,((x,y)\in F\rightarrow (x,y)\in G)}{%
    \FormulaRefAuto{%
      F\colon A\to B\dsep G\colon A\to B\dsep
      \forall x\in A\,(F(x)=G(x))
      \vdash \forall (x,y)\,((x,y)\in F\rightarrow (x,y)\in G)%
    }{1,2,3}}
  \proofstep{1,2,3}{\forall (x,y)\,((x,y)\in G\rightarrow (x,y)\in F)}{%
    \FormulaRefAuto{%
      F\colon A\to B\dsep G\colon A\to B\dsep
      \forall x\in A\,(F(x)=G(x))
      \vdash \forall (x,y)\,((x,y)\in G\rightarrow (x,y)\in F)%
    }{1,2,3}}
  \proofstep{1,2,3}{\forall (x,y)\,((x,y)\in F\leftrightarrow (x,y)\in G)}{%
    \FormulaRefAuto{%
      \forall x\,(P(x)\leftrightarrow Q(x))
      \eqvdash
      \forall x\,(P(x)\rightarrow Q(x))\land
      \forall x\,(Q(x)\rightarrow P(x))%
    }{4,5}}
\closeproofpart
\proofpart{\(\dashv\)}
  \proofstep{1}{F\colon A\to B}{\rA}
  \proofstep{2}{G\colon A\to B}{\rA}
  \proofstep{3}{\forall (x,y)\,((x,y)\in F\leftrightarrow (x,y)\in G)}{\rA}
  \proofstep{4}{x\in A}{\rA}
  \proofstep{1,4}{(x,F(x))\in F}{%
    \FormulaRefAuto{F\colon A\to B\dsep x\in A\vdash (x,F(x))\in F}{1,4}}
  \proofstep{3}{(x,F(x))\in F\leftrightarrow (x,F(x))\in G}{\rUI{3}}
  \proofstep{1,3,4}{(x,F(x))\in G}{%
    \FormulaRefAuto{P\leftrightarrow Q\dsep P\vdash Q}{6,5}}
  \proofstep{2,4}{(x,G(x))\in G}{%
    \FormulaRefAuto{F\colon A\to B\dsep x\in A\vdash (x,F(x))\in F}{2,4}}
  \proofstep{1,2,3,4}{F(x)=G(x)}{%
    \FormulaRefAuto{(x,y)\in F\dsep (x,z)\in F\vdash y=z}{7,8}}
  \proofstep{1,2,3}{\forall x\in A\,(F(x)=G(x))}{\rUI{\rRI{4,9}}}
\closeproofpart
\end{tabproofsplit}

\FormulaThmDeltaK[Funktionsextensionalität]{%
  F\colon A\to B\dsep G\colon A\to B\dsep
  \forall x\in A\,(F(x)=G(x))\vdash F=G%
}{FunctionExtensionality}{
\DeltaRow{Mengen}{A\dsep B\dsep x}
\DeltaRow{Funktionen}{F\dsep G}
}
\begin{tabproof}
  \proofstep{1}{F\colon A\to B}{\rA}
  \proofstep{2}{G\colon A\to B}{\rA}
  \proofstep{3}{\forall x\in A\,(F(x)=G(x))}{\rA}
  \proofstep{1,2}{%
    \forall x\in A\,(F(x)=G(x))\leftrightarrow
    \forall (x,y)\,((x,y)\in F\leftrightarrow (x,y)\in G)%
  }{%
    \FormulaRefAuto{%
      F\colon A\to B\dsep G\colon A\to B\vdash
      \Bigl(
        \forall x\in A\,(F(x)=G(x))
        \leftrightarrow
        \forall (x,y)\,((x,y)\in F\leftrightarrow (x,y)\in G)
      \Bigr)%
    }{1,2}}
  \proofstep{1,2,3}{\forall (x,y)\,((x,y)\in F\leftrightarrow (x,y)\in G)}{%
    \FormulaRefAuto{P\leftrightarrow Q\dsep P\vdash Q}{4,3}}
  \proofstep{1,2,3}{F=G}{%
    \FormulaRefAuto{\forall x\,(x\in A\leftrightarrow x\in B)\eqvdash A=B}{5}}
\end{tabproof}

\FormulaThmDeltaK[Auswertung gleicher Funktionen]{%
  F\colon A\to B\dsep G\colon A\to B\dsep x\in A\dsep F=G
  \vdash F(x)=G(x)%
}{FunctionEqualityEvaluation}{
\DeltaRow{Mengen}{A\dsep B\dsep x}
\DeltaRow{Funktionen}{F\dsep G}
}
\begin{tabproof}
    \proofstep{1}{F\colon A\to B}{\rA}
    \proofstep{2}{G\colon A\to B}{\rA}
    \proofstep{3}{x\in A}{\rA}
    \proofstep{4}{F=G}{\rA}
    \proofstep{}{F(x)=F(x)}{\rII}
    \proofstep{1,2,3,4}{F(x)=G(x)}{\rIE{4,5}}
\end{tabproof}

\section{Die Bildmenge}

\subsection{Bildmengen von Teilmengen}

\FormulaThmDeltaK[Eindeutige Existenz der Bildmenge einer Teilmenge]{%
  C\subseteq A\dsep F\colon A\to B
  \vdash
  \exists!D\,\forall y\,
  \Bigl(
    y\in D\leftrightarrow
    \exists x\in C\,\bigl(y=F(x)\bigr)
  \Bigr)%
}{SubsetFunctionImageSetUniqueExists}{%
  \DeltaRow{Mengen}{A\dsep B\dsep C\dsep D\dsep u\dsep x\dsep y}
  \DeltaRow{Funktionen}{F}
}
\begin{tabproof}
  \proofstep{1}{C\subseteq A}{\rA}
  \proofstep{2}{F\colon A\to B}{\rA}
  \proofstep{3}{\exists x\in C\,\bigl(y=F(x)\bigr)}{\rA}
  \proofstep{4}{u\in C\land y=F(u)}{\rA}
  \proofstep{4}{u\in C}{\rAEa{4}}
  \proofstep{4}{y=F(u)}{\rAEb{4}}
  \proofstep{1,4}{u\in A}{%
    \FormulaRefAuto{A\subseteq B,\,x\in A\vdash x\in B}{1,5}}
  \proofstep{1,2,4}{F(u)\in B}{%
    \FormulaRefAuto{F\colon A\to B\dsep x\in A\vdash F(x)\in B}{2,7}}
  \proofstep{1,2,4}{y\in B}{\rIE{6,8}}
  \proofstep{1,2,3}{y\in B}{\rEE{3,4,9}}
  \proofstep{1,2}{%
    \forall y\,
    \Bigl(
      \exists x\in C\,\bigl(y=F(x)\bigr)
      \rightarrow y\in B
    \Bigr)
  }{\rUI{\rRI{3,10}}}
  \proofstep{1,2}{%
    \exists!D\,\forall y\,
    \Bigl(
      y\in D\leftrightarrow
      \exists x\in C\,\bigl(y=F(x)\bigr)
    \Bigr)
  }{%
    \FormulaRefAuto{%
      \forall x(P(x)\rightarrow x\in A)
      \vdash
      \exists! B(\forall x(x\in B\leftrightarrow P(x)))
    }{11}}
\end{tabproof}

%%begin novalidate
\FormulaDefDelta[Bildmenge einer Teilmenge]
{%
  C\subseteq A\dsep F\colon A\to B\vdash
  F[C]\coloneqq \iota D\Bigl(\forall y\,\bigl(y\in D \leftrightarrow \exists x\in C\,y=F(x)\bigr)\Bigr)
}
{
  \DeltaRow{Mengen}{A\dsep B\dsep C\dsep x}
  \DeltaRow{Funktionen}{F}
}
%%end novalidate

\FormulaThmDelta[Charakterisierung der Bildmenge einer Teilmenge]{%
  C\subseteq A\dsep F\colon A\to B\vdash y\in F[C]\leftrightarrow \exists x\in C\,y=F(x)
}
{
  \DeltaRow{Mengen}{A\dsep B\dsep C\dsep x\dsep y}
  \DeltaRow{Funktionen}{F}
}
\begin{tabproof}
  \proofstep{1}{C\subseteq A}{\rA}
  \proofstep{2}{F\colon A\to B}{\rA}
  \proofstep{1,2}{\forall y\,\bigl(y\in F[C]\leftrightarrow \exists x\in C\,y=F(x)\bigr)}{%
    \FormulaRefAuto{C\subseteq A\dsep F\colon A\to B\vdash F[C]\coloneqq \iota D\Bigl(\forall y\,\bigl(y\in D \leftrightarrow \exists x\in C\,y=F(x)\bigr)\Bigr)}{1,2}}
  \proofstep{1,2}{y\in F[C]\leftrightarrow \exists x\in C\,y=F(x)}{\rUE{3}}
\end{tabproof}

\FormulaThmDelta[Aus einem Bildelement einer Teilmenge folgt ein Urbild in der Teilmenge]{%
  C\subseteq A\dsep F\colon A\to B\dsep y\in F[C]\vdash \exists x\in C\,y=F(x)
}
{
  \DeltaRow{Mengen}{A\dsep B\dsep C\dsep x\dsep y}
  \DeltaRow{Funktionen}{F}
}
\begin{tabproof}
  \proofstep{1}{C\subseteq A}{\rA}
  \proofstep{2}{F\colon A\to B}{\rA}
  \proofstep{3}{y\in F[C]}{\rA}
  \proofstep{1,2}{y\in F[C]\leftrightarrow \exists x\in C\,y=F(x)}{%
    \FormulaRefAuto{C\subseteq A\dsep F\colon A\to B\vdash y\in F[C]\leftrightarrow \exists x\in C\,y=F(x)}{1,2}}
  \proofstep{1,2,3}{\exists x\in C\,y=F(x)}{%
    \FormulaRefAuto{P\leftrightarrow Q, P \vdash Q}{4,3}}
\end{tabproof}

\FormulaThmDelta[Ein Urbild in einer Teilmenge liefert ein Bildelement der Teilmenge]{%
  C\subseteq A\dsep F\colon A\to B\dsep \exists x\in C\,y=F(x)\vdash y\in F[C]
}
{
  \DeltaRow{Mengen}{A\dsep B\dsep C\dsep x\dsep y}
  \DeltaRow{Funktionen}{F}
}
\begin{tabproof}
  \proofstep{1}{C\subseteq A}{\rA}
  \proofstep{2}{F\colon A\to B}{\rA}
  \proofstep{3}{\exists x\in C\,y=F(x)}{\rA}
  \proofstep{1,2}{y\in F[C]\leftrightarrow \exists x\in C\,y=F(x)}{%
    \FormulaRefAuto{C\subseteq A\dsep F\colon A\to B\vdash y\in F[C]\leftrightarrow \exists x\in C\,y=F(x)}{1,2}}
  \proofstep{1,2,3}{y\in F[C]}{%
    \FormulaRefAuto{P\leftrightarrow Q, Q \vdash P}{4,3}}
\end{tabproof}

\FormulaThmDelta[Elemente einer Teilmenge liegen im Bild der Teilmenge]{%
  C\subseteq A\dsep F\colon A\to B\dsep x\in C\vdash F(x)\in F[C]
}
{
  \DeltaRow{Mengen}{A\dsep B\dsep C\dsep x\dsep z}
  \DeltaRow{Funktionen}{F}
}
\begin{tabproof}
  \proofstep{1}{C\subseteq A}{\rA}
  \proofstep{2}{F\colon A\to B}{\rA}
  \proofstep{3}{x\in C}{\rA}
  \proofstep{}{F(x)=F(x)}{\rII}
  \proofstep{3}{\exists z\in C\,F(x)=F(z)}{\rEI{\rAI{3,4}}}
  \proofstep{1,2,3}{F(x)\in F[C]}{%
    \FormulaRefAuto{C\subseteq A\dsep F\colon A\to B\dsep \exists x\in C\,y=F(x)\vdash y\in F[C]}{1,2,5}}
\end{tabproof}

\paragraph{Spezialfall \(C=A\)}

\FormulaThmDelta{F\colon A\to B\vdash y\in F[A]\leftrightarrow \exists x\in A\,y=F(x)}
{
  \DeltaRow{Mengen}{A\dsep B\dsep x\dsep y}
  \DeltaRow{Funktionen}{F}
}
\begin{tabproof}
  \proofstep{}{A\subseteq A}{\FormulaRefAuto{A\subseteq A}}
  \proofstep{2}{F\colon A\to B}{\rA}
  \proofstep{1,2}{y\in F[A]\leftrightarrow \exists x\in A\,y=F(x)}{%
    \FormulaRefAuto{C\subseteq A\dsep F\colon A\to B\vdash y\in F[C]\leftrightarrow \exists x\in C\,y=F(x)}{1,2}}
\end{tabproof}

\FormulaThmDelta[Aus einem Bildelement folgt ein Urbild]{%
  F\colon A\to B\dsep y\in F[A]\vdash \exists x\in A\,y=F(x)
}{
  \DeltaRow{Mengen}{A\dsep B\dsep x\dsep y}
  \DeltaRow{Funktionen}{F}
}
\begin{tabproof}
  \proofstep{}{A\subseteq A}{\FormulaRefAuto{A\subseteq A}}
  \proofstep{2}{F\colon A\to B}{\rA}
  \proofstep{3}{y\in F[A]}{\rA}
  \proofstep{1,2,3}{\exists x\in A\,y=F(x)}{%
    \FormulaRefAuto{C\subseteq A\dsep F\colon A\to B\dsep y\in F[C]\vdash \exists x\in C\,y=F(x)}{1,2,3}}
\end{tabproof}

\FormulaThmDelta[Ein Urbild liefert ein Bildelement]{%
  F\colon A\to B\dsep \exists x\in A\,y=F(x)\vdash y\in F[A]
}{
  \DeltaRow{Mengen}{A\dsep B\dsep x\dsep y}
  \DeltaRow{Funktionen}{F}
}
\begin{tabproof}
  \proofstep{}{A\subseteq A}{\FormulaRefAuto{A\subseteq A}}
  \proofstep{2}{F\colon A\to B}{\rA}
  \proofstep{3}{\exists x\in A\,y=F(x)}{\rA}
  \proofstep{1,2,3}{y\in F[A]}{%
    \FormulaRefAuto{C\subseteq A\dsep F\colon A\to B\dsep \exists x\in C\,y=F(x)\vdash y\in F[C]}{1,2,3}}
\end{tabproof}

\subsection{Rechenregeln für Bildmengen}

\FormulaThmDelta[Bild einer Einermenge]{%
  x\in A \dsep F\colon A\to B \vdash F[\{x\}]=\{F(x)\}
}{%
  \DeltaRow{Mengen}{A\dsep B\dsep x\dsep y\dsep z}
  \DeltaRow{Funktionen}{F}
}
\begin{tabproofsplitwide}
  \proofpartwideind[Teilmengenrichtung \(F[\{x\}]\subseteq \{F(x)\}\)]{%
    x\in A \dsep F\colon A\to B \vdash F[\{x\}]\subseteq \{F(x)\}
  }[%
    x\in A \dsep F\colon A\to B \vdash F[\{x\}]\subseteq \{F(x)\}
  ]
    \proofstepwidestar[1]{x\in A}{\rA}
    \proofstepwidestar[2]{F\colon A\to B}{\rA}
    \proofstepwidestar[3]{y\in F[\{x\}]}{\rA}
    \proofstepwidestar[1]{\{x\}\subseteq A}{\FormulaRefAuto{a\in A \vdash \{a\}\subseteq A}{1}}
    \proofstepwidestar[1,2,3]{\exists z\in \{x\}\,y=F(z)}{%
      \FormulaRefAuto{C\subseteq A\dsep F\colon A\to B\dsep y\in F[C]\vdash \exists x\in C\,y=F(x)}{4,2,3}}

    \proofstepwidestar[6]{z\in \{x\}\land y=F(z)}{\rA}
    \proofstepwidestar[6]{z\in \{x\}}{\rAEa{6}}
    \proofstepwidestar[6]{y=F(z)}{\rAEb{6}}
    \proofstepwidestar[6]{z=x}{\FormulaRefAuto{x\in \{a\}\eqvdash x=a}{7}}
    \proofstepwidestar[6]{x=z}{\FormulaRefAuto{a=b\vdash b=a}{9}}
    \proofstepwidestar[6]{y=F(x)}{%
      \FormulaRefAuto{x=y \dsep a=F(y) \vdash a=F(x)}{10,8}}
    \proofstepwidestar[6]{y\in \{F(x)\}}{\FormulaRefAuto{x\in \{a\}\eqvdash x=a}{11}}

    \proofstepwidestar[1,2,3]{y\in \{F(x)\}}{\rEE{5,6,12}}
    \proofstepwidestar[1,2]{\forall y\,(y\in F[\{x\}]\rightarrow y\in \{F(x)\})}{\rUI{\rRI{3,13}}}
    \proofstepwidestar[1,2]{F[\{x\}]\subseteq \{F(x)\}}{%
      \FormulaRefAuto{A\subseteq B \coloneq \forall x\,(x\in A \rightarrow x\in B)}{14}}
  \closeproofpartwideind

  \proofpartwideind[Teilmengenrichtung \(\{F(x)\}\subseteq F[\{x\}]\)]{%
    x\in A \dsep F\colon A\to B \vdash \{F(x)\}\subseteq F[\{x\}]
  }[%
    x\in A \dsep F\colon A\to B \vdash \{F(x)\}\subseteq F[\{x\}]
  ]
    \proofstepwidestar[1]{x\in A}{\rA}
    \proofstepwidestar[2]{F\colon A\to B}{\rA}
    \proofstepwidestar[3]{y\in \{F(x)\}}{\rA}
    \proofstepwidestar[3]{y=F(x)}{\FormulaRefAuto{x\in \{a\}\eqvdash x=a}{3}}
    \proofstepwidestar[]{x=x}{\rII}
    \proofstepwidestar[]{x\in \{x\}}{\FormulaRefAuto{x\in \{a\}\eqvdash x=a}{5}}
    \proofstepwidestar[3]{x\in \{x\}\land y=F(x)}{\rAI{6,4}}
    \proofstepwidestar[3]{\exists z\in \{x\}\,y=F(z)}{\rEI{7}}
    \proofstepwidestar[1]{\{x\}\subseteq A}{\FormulaRefAuto{a\in A \vdash \{a\}\subseteq A}{1}}
    \proofstepwidestar[1,2,3]{y\in F[\{x\}]}{%
      \FormulaRefAuto{C\subseteq A\dsep F\colon A\to B\dsep \exists x\in C\,y=F(x)\vdash y\in F[C]}{9,2,8}}
    \proofstepwidestar[1,2]{\forall y\,(y\in \{F(x)\}\rightarrow y\in F[\{x\}])}{\rUI{\rRI{3,10}}}
    \proofstepwidestar[1,2]{\{F(x)\}\subseteq F[\{x\}]}{%
      \FormulaRefAuto{A\subseteq B \coloneq \forall x\,(x\in A \rightarrow x\in B)}{11}}
  \closeproofpartwideind

  \proofpartwide{Zusammenfassung}
    \proofstepwidestar[1]{x\in A}{\rA}
    \proofstepwidestar[2]{F\colon A\to B}{\rA}
    \proofstepwidestar[1,2]{F[\{x\}]\subseteq \{F(x)\}}{%
      \FormulaRefAuto{x\in A \dsep F\colon A\to B \vdash F[\{x\}]\subseteq \{F(x)\}}{1,2}}
    \proofstepwidestar[1,2]{\{F(x)\}\subseteq F[\{x\}]}{%
      \FormulaRefAuto{x\in A \dsep F\colon A\to B \vdash \{F(x)\}\subseteq F[\{x\}]}{1,2}}
    \proofstepwidestar[1,2]{F[\{x\}]=\{F(x)\}}{%
      \FormulaRefAuto{A\subseteq B\dsep B\subseteq A \vdash A=B}{3,4}}
  \closeproofpartwide
\end{tabproofsplitwide}

\FormulaThmDelta[Monotonie der Bildmenge]{%
  F\colon M\to N \dsep A\subseteq B \dsep B\subseteq M \vdash F[A]\subseteq F[B]
}{%
  \DeltaRow{Mengen}{A\dsep B\dsep M\dsep N\dsep x\dsep y}
  \DeltaRow{Funktionen}{F}
}
\begin{tabproof}
  \proofstep{1}{F\colon M\to N}{\rA}
  \proofstep{2}{A\subseteq B}{\rA}
  \proofstep{3}{B\subseteq M}{\rA}
  \proofstep{4}{y\in F[A]}{\rA}
  \proofstep{2,3}{A\subseteq M}{%
    \FormulaRefAuto{A\subseteq B,\, B\subseteq C \vdash A\subseteq C}{2,3}}
  \proofstep{1,2,3,4}{\exists x\in A\,y=F(x)}{%
    \FormulaRefAuto{C\subseteq A\dsep F\colon A\to B\dsep y\in F[C]\vdash \exists x\in C\,y=F(x)}{5,1,4}}
  \proofstep{7}{x\in A\land y=F(x)}{\rA}
  \proofstep{7}{x\in A}{\rAEa{7}}
  \proofstep{7}{y=F(x)}{\rAEb{7}}
  \proofstep{2,8}{x\in B}{\FormulaRefAuto{A\subseteq B\dsep x\in A\vdash x\in B}{2,8}}
  \proofstep{2,7}{x\in B\land y=F(x)}{\rAI{10,9}}
  \proofstep{2,7}{\exists x\in B\,y=F(x)}{\rEI{11}}
  \proofstep{1,2,3,7}{y\in F[B]}{%
    \FormulaRefAuto{C\subseteq A\dsep F\colon A\to B\dsep \exists x\in C\,y=F(x)\vdash y\in F[C]}{3,1,12}}
  \proofstep{1,2,3,4}{y\in F[B]}{\rEE{6,7,13}}
  \proofstep{1,2,3}{\forall y\,(y\in F[A]\rightarrow y\in F[B])}{\rUI{\rRI{4,14}}}
  \proofstep{1,2,3}{F[A]\subseteq F[B]}{%
    \FormulaRefAuto{A\subseteq B \coloneq \forall x\,(x\in A \rightarrow x\in B)}{15}}
\end{tabproof}

\FormulaThmDelta[Bild eines Elements aus einer Vereinigung]{%
  F\colon M\to N \dsep A\subseteq M \dsep B\subseteq M \dsep x\in A\cup B \vdash F(x)\in F[A]\cup F[B]
}{%
  \DeltaRow{Mengen}{A\dsep B\dsep M\dsep N\dsep x}
  \DeltaRow{Funktionen}{F}
}
\begin{tabproof}
  \proofstep{1}{F\colon M\to N}{\rA}
  \proofstep{2}{A\subseteq M}{\rA}
  \proofstep{3}{B\subseteq M}{\rA}
  \proofstep{4}{x\in A\cup B}{\rA}
  \proofstep{4}{x\in A\lor x\in B}{%
    \FormulaRefAuto{x\in A\cup B\eqvdash x\in A\lor x\in B}{4}}

  \proofstep{6}{x\in A}{\rA}
  \proofstep{1,2,6}{F(x)\in F[A]}{%
    \FormulaRefAuto{C\subseteq A\dsep F\colon A\to B\dsep x\in C\vdash F(x)\in F[C]}{2,1,6}}
  \proofstep{1,2,6}{F(x)\in F[A]\cup F[B]}{%
    \FormulaRefAuto{y\in A\vdash y\in A\cup B}{7}}

  \proofstep{9}{x\in B}{\rA}
  \proofstep{1,3,9}{F(x)\in F[B]}{%
    \FormulaRefAuto{C\subseteq A\dsep F\colon A\to B\dsep x\in C\vdash F(x)\in F[C]}{3,1,9}}
  \proofstep{1,3,9}{F(x)\in F[A]\cup F[B]}{%
    \FormulaRefAuto{y\in B\vdash y\in A\cup B}{10}}

  \proofstep{1,2,3,4}{F(x)\in F[A]\cup F[B]}{\rOE{5,6,8,9,11}}
\end{tabproof}

\FormulaThmDelta[Bild einer Vereinigung]{%
  F\colon M\to N \dsep A\subseteq M \dsep B\subseteq M \vdash F[A\cup B]=F[A]\cup F[B]
}{%
  \DeltaRow{Mengen}{A\dsep B\dsep M\dsep N\dsep x\dsep y}
  \DeltaRow{Funktionen}{F}
}
\begin{tabproofsplitwide}
  \proofpartwideind[Teilmengenrichtung \(F[A\cup B]\subseteq F[A]\cup F[B]\)]{%
    F\colon M\to N \dsep A\subseteq M \dsep B\subseteq M \vdash F[A\cup B]\subseteq F[A]\cup F[B]
  }[%
    F\colon M\to N \dsep A\subseteq M \dsep B\subseteq M \vdash F[A\cup B]\subseteq F[A]\cup F[B]
  ]
    \proofstepwidestar[1]{F\colon M\to N}{\rA}
    \proofstepwidestar[2]{A\subseteq M}{\rA}
    \proofstepwidestar[3]{B\subseteq M}{\rA}
    \proofstepwidestar[4]{y\in F[A\cup B]}{\rA}
    \proofstepwidestar[2,3]{A\cup B\subseteq M}{%
      \FormulaRefAuto{A\subseteq M\dsep B\subseteq M \vdash A\cup B\subseteq M}{2,3}}
    \proofstepwidestar[1,2,3,4]{\exists x\in A\cup B\,y=F(x)}{%
      \FormulaRefAuto{C\subseteq A\dsep F\colon A\to B\dsep y\in F[C]\vdash \exists x\in C\,y=F(x)}{5,1,4}}
    \proofstepwidestar[7]{x\in A\cup B\land y=F(x)}{\rA}
    \proofstepwidestar[7]{x\in A\cup B}{\rAEa{7}}
    \proofstepwidestar[7]{y=F(x)}{\rAEb{7}}
    \proofstepwidestar[1,2,3,8]{F(x)\in F[A]\cup F[B]}{%
      \FormulaRefAuto{F\colon M\to N \dsep A\subseteq M \dsep B\subseteq M \dsep x\in A\cup B \vdash F(x)\in F[A]\cup F[B]}{1,2,3,8}}
    \proofstepwidestar[1,2,3,7]{y\in F[A]\cup F[B]}{\rIE{9,10}}
    \proofstepwidestar[1,2,3,4]{y\in F[A]\cup F[B]}{\rEE{6,7,11}}
    \proofstepwidestar[1,2,3]{\forall y\,\bigl(y\in F[A\cup B]\rightarrow y\in F[A]\cup F[B]\bigr)}{\rUI{\rRI{4,12}}}
    \proofstepwidestar[1,2,3]{F[A\cup B]\subseteq F[A]\cup F[B]}{%
      \FormulaRefAuto{A\subseteq B \coloneq \forall x\,(x\in A \rightarrow x\in B)}{13}}
  \closeproofpartwideind

  \proofpartwideind[Teilmengenrichtung \(F[A]\cup F[B]\subseteq F[A\cup B]\)]{%
    F\colon M\to N \dsep A\subseteq M \dsep B\subseteq M \vdash F[A]\cup F[B]\subseteq F[A\cup B]
  }[%
    F\colon M\to N \dsep A\subseteq M \dsep B\subseteq M \vdash F[A]\cup F[B]\subseteq F[A\cup B]
  ]
    \proofstepwidestar[1]{F\colon M\to N}{\rA}
    \proofstepwidestar[2]{A\subseteq M}{\rA}
    \proofstepwidestar[3]{B\subseteq M}{\rA}
    \proofstepwidestar[2]{A\subseteq A\cup B}{\FormulaRefAuto{A\subseteq A\cup B}{2}}
    \proofstepwidestar[3]{B\subseteq A\cup B}{\FormulaRefAuto{B\subseteq A\cup B}{3}}
    \proofstepwidestar[2,3]{A\cup B\subseteq M}{%
      \FormulaRefAuto{A\subseteq M\dsep B\subseteq M \vdash A\cup B\subseteq M}{2,3}}
    \proofstepwidestar[1,2,3]{F[A]\subseteq F[A\cup B]}{%
      \FormulaRefAuto{F\colon M\to N \dsep A\subseteq B \dsep B\subseteq M \vdash F[A]\subseteq F[B]}{1,4,6}}
    \proofstepwidestar[1,2,3]{F[B]\subseteq F[A\cup B]}{%
      \FormulaRefAuto{F\colon M\to N \dsep A\subseteq B \dsep B\subseteq M \vdash F[A]\subseteq F[B]}{1,5,6}}
    \proofstepwidestar[1,2,3]{F[A]\cup F[B]\subseteq F[A\cup B]}{%
      \FormulaRefAuto{X\subseteq Z\dsep Y\subseteq Z \vdash X\cup Y\subseteq Z}{7,8}}
  \closeproofpartwideind

  \proofpartwide{Zusammenfassung}
    \proofstepwidestar[1]{F\colon M\to N}{\rA}
    \proofstepwidestar[2]{A\subseteq M}{\rA}
    \proofstepwidestar[3]{B\subseteq M}{\rA}
    \proofstepwidestar[1,2,3]{F[A\cup B]\subseteq F[A]\cup F[B]}{%
      \FormulaRefAuto{F\colon M\to N \dsep A\subseteq M \dsep B\subseteq M \vdash F[A\cup B]\subseteq F[A]\cup F[B]}{1,2,3}}
    \proofstepwidestar[1,2,3]{F[A]\cup F[B]\subseteq F[A\cup B]}{%
      \FormulaRefAuto{F\colon M\to N \dsep A\subseteq M \dsep B\subseteq M \vdash F[A]\cup F[B]\subseteq F[A\cup B]}{1,2,3}}
    \proofstepwidestar[1,2,3]{F[A\cup B]=F[A]\cup F[B]}{%
      \FormulaRefAuto{A\subseteq B\dsep B\subseteq A \vdash A=B}{4,5}}
  \closeproofpartwide
\end{tabproofsplitwide}

\FormulaThmDelta[Bild eines Durchschnitts liegt im Durchschnitt der Bilder]{%
  F\colon M\to N \dsep A\subseteq M \dsep B\subseteq M \vdash F[A\cap B]\subseteq F[A]\cap F[B]
}{%
  \DeltaRow{Mengen}{A\dsep B\dsep M\dsep N}
  \DeltaRow{Funktionen}{F}
}
\begin{tabproof}
  \proofstep{1}{F\colon M\to N}{\rA}
  \proofstep{2}{A\subseteq M}{\rA}
  \proofstep{3}{B\subseteq M}{\rA}
  \proofstep{}{A\cap B\subseteq A}{\FormulaRefAuto{A\cap B\subseteq A}}
  \proofstep{}{A\cap B\subseteq B}{\FormulaRefAuto{A\cap B\subseteq B}}
  \proofstep{1,2}{F[A\cap B]\subseteq F[A]}{%
    \FormulaRefAuto{F\colon M\to N \dsep A\subseteq B \dsep B\subseteq M \vdash F[A]\subseteq F[B]}{1,4,2}}
  \proofstep{1,3}{F[A\cap B]\subseteq F[B]}{%
    \FormulaRefAuto{F\colon M\to N \dsep A\subseteq B \dsep B\subseteq M \vdash F[A]\subseteq F[B]}{1,5,3}}
  \proofstep{1,2,3}{F[A\cap B]\subseteq F[A]\cap F[B]}{%
    \FormulaRefAuto{A\subseteq B \dsep A\subseteq C \vdash A\subseteq B\cap C}{6,7}}
\end{tabproof}

\FormulaThmDelta[Das Bild der leeren Menge ist leer]{%
  F\colon A\to B \vdash F[\varnothing]=\varnothing
}{%
  \DeltaRow{Mengen}{A\dsep B\dsep y}
  \DeltaRow{Funktionen}{F}
}
\begin{tabproofsplitwide}
  \proofpartwideind[Teilmengenrichtung \(F[\varnothing]\subseteq \varnothing\)]{%
    F\colon A\to B \vdash F[\varnothing]\subseteq \varnothing
  }[%
    F\colon A\to B \vdash F[\varnothing]\subseteq \varnothing
  ]
    \proofstepwidestar[1]{F\colon A\to B}{\rA}
    \proofstepwidestar[2]{y\in F[\varnothing]}{\rA}
    \proofstepwidestar[]{\varnothing\subseteq A}{\FormulaRefAuto{\varnothing\subseteq A}}
    \proofstepwidestar[1,2]{\exists x\in \varnothing\,y=F(x)}{%
      \FormulaRefAuto{C\subseteq A\dsep F\colon A\to B\dsep y\in F[C]\vdash \exists x\in C\,y=F(x)}{3,1,2}}
    \proofstepwidestar[1,2]{y\in \varnothing}{%
      \FormulaRefAuto{\exists x\in \varnothing\,P(x)\vdash Q}{4}}
    \proofstepwidestar[1]{\forall y\,(y\in F[\varnothing]\rightarrow y\in \varnothing)}{%
      \rUI{\rRI{2,5}}}
    \proofstepwidestar[1]{F[\varnothing]\subseteq \varnothing}{%
      \FormulaRefAuto{A\subseteq B \coloneq \forall x\,(x\in A \rightarrow x\in B)}{6}}
  \closeproofpartwideind

  \proofpartwideind[Teilmengenrichtung \(\varnothing\subseteq F[\varnothing]\)]{%
    F\colon A\to B \vdash \varnothing\subseteq F[\varnothing]
  }[%
    F\colon A\to B \vdash \varnothing\subseteq F[\varnothing]
  ]
    \proofstepwidestar[1]{F\colon A\to B}{\rA}
    \proofstepwidestar[1]{\varnothing\subseteq F[\varnothing]}{\FormulaRefAuto{\varnothing\subseteq A}}
  \closeproofpartwideind

  \proofpartwide{Zusammenfassung}
    \proofstepwidestar[1]{F\colon A\to B}{\rA}
    \proofstepwidestar[1]{F[\varnothing]\subseteq \varnothing}{%
      \FormulaRefAuto{F\colon A\to B \vdash F[\varnothing]\subseteq \varnothing}{1}}
    \proofstepwidestar[1]{\varnothing\subseteq F[\varnothing]}{%
      \FormulaRefAuto{F\colon A\to B \vdash \varnothing\subseteq F[\varnothing]}{1}}
    \proofstepwidestar[1]{F[\varnothing]=\varnothing}{%
      \FormulaRefAuto{A\subseteq B\dsep B\subseteq A \vdash A=B}{2,3}}
  \closeproofpartwide
\end{tabproofsplitwide}

\subsection{Bildmengen im Kontext des Wertebereichs}

\FormulaThmDelta{F\colon A\to B\vdash F\subseteq A\times F[A]}
{
  \DeltaRow{Mengen}{A\dsep B\dsep x\dsep y\dsep z}
  \DeltaRow{Funktionen}{F}
}
\begin{tabproof}
  \proofstep{1}{F\colon A\to B}{\rA}
  \proofstep{2}{(x,y)\in F}{\rA}
  \proofstep{1,2}{x\in A}{\FormulaRefAuto{F\colon A\to B\dsep (x,y)\in F\vdash x\in A}{1,2}}
  \proofstep{1,2}{y=F(x)}{\FormulaRefAuto{F\colon A\to B\dsep (x,y)\in F\vdash y=F(x)}{1,2}}
  \proofstep{}{F(x)=F(x)}{\rII}
  \proofstep{3}{\exists z\in A\,F(x)=F(z)}{\rEI{\rAI{3,5}}}
  \proofstep{1,3}{F(x)\in F[A]}{%
    \FormulaRefAuto{F\colon A\to B\dsep \exists x\in A\,y=F(x)\vdash y\in F[A]}{1,6}}
  \proofstep{1,2,3}{y\in F[A]}{\rIE{4,7}}
  \proofstep{1,2,3}{(x,y)\in A\times F[A]}{%
    \FormulaRefAuto{a\in A\dsep b\in B\vdash (a,b)\in A\times B}{3,8}}
  \proofstep{1}{F\subseteq A\times F[A]}{%
    \FormulaRefAuto{A \subseteq B \coloneq \forall x\,(x\in A \rightarrow x\in B)}{\rUI{\rRI{2,9}}}}
\end{tabproof}

\FormulaThmDelta[Das Bild einer Teilmenge liegt im Zielbereich]{%
  C\subseteq A \dsep F\colon A\to B \vdash F[C]\subseteq B
}{%
  \DeltaRow{Mengen}{A\dsep B\dsep C\dsep x\dsep y}
  \DeltaRow{Funktionen}{F}
}
\begin{tabproof}
  \proofstep{1}{C\subseteq A}{\rA}
  \proofstep{2}{F\colon A\to B}{\rA}
  \proofstep{3}{y\in F[C]}{\rA}
  \proofstep{1,2,3}{\exists x\in C\,y=F(x)}{%
    \FormulaRefAuto{C\subseteq A\dsep F\colon A\to B\dsep y\in F[C]\vdash \exists x\in C\,y=F(x)}{1,2,3}}
  \proofstep{5}{x\in C\land y=F(x)}{\rA}
  \proofstep{5}{x\in C}{\rAEa{5}}
  \proofstep{5}{y=F(x)}{\rAEb{5}}
  \proofstep{1,6}{x\in A}{%
    \FormulaRefAuto{A\subseteq B\dsep x\in A\vdash x\in B}{1,6}}
  \proofstep{1,2,5}{F(x)\in B}{%
    \FormulaRefAuto{F\colon A\to B\dsep x\in A\vdash F(x)\in B}{2,8}}
  \proofstep{1,2,5}{y\in B}{\rIE{7,9}}
  \proofstep{1,2,3}{y\in B}{\rEE{4,5,10}}
  \proofstep{1,2}{\forall y\,(y\in F[C]\rightarrow y\in B)}{\rUI{\rRI{3,11}}}
  \proofstep{1,2}{F[C]\subseteq B}{%
    \FormulaRefAuto{A\subseteq B \coloneq \forall x\,(x\in A \rightarrow x\in B)}{12}}
\end{tabproof}

\FormulaThmDeltaK[Bilder nichtleerer Teilmengen sind nichtleer]{%
  f\colon A\to B\dsep X\in\PowNE{A}\vdash f[X]\in\PowNE{B}%
}{NonemptyDirectImage}{%
  \DeltaRow{Mengen}{A\dsep B\dsep X\dsep x}
  \DeltaRow{Funktionen}{f}
}
\begin{tabproof}
  \proofstep{1}{f\colon A\to B}{\rA}
  \proofstep{2}{X\in\PowNE{A}}{\rA}
  \proofstep{2}{X\subseteq A}{%
    \rAEa{\FormulaRefAuto{NonemptyPowersetMembership}[thm]{2}}}
  \proofstep{2}{X\neq\varnothing}{%
    \rAEb{\FormulaRefAuto{NonemptyPowersetMembership}[thm]{2}}}
  \proofstep{2}{\exists x\,(x\in X)}{%
    \FormulaRefAuto{A\neq\varnothing\eqvdash
      \exists x\,(x\in A)}[thm]{4}}
  \proofstep{6}{x\in X}{\rA}
  \proofstep{1,2,6}{f(x)\in f[X]}{%
    \FormulaRefAuto{C\subseteq A\dsep F\colon A\to B\dsep
      x\in C\vdash F(x)\in F[C]}[thm]{3,1,6}}
  \proofstep{1,2,6}{f[X]\neq\varnothing}{%
    \FormulaRefAuto{a\in A\vdash A\neq\varnothing}[thm]{7}}
  \proofstep{1,2}{f[X]\neq\varnothing}{\rEE{5,6,8}}
  \proofstep{1,2}{f[X]\subseteq B}{%
    \FormulaRefAuto{C\subseteq A\dsep F\colon A\to B
      \vdash F[C]\subseteq B}[thm]{3,1}}
  \proofstep{1,2}{f[X]\subseteq B\land f[X]\neq\varnothing}{%
    \rAI{10,9}}
  \proofstep{1,2}{f[X]\in\PowNE{B}}{%
    \FormulaRefAuto{NonemptyPowersetMembership}[thm]{11}}
\end{tabproof}

\FormulaThmDelta[Bilder liegen im Zielbereich]{%
  F\colon A\to B \vdash F[A]\subseteq B
}{
  \DeltaRow{Mengen}{A\dsep B}
  \DeltaRow{Funktionen}{F}
}
\begin{tabproof}
  \proofstep{1}{F\colon A\to B}{\rA}
  \proofstep{}{A\subseteq A}{\FormulaRefAuto{A\subseteq A}}
  \proofstep{1}{F[A]\subseteq B}{%
    \FormulaRefAuto{C\subseteq A \dsep F\colon A\to B \vdash F[C]\subseteq B}{2,1}}
\end{tabproof}

\FormulaThmDeltaK[Bilder in disjunkten Zielmengen sind disjunkt]{%
  F\colon A\to C\dsep G\colon B\to D\dsep C\cap D=\varnothing
  \vdash F[A]\cap G[B]=\varnothing
}{DisjointTargetImagesDisjoint}{%
  \DeltaRow{Mengen}{A\dsep B\dsep C\dsep D}
  \DeltaRow{Funktionen}{F\dsep G}
}
\begin{tabproof}
  \proofstep{1}{F\colon A\to C}{\rA}
  \proofstep{2}{G\colon B\to D}{\rA}
  \proofstep{3}{C\cap D=\varnothing}{\rA}
  \proofstep{1}{F[A]\subseteq C}{%
    \FormulaRefAuto{F\colon A\to B \vdash F[A]\subseteq B}{1}}
  \proofstep{2}{G[B]\subseteq D}{%
    \FormulaRefAuto{F\colon A\to B \vdash F[A]\subseteq B}{2}}
  \proofstep{1,2}{F[A]\cap G[B]\subseteq C\cap D}{%
    \FormulaRefAuto{A \subseteq B,\, C \subseteq D \vdash A \cap C \subseteq B \cap D}{4,5}}
  \proofstep{1,2,3}{F[A]\cap G[B]\subseteq\varnothing}{\rIE{3,6}}
  \proofstep{1,2,3}{F[A]\cap G[B]=\varnothing}{%
    \FormulaRefAuto{A\subseteq\varnothing\vdash A=\varnothing}{7}}
\end{tabproof}



\section{Urbild}
%%begin novalidate
\FormulaDefDelta[Urbildmenge]{%
F\colon A\to B\dsep C\subseteq B \vdash
F^{-1}[C] \coloneqq \{\,x\in A \mid F(x)\in C\,\}%
}{
  \DeltaRow{Mengen}{A\dsep B\dsep C\dsep x}
  \DeltaRow{Funktionen}{F}
}
%%end novalidate

\FormulaThmDelta{%
F\colon A\to B\dsep C\subseteq B
\vdash \bigl(x\in F^{-1}[C] \leftrightarrow
(x\in A \land F(x)\in C)\bigr)
}{
  \DeltaRow{Mengen}{A\dsep B\dsep C\dsep x}
  \DeltaRow{Funktionen}{F}
}
\begin{tabproof}
  \proofstep{1}{F\colon A\to B}{\rA}
  \proofstep{2}{C\subseteq B}{\rA}

  \proofstep{1,2}{F^{-1}[C]=\{\,u\in A \mid F(u)\in C\,\}}{%
    \FormulaRefAuto{F\colon A\to B\dsep C\subseteq B
    \vdash F^{-1}[C] \coloneqq \{\,u\in A \mid F(u)\in C\,\}}{1,2}}

  \proofstep{}{x\in \{\,u\in A \mid F(u)\in C\,\}\leftrightarrow (x\in A \land F(x)\in C)}{%
    \FormulaRefAuto{x\in \{u\in A \mid P(u)\}\eqvdash x\in A \land P(x)}}

  \proofstep{1,2}{x\in F^{-1}[C]\leftrightarrow (x\in A \land F(x)\in C)}{\rIE{3,4}}
\end{tabproof}

\FormulaThmDelta[Das Urbild der leeren Menge ist leer]{%
  F\colon A\to B \vdash F^{-1}[\varnothing]=\varnothing
}{%
  \DeltaRow{Mengen}{A\dsep B\dsep x}
  \DeltaRow{Funktionen}{F}
}
\begin{tabproof}
  \proofstep{1}{F\colon A\to B}{\rA}
  \proofstep{}{ \varnothing\subseteq B }{\FormulaRefAuto{\varnothing\subseteq A}}
  \proofstep{3}{x\in F^{-1}[\varnothing]}{\rA}
  \proofstep{1,2}{x\in F^{-1}[\varnothing]\leftrightarrow (x\in A \land F(x)\in \varnothing)}{%
    \FormulaRefAuto{F\colon A\to B\dsep C\subseteq B
    \vdash \bigl(x\in F^{-1}[C] \leftrightarrow
    (x\in A \land F(x)\in C)\bigr)}{1,2}}
  \proofstep{1,3}{x\in A \land F(x)\in \varnothing}{%
    \FormulaRefAuto{P \leftrightarrow Q, P \vdash Q}{4,3}}
  \proofstep{1,3}{F(x)\in \varnothing}{\rAEb{5}}
  \proofstep{}{F(x)\notin \varnothing}{\FormulaRefAuto{x \not\in \varnothing}}
  \proofstep{1,3}{x\in \varnothing}{\FormulaRefAuto{P\dsep \neg P \vdash Q}{6,7}}
  \proofstep{1}{ \forall x\,(x\in F^{-1}[\varnothing]\rightarrow x\in \varnothing)}{\rUI{\rRI{3,8}}}
  \proofstep{1}{F^{-1}[\varnothing]\subseteq \varnothing}{%
    \FormulaRefAuto{A\subseteq B \coloneq \forall x\,(x\in A \rightarrow x\in B)}{9}}
  \proofstep{1}{F^{-1}[\varnothing]=\varnothing}{%
    \FormulaRefAuto{A\subseteq\varnothing\vdash A=\varnothing}{10}}
\end{tabproof}

\FormulaThmDelta[Urbild einer Einermenge]{%
  F\colon A\to B\dsep y\in B
  \vdash F^{-1}[\{y\}] = \{\,x\in A \mid F(x)=y\,\}
}{%
  \DeltaRow{Mengen}{A\dsep B\dsep x\dsep y}
  \DeltaRow{Funktionen}{F}
}
\begin{tabproofsplitwide}
  \proofpartwideind[Teilmengenrichtung \(F^{-1}[\{y\}]\subseteq \{\,x\in A \mid F(x)=y\,\}\)]{%
    F\colon A\to B\dsep y\in B
    \vdash F^{-1}[\{y\}]\subseteq \{\,x\in A \mid F(x)=y\,\}
  }[%
    F\colon A\to B\dsep y\in B
    \vdash F^{-1}[\{y\}]\subseteq \{\,x\in A \mid F(x)=y\,\}
  ]
    \proofstepwidestar[1]{F\colon A\to B}{\rA}
    \proofstepwidestar[2]{y\in B}{\rA}
    \proofstepwidestar[2]{\{y\}\subseteq B}{\FormulaRefAuto{x\in A \vdash \{x\}\subseteq A}{2}}
    \proofstepwidestar[4]{x\in F^{-1}[\{y\}]}{\rA}
    \proofstepwidestar[1,3]{x\in F^{-1}[\{y\}]\leftrightarrow (x\in A \land F(x)\in \{y\})}{%
      \FormulaRefAuto{F\colon A\to B\dsep C\subseteq B
      \vdash \bigl(x\in F^{-1}[C] \leftrightarrow
      (x\in A \land F(x)\in C)\bigr)}{1,3}}
    \proofstepwidestar[1,3,4]{x\in A \land F(x)\in \{y\}}{%
      \FormulaRefAuto{P \leftrightarrow Q, P \vdash Q}{5,4}}
    \proofstepwidestar[1,3,4]{x\in A}{\rAEa{6}}
    \proofstepwidestar[1,3,4]{F(x)\in \{y\}}{\rAEb{6}}
    \proofstepwidestar[1,3,4]{F(x)=y}{\FormulaRefAuto{x\in \{a\}\eqvdash x=a}{8}}
    \proofstepwidestar[1,3,4]{x\in A \land F(x)=y}{\rAI{7,9}}
    \proofstepwidestar[1,3,4]{x\in \{\,x\in A \mid F(x)=y\,\}}{%
      \FormulaRefAuto{x\in \{u\in A \mid P(u)\}\eqvdash x\in A \land P(x)}{10}}
    \proofstepwidestar[1,2]{\forall x\,\bigl(x\in F^{-1}[\{y\}]\rightarrow x\in \{\,x\in A \mid F(x)=y\,\}\bigr)}{\rUI{\rRI{4,11}}}
    \proofstepwidestar[1,2]{F^{-1}[\{y\}]\subseteq \{\,x\in A \mid F(x)=y\,\}}{%
      \FormulaRefAuto{A\subseteq B \coloneq \forall x\,(x\in A \rightarrow x\in B)}{12}}
  \closeproofpartwideind

  \proofpartwideind[Teilmengenrichtung \(\{\,x\in A \mid F(x)=y\,\}\subseteq F^{-1}[\{y\}]\)]{%
    F\colon A\to B\dsep y\in B
    \vdash \{\,x\in A \mid F(x)=y\,\}\subseteq F^{-1}[\{y\}]
  }[%
    F\colon A\to B\dsep y\in B
    \vdash \{\,x\in A \mid F(x)=y\,\}\subseteq F^{-1}[\{y\}]
  ]
    \proofstepwidestar[1]{F\colon A\to B}{\rA}
    \proofstepwidestar[2]{y\in B}{\rA}
    \proofstepwidestar[2]{\{y\}\subseteq B}{\FormulaRefAuto{x\in A \vdash \{x\}\subseteq A}{2}}
    \proofstepwidestar[4]{x\in \{\,x\in A \mid F(x)=y\,\}}{\rA}
    \proofstepwidestar[4]{x\in A \land F(x)=y}{%
      \FormulaRefAuto{x\in \{u\in A \mid P(u)\}\eqvdash x\in A \land P(x)}{4}}
    \proofstepwidestar[4]{x\in A}{\rAEa{5}}
    \proofstepwidestar[4]{F(x)=y}{\rAEb{5}}
    \proofstepwidestar[4]{F(x)\in \{y\}}{\FormulaRefAuto{x\in \{a\}\eqvdash x=a}{7}}
    \proofstepwidestar[4]{x\in A \land F(x)\in \{y\}}{\rAI{6,8}}
    \proofstepwidestar[1,3]{x\in F^{-1}[\{y\}]\leftrightarrow (x\in A \land F(x)\in \{y\})}{%
      \FormulaRefAuto{F\colon A\to B\dsep C\subseteq B
      \vdash \bigl(x\in F^{-1}[C] \leftrightarrow
      (x\in A \land F(x)\in C)\bigr)}{1,3}}
    \proofstepwidestar[1,3,4]{x\in F^{-1}[\{y\}]}{\FormulaRefAuto{P \leftrightarrow Q, Q \vdash P}{10,9}}
    \proofstepwidestar[1,2]{\forall x\,\bigl(x\in \{\,x\in A \mid F(x)=y\,\}\rightarrow x\in F^{-1}[\{y\}]\bigr)}{\rUI{\rRI{4,11}}}
    \proofstepwidestar[1,2]{\{\,x\in A \mid F(x)=y\,\}\subseteq F^{-1}[\{y\}]}{%
      \FormulaRefAuto{A\subseteq B \coloneq \forall x\,(x\in A \rightarrow x\in B)}{12}}
  \closeproofpartwideind

  \proofpartwide{Zusammenfassung}
    \proofstepwidestar[1]{F\colon A\to B}{\rA}
    \proofstepwidestar[2]{y\in B}{\rA}
    \proofstepwidestar[1,2]{F^{-1}[\{y\}]\subseteq \{\,x\in A \mid F(x)=y\,\}}{%
      \FormulaRefAuto{F\colon A\to B\dsep y\in B
      \vdash F^{-1}[\{y\}]\subseteq \{\,x\in A \mid F(x)=y\,\}}{1,2}}
    \proofstepwidestar[1,2]{\{\,x\in A \mid F(x)=y\,\}\subseteq F^{-1}[\{y\}]}{%
      \FormulaRefAuto{F\colon A\to B\dsep y\in B
      \vdash \{\,x\in A \mid F(x)=y\,\}\subseteq F^{-1}[\{y\}]}{1,2}}
    \proofstepwidestar[1,2]{F^{-1}[\{y\}] = \{\,x\in A \mid F(x)=y\,\}}{%
      \FormulaRefAuto{A\subseteq B\dsep B\subseteq A \vdash A=B}{2,3}}
  \closeproofpartwide
\end{tabproofsplitwide}

\subsection{Urbilder liegen im Definitionsbereich}


\FormulaThmDelta[Urbilder liegen im Definitionsbereich]{%
  F\colon A\to B\dsep C\subseteq B \vdash F^{-1}[C]\subseteq A
}{
  \DeltaRow{Mengen}{A\dsep B\dsep C\dsep x}
  \DeltaRow{Funktionen}{F}
}
\begin{tabproof}
  \proofstep{1}{F\colon A\to B}{\rA}
  \proofstep{2}{C\subseteq B}{\rA}
  \proofstep{3}{x\in F^{-1}[C]}{\rA}
  \proofstep{1,2}{x\in F^{-1}[C]\leftrightarrow (x\in A \land F(x)\in C)}{%
    \FormulaRefAuto{F\colon A\to B\dsep C\subseteq B
    \vdash \bigl(x\in F^{-1}[C] \leftrightarrow
    (x\in A \land F(x)\in C)\bigr)}{1,2}}
  \proofstep{1,2,3}{x\in A \land F(x)\in C}{%
    \FormulaRefAuto{P \leftrightarrow Q, P \vdash Q}{4,3}}
  \proofstep{1,2,3}{x\in A}{\rAEa{5}}
  \proofstep{1,2}{\forall x\,\bigl(x\in F^{-1}[C]\rightarrow x\in A\bigr)}{\rUI{\rRI{3,6}}}
  \proofstep{1,2}{F^{-1}[C]\subseteq A}{%
    \FormulaRefAuto{A\subseteq B \coloneq \forall x\,(x\in A \rightarrow x\in B)}{7}}
\end{tabproof}

\FormulaThmDeltaK[Urbildgleichheit aus punktweiser Charakterisierung]{%
  \begin{aligned}[t]
    &F\colon A\to B\dsep T\subseteq B\dsep S\subseteq A\\
    &\dsep \forall x\in A\,\bigl(F(x)\in T\leftrightarrow x\in S\bigr)
    \vdash F^{-1}[T]=S
  \end{aligned}
}{PreimageEqFromPointwiseIff}{%
  \DeltaRow{Mengen}{A\dsep B\dsep S\dsep T\dsep x}
  \DeltaRow{Funktionen}{F}
}
\begin{tabproofsplitwide}
  \proofpartwideind[Teilmengenrichtung \(F^{-1}[T]\subseteq S\)]{%
    \begin{aligned}[t]
      &F\colon A\to B\dsep T\subseteq B\dsep S\subseteq A\\
      &\dsep \forall x\in A\,\bigl(F(x)\in T\leftrightarrow x\in S\bigr)
      \vdash F^{-1}[T]\subseteq S
    \end{aligned}
  }[%
    F\colon A\to B\dsep T\subseteq B\dsep S\subseteq A
    \dsep \forall x\in A\,\bigl(F(x)\in T\leftrightarrow x\in S\bigr)
    \vdash F^{-1}[T]\subseteq S
  ]
    \proofstepwidestar[1]{F\colon A\to B}{\rA}
    \proofstepwidestar[2]{T\subseteq B}{\rA}
    \proofstepwidestar[3]{S\subseteq A}{\rA}
    \proofstepwidestar[4]{%
      \forall x\in A\,\bigl(F(x)\in T\leftrightarrow x\in S\bigr)}{\rA}
    \proofstepwidestar[5]{x\in F^{-1}[T]}{\rA}
    \proofstepwidestar[1,2]{x\in F^{-1}[T]\leftrightarrow
      (x\in A\land F(x)\in T)}{%
      \FormulaRefAuto{F\colon A\to B\dsep C\subseteq B
      \vdash \bigl(x\in F^{-1}[C] \leftrightarrow
      (x\in A \land F(x)\in C)\bigr)}{1,2}}
    \proofstepwidestar[1,2,5]{x\in A\land F(x)\in T}{%
      \FormulaRefAuto{P \leftrightarrow Q, P \vdash Q}{6,5}}
    \proofstepwidestar[1,2,5]{x\in A}{\rAEa{7}}
    \proofstepwidestar[1,2,5]{F(x)\in T}{\rAEb{7}}
    \proofstepwidestar[4]{x\in A\rightarrow
      \bigl(F(x)\in T\leftrightarrow x\in S\bigr)}{\rUE{4}}
    \proofstepwidestar[1,2,4,5]{F(x)\in T\leftrightarrow x\in S}{%
      \rRE{8,10}}
    \proofstepwidestar[1,2,4,5]{x\in S}{%
      \FormulaRefAuto{P \leftrightarrow Q, P \vdash Q}{11,9}}
    \proofstepwidestar[1,2,4]{F^{-1}[T]\subseteq S}{%
      \FormulaRefAuto{A \subseteq B \coloneq
      \forall x\,(x\in A \rightarrow x\in B)}{\rUI{\rRI{5,12}}}}
  \closeproofpartwideind

  \proofpartwideind[Teilmengenrichtung \(S\subseteq F^{-1}[T]\)]{%
    \begin{aligned}[t]
      &F\colon A\to B\dsep T\subseteq B\dsep S\subseteq A\\
      &\dsep \forall x\in A\,\bigl(F(x)\in T\leftrightarrow x\in S\bigr)
      \vdash S\subseteq F^{-1}[T]
    \end{aligned}
  }[%
    F\colon A\to B\dsep T\subseteq B\dsep S\subseteq A
    \dsep \forall x\in A\,\bigl(F(x)\in T\leftrightarrow x\in S\bigr)
    \vdash S\subseteq F^{-1}[T]
  ]
    \proofstepwidestar[1]{F\colon A\to B}{\rA}
    \proofstepwidestar[2]{T\subseteq B}{\rA}
    \proofstepwidestar[3]{S\subseteq A}{\rA}
    \proofstepwidestar[4]{%
      \forall x\in A\,\bigl(F(x)\in T\leftrightarrow x\in S\bigr)}{\rA}
    \proofstepwidestar[5]{x\in S}{\rA}
    \proofstepwidestar[3,5]{x\in A}{%
      \FormulaRefAuto{A\subseteq B\dsep x\in A\vdash x\in B}{3,5}}
    \proofstepwidestar[4]{x\in A\rightarrow
      \bigl(F(x)\in T\leftrightarrow x\in S\bigr)}{\rUE{4}}
    \proofstepwidestar[3,4,5]{F(x)\in T\leftrightarrow x\in S}{%
      \rRE{6,7}}
    \proofstepwidestar[3,4,5]{F(x)\in T}{%
      \FormulaRefAuto{P \leftrightarrow Q, Q \vdash P}{8,5}}
    \proofstepwidestar[3,4,5]{x\in A\land F(x)\in T}{\rAI{6,9}}
    \proofstepwidestar[1,2]{x\in F^{-1}[T]\leftrightarrow
      (x\in A\land F(x)\in T)}{%
      \FormulaRefAuto{F\colon A\to B\dsep C\subseteq B
      \vdash \bigl(x\in F^{-1}[C] \leftrightarrow
      (x\in A \land F(x)\in C)\bigr)}{1,2}}
    \proofstepwidestar[1,2,3,4,5]{x\in F^{-1}[T]}{%
      \FormulaRefAuto{P \leftrightarrow Q, Q \vdash P}{11,10}}
    \proofstepwidestar[1,2,3,4]{S\subseteq F^{-1}[T]}{%
      \FormulaRefAuto{A \subseteq B \coloneq
      \forall x\,(x\in A \rightarrow x\in B)}{\rUI{\rRI{5,12}}}}
  \closeproofpartwideind

  \proofpartwide{Schluss}
    \proofstepwidestar[1]{F\colon A\to B}{\rA}
    \proofstepwidestar[2]{T\subseteq B}{\rA}
    \proofstepwidestar[3]{S\subseteq A}{\rA}
    \proofstepwidestar[4]{%
      \forall x\in A\,\bigl(F(x)\in T\leftrightarrow x\in S\bigr)}{\rA}
    \proofstepwidestar[1,2,3,4]{F^{-1}[T]\subseteq S}{%
      \FormulaRefAuto{F\colon A\to B\dsep T\subseteq B\dsep
      S\subseteq A\dsep
      \forall x\in A\,\bigl(F(x)\in T\leftrightarrow x\in S\bigr)
      \vdash F^{-1}[T]\subseteq S}{1,2,3,4}}
    \proofstepwidestar[1,2,3,4]{S\subseteq F^{-1}[T]}{%
      \FormulaRefAuto{F\colon A\to B\dsep T\subseteq B\dsep
      S\subseteq A\dsep
      \forall x\in A\,\bigl(F(x)\in T\leftrightarrow x\in S\bigr)
      \vdash S\subseteq F^{-1}[T]}{1,2,3,4}}
    \proofstepwidestar[1,2,3,4]{F^{-1}[T]=S}{%
      \FormulaRefAuto{A\subseteq B\dsep B\subseteq A \vdash A=B}{5,6}}
  \closeproofpartwide
\end{tabproofsplitwide}

\subsection{Rechenregeln für Urbilder}

\FormulaThmDelta[Monotonie des Urbilds]{%
  F\colon A\to B \dsep C\subseteq D \dsep D\subseteq B \vdash F^{-1}[C]\subseteq F^{-1}[D]
}{%
  \DeltaRow{Mengen}{A\dsep B\dsep C\dsep D\dsep x}
  \DeltaRow{Funktionen}{F}
}
\begin{tabproof}
  \proofstep{1}{F\colon A\to B}{\rA}
  \proofstep{2}{C\subseteq D}{\rA}
  \proofstep{3}{D\subseteq B}{\rA}
  \proofstep{2,3}{C\subseteq B}{\FormulaRefAuto{A\subseteq B\dsep B\subseteq C \vdash A\subseteq C}{2,3}}

  \proofstep{5}{x\in F^{-1}[C]}{\rA}
  \proofstep{1,4}{x\in F^{-1}[C]\leftrightarrow (x\in A \land F(x)\in C)}{%
    \FormulaRefAuto{F\colon A\to B\dsep C\subseteq B
    \vdash \bigl(x\in F^{-1}[C] \leftrightarrow
    (x\in A \land F(x)\in C)\bigr)}{1,4}}
  \proofstep{1,4,5}{x\in A \land F(x)\in C}{%
    \FormulaRefAuto{P \leftrightarrow Q, P \vdash Q}{6,5}}
  \proofstep{1,4,5}{x\in A}{\rAEa{7}}
  \proofstep{1,4,5}{F(x)\in C}{\rAEb{7}}
  \proofstep{1,2,4,5}{F(x)\in D}{\FormulaRefAuto{A\subseteq B\dsep x\in A\vdash x\in B}{2,9}}
  \proofstep{1,2,4,5}{x\in A \land F(x)\in D}{\rAI{8,10}}
  \proofstep{1,3}{x\in F^{-1}[D]\leftrightarrow (x\in A \land F(x)\in D)}{%
    \FormulaRefAuto{F\colon A\to B\dsep C\subseteq B
    \vdash \bigl(x\in F^{-1}[C] \leftrightarrow
    (x\in A \land F(x)\in C)\bigr)}{1,3}}
  \proofstep{1,2,3,5}{x\in F^{-1}[D]}{\FormulaRefAuto{P \leftrightarrow Q, Q \vdash P}{12,11}}

  \proofstep{1,2,3}{\forall x\,\bigl(x\in F^{-1}[C]\rightarrow x\in F^{-1}[D]\bigr)}{\rUI{\rRI{5,13}}}
  \proofstep{1,2,3}{F^{-1}[C]\subseteq F^{-1}[D]}{%
    \FormulaRefAuto{A\subseteq B \coloneq \forall x\,(x\in A \rightarrow x\in B)}{14}}
\end{tabproof}

\FormulaThmDelta[Element eines Urbilds einer Vereinigung]{%
  F\colon A\to B \dsep C\subseteq B \dsep D\subseteq B \dsep x\in F^{-1}[C\cup D] \vdash x\in F^{-1}[C]\cup F^{-1}[D]
}{%
  \DeltaRow{Mengen}{A\dsep B\dsep C\dsep D\dsep x}
  \DeltaRow{Funktionen}{F}
}
\begin{tabproof}
  \proofstep{1}{F\colon A\to B}{\rA}
  \proofstep{2}{C\subseteq B}{\rA}
  \proofstep{3}{D\subseteq B}{\rA}
  \proofstep{4}{x\in F^{-1}[C\cup D]}{\rA}
  \proofstep{2,3}{C\cup D\subseteq B}{\FormulaRefAuto{A\subseteq M\dsep B\subseteq M \vdash A\cup B\subseteq M}{2,3}}
  \proofstep{5}{x\in F^{-1}[C\cup D]\leftrightarrow (x\in A \land F(x)\in C\cup D)}{%
    \FormulaRefAuto{F\colon A\to B\dsep C\subseteq B
    \vdash \bigl(x\in F^{-1}[C] \leftrightarrow
    (x\in A \land F(x)\in C)\bigr)}{1,5}}
  \proofstep{4,5}{x\in A \land F(x)\in C\cup D}{%
    \FormulaRefAuto{P \leftrightarrow Q, P \vdash Q}{6,4}}
  \proofstep{4,5}{x\in A}{\rAEa{7}}
  \proofstep{4,5}{F(x)\in C\cup D}{\rAEb{7}}
  \proofstep{4,5}{F(x)\in C\lor F(x)\in D}{%
    \FormulaRefAuto{x\in A\cup B\eqvdash x\in A\lor x\in B}{9}}

  \proofstep{11}{F(x)\in C}{\rA}
  \proofstep{1,2,4,5,11}{x\in F^{-1}[C]\cup F^{-1}[D]}{%
    \FormulaRefAuto{y\in A\vdash y\in A\cup B}{%
      \FormulaRefAuto{P \leftrightarrow Q, Q \vdash P}{%
        \FormulaRefAuto{F\colon A\to B\dsep C\subseteq B
          \vdash \bigl(x\in F^{-1}[C] \leftrightarrow
          (x\in A \land F(x)\in C)\bigr)}{1,2},%
        \rAI{8,11}}}}

  \proofstep{13}{F(x)\in D}{\rA}
  \proofstep{1,3,4,5,13}{x\in F^{-1}[C]\cup F^{-1}[D]}{%
    \FormulaRefAuto{y\in B\vdash y\in A\cup B}{%
      \FormulaRefAuto{P \leftrightarrow Q, Q \vdash P}{%
        \FormulaRefAuto{F\colon A\to B\dsep C\subseteq B
          \vdash \bigl(x\in F^{-1}[C] \leftrightarrow
          (x\in A \land F(x)\in C)\bigr)}{1,3},%
        \rAI{8,13}}}}

  \proofstep{1,2,3,4}{x\in F^{-1}[C]\cup F^{-1}[D]}{\rOE{10,11,12,13,14}}
\end{tabproof}

\FormulaThmDelta[Urbild einer Vereinigung]{%
  F\colon A\to B \dsep C\subseteq B \dsep D\subseteq B \vdash F^{-1}[C\cup D]=F^{-1}[C]\cup F^{-1}[D]
}{%
  \DeltaRow{Mengen}{A\dsep B\dsep C\dsep D\dsep x}
  \DeltaRow{Funktionen}{F}
}
\begin{tabproofsplitwide}

  \proofpartwideind[Teilmengenrichtung \(F^{-1}[C\cup D]\subseteq F^{-1}[C]\cup F^{-1}[D]\)]{%
    F\colon A\to B \dsep C\subseteq B \dsep D\subseteq B \vdash F^{-1}[C\cup D]\subseteq F^{-1}[C]\cup F^{-1}[D]
  }[%
    F\colon A\to B \dsep C\subseteq B \dsep D\subseteq B \vdash F^{-1}[C\cup D]\subseteq F^{-1}[C]\cup F^{-1}[D]
  ]
    \proofstepwidestar[1]{F\colon A\to B}{\rA}
    \proofstepwidestar[2]{C\subseteq B}{\rA}
    \proofstepwidestar[3]{D\subseteq B}{\rA}
    \proofstepwidestar[4]{x\in F^{-1}[C\cup D]}{\rA}
    \proofstepwidestar[1,2,3,4]{x\in F^{-1}[C]\cup F^{-1}[D]}{%
      \FormulaRefAuto{F\colon A\to B \dsep C\subseteq B \dsep D\subseteq B \dsep x\in F^{-1}[C\cup D] \vdash x\in F^{-1}[C]\cup F^{-1}[D]}{1,2,3,4}}
    \proofstepwidestar[1,2,3]{\forall x\,\bigl(x\in F^{-1}[C\cup D]\rightarrow x\in F^{-1}[C]\cup F^{-1}[D]\bigr)}{\rUI{\rRI{4,5}}}
    \proofstepwidestar[1,2,3]{F^{-1}[C\cup D]\subseteq F^{-1}[C]\cup F^{-1}[D]}{%
      \FormulaRefAuto{A\subseteq B \coloneq \forall x\,(x\in A \rightarrow x\in B)}{6}}
  \closeproofpartwideind

\proofpartwideind[Teilmengenrichtung \(F^{-1}[C]\cup F^{-1}[D]\subseteq F^{-1}[C\cup D]\)]{%
  F\colon A\to B \dsep C\subseteq B \dsep D\subseteq B \vdash F^{-1}[C]\cup F^{-1}[D]\subseteq F^{-1}[C\cup D]
}[%
  F\colon A\to B \dsep C\subseteq B \dsep D\subseteq B \vdash F^{-1}[C]\cup F^{-1}[D]\subseteq F^{-1}[C\cup D]
]
  \proofstepwidestar[1]{F\colon A\to B}{\rA}
  \proofstepwidestar[2]{C\subseteq B}{\rA}
  \proofstepwidestar[3]{D\subseteq B}{\rA}
  \proofstepwidestar[2]{C\subseteq C\cup D}{\FormulaRefAuto{A\subseteq A\cup B}{2}}
  \proofstepwidestar[3]{D\subseteq C\cup D}{\FormulaRefAuto{B\subseteq A\cup B}{3}}
  \proofstepwidestar[2,3]{C\cup D\subseteq B}{%
    \FormulaRefAuto{A\subseteq M\dsep B\subseteq M \vdash A\cup B\subseteq M}{2,3}}
  \proofstepwidestar[1,2,3]{F^{-1}[C]\subseteq F^{-1}[C\cup D]}{%
    \FormulaRefAuto{F\colon A\to B \dsep C\subseteq D \dsep D\subseteq B \vdash F^{-1}[C]\subseteq F^{-1}[D]}{1,4,6}}
  \proofstepwidestar[1,2,3]{F^{-1}[D]\subseteq F^{-1}[C\cup D]}{%
    \FormulaRefAuto{F\colon A\to B \dsep C\subseteq D \dsep D\subseteq B \vdash F^{-1}[C]\subseteq F^{-1}[D]}{1,5,6}}
  \proofstepwidestar[1,2,3]{F^{-1}[C]\cup F^{-1}[D]\subseteq F^{-1}[C\cup D]}{%
    \FormulaRefAuto{X\subseteq Z\dsep Y\subseteq Z \vdash X\cup Y\subseteq Z}{7,8}}
\closeproofpartwideind

  \proofpartwide{Zusammenfassung}
    \proofstepwidestar[1]{F\colon A\to B}{\rA}
    \proofstepwidestar[2]{C\subseteq B}{\rA}
    \proofstepwidestar[3]{D\subseteq B}{\rA}
    \proofstepwidestar[1,2,3]{F^{-1}[C\cup D]\subseteq F^{-1}[C]\cup F^{-1}[D]}{%
      \FormulaRefAuto{F\colon A\to B \dsep C\subseteq B \dsep D\subseteq B \vdash F^{-1}[C\cup D]\subseteq F^{-1}[C]\cup F^{-1}[D]}{1,2,3}}
    \proofstepwidestar[1,2,3]{F^{-1}[C]\cup F^{-1}[D]\subseteq F^{-1}[C\cup D]}{%
      \FormulaRefAuto{F\colon A\to B \dsep C\subseteq B \dsep D\subseteq B \vdash F^{-1}[C]\cup F^{-1}[D]\subseteq F^{-1}[C\cup D]}{1,2,3}}
    \proofstepwidestar[1,2,3]{F^{-1}[C\cup D]=F^{-1}[C]\cup F^{-1}[D]}{%
      \FormulaRefAuto{A\subseteq B\dsep B\subseteq A \vdash A=B}{4,5}}
  \closeproofpartwide

\end{tabproofsplitwide}

\FormulaThmDelta[Urbild eines Durchschnitts liegt im Durchschnitt der Urbilder]{%
  F\colon A\to B \dsep C\subseteq B \dsep D\subseteq B \vdash F^{-1}[C\cap D]\subseteq F^{-1}[C]\cap F^{-1}[D]
}{%
  \DeltaRow{Mengen}{A\dsep B\dsep C\dsep D\dsep x}
  \DeltaRow{Funktionen}{F}
}
\begin{tabproof}
  \proofstep{1}{F\colon A\to B}{\rA}
  \proofstep{2}{C\subseteq B}{\rA}
  \proofstep{3}{D\subseteq B}{\rA}
  \proofstep{}{C\cap D\subseteq C}{\FormulaRefAuto{A\cap B\subseteq A}}
  \proofstep{}{C\cap D\subseteq D}{\FormulaRefAuto{A\cap B\subseteq B}}
  \proofstep{1,2}{F^{-1}[C\cap D]\subseteq F^{-1}[C]}{%
    \FormulaRefAuto{F\colon A\to B \dsep C\subseteq D \dsep D\subseteq B \vdash F^{-1}[C]\subseteq F^{-1}[D]}{1,4,2}}
  \proofstep{1,3}{F^{-1}[C\cap D]\subseteq F^{-1}[D]}{%
    \FormulaRefAuto{F\colon A\to B \dsep C\subseteq D \dsep D\subseteq B \vdash F^{-1}[C]\subseteq F^{-1}[D]}{1,5,3}}

  \proofstep{8}{x\in F^{-1}[C\cap D]}{\rA}
  \proofstep{6,8}{x\in F^{-1}[C]}{\FormulaRefAuto{A\subseteq B\dsep x\in A\vdash x\in B}{6,8}}
  \proofstep{7,8}{x\in F^{-1}[D]}{\FormulaRefAuto{A\subseteq B\dsep x\in A\vdash x\in B}{7,8}}
  \proofstep{6,7,8}{x\in F^{-1}[C]\land x\in F^{-1}[D]}{\rAI{9,10}}
  \proofstep{6,7,8}{x\in F^{-1}[C]\cap F^{-1}[D]}{%
    \FormulaRefAuto{x\in A\cap B\eqvdash x\in A\land x\in B}{11}}
  \proofstep{1,2,3}{\forall x\,\bigl(x\in F^{-1}[C\cap D]\rightarrow x\in F^{-1}[C]\cap F^{-1}[D]\bigr)}{\rUI{\rRI{8,12}}}
  \proofstep{1,2,3}{F^{-1}[C\cap D]\subseteq F^{-1}[C]\cap F^{-1}[D]}{%
    \FormulaRefAuto{A\subseteq B \coloneq \forall x\,(x\in A \rightarrow x\in B)}{13}}
\end{tabproof}

\FormulaThmDelta[Gemeinsame Urbildwerte liegen im Urbild des Durchschnitts]{%
  F\colon A\to B \dsep C\subseteq B \dsep D\subseteq B \dsep x\in F^{-1}[C] \dsep x\in F^{-1}[D] \vdash x\in F^{-1}[C\cap D]
}{%
  \DeltaRow{Mengen}{A\dsep B\dsep C\dsep D\dsep x}
  \DeltaRow{Funktionen}{F}
}
\begin{tabproof}
  \proofstep{1}{F\colon A\to B}{\rA}
  \proofstep{2}{C\subseteq B}{\rA}
  \proofstep{3}{D\subseteq B}{\rA}
  \proofstep{4}{x\in F^{-1}[C]}{\rA}
  \proofstep{5}{x\in F^{-1}[D]}{\rA}

  \proofstep{1,2}{x\in F^{-1}[C]\leftrightarrow (x\in A \land F(x)\in C)}{%
    \FormulaRefAuto{F\colon A\to B\dsep C\subseteq B
    \vdash \bigl(x\in F^{-1}[C] \leftrightarrow
    (x\in A \land F(x)\in C)\bigr)}{1,2}}
  \proofstep{1,2,4}{x\in A \land F(x)\in C}{%
    \FormulaRefAuto{P \leftrightarrow Q, P \vdash Q}{6,4}}
  \proofstep{1,2,4}{x\in A}{\rAEa{7}}
  \proofstep{1,2,4}{F(x)\in C}{\rAEb{7}}

  \proofstep{1,3}{x\in F^{-1}[D]\leftrightarrow (x\in A \land F(x)\in D)}{%
    \FormulaRefAuto{F\colon A\to B\dsep C\subseteq B
    \vdash \bigl(x\in F^{-1}[C] \leftrightarrow
    (x\in A \land F(x)\in C)\bigr)}{1,3}}
  \proofstep{1,3,5}{x\in A \land F(x)\in D}{%
    \FormulaRefAuto{P \leftrightarrow Q, P \vdash Q}{10,5}}
  \proofstep{1,3,5}{F(x)\in D}{\rAEb{11}}

  \proofstep{1,2,3,4,5}{F(x)\in C \land F(x)\in D}{\rAI{9,12}}
  \proofstep{1,2,3,4,5}{F(x)\in C\cap D}{%
    \FormulaRefAuto{x\in A\cap B\eqvdash x\in A\land x\in B}{13}}
  \proofstep{3}{C\cap D\subseteq B}{%
    \FormulaRefAuto{B\subseteq M \vdash A\cap B\subseteq M}{3}}
  \proofstep{1,2,3,4,5}{x\in A \land F(x)\in C\cap D}{\rAI{8,14}}
  \proofstep{1,15}{x\in F^{-1}[C\cap D]\leftrightarrow (x\in A \land F(x)\in C\cap D)}{%
    \FormulaRefAuto{F\colon A\to B\dsep C\subseteq B
    \vdash \bigl(x\in F^{-1}[C] \leftrightarrow
    (x\in A \land F(x)\in C)\bigr)}{1,15}}
  \proofstep{1,2,3,4,5}{x\in F^{-1}[C\cap D]}{\FormulaRefAuto{P \leftrightarrow Q, Q \vdash P}{17,16}}
\end{tabproof}

\FormulaThmDelta[Urbild eines Durchschnitts]{%
  F\colon A\to B \dsep C\subseteq B \dsep D\subseteq B \vdash F^{-1}[C\cap D]=F^{-1}[C]\cap F^{-1}[D]
}{%
  \DeltaRow{Mengen}{A\dsep B\dsep C\dsep D\dsep x}
  \DeltaRow{Funktionen}{F}
}
\begin{tabproofsplitwide}

  \proofpartwideind[Teilmengenrichtung \(F^{-1}[C]\cap F^{-1}[D]\subseteq F^{-1}[C\cap D]\)]{%
    F\colon A\to B \dsep C\subseteq B \dsep D\subseteq B \vdash F^{-1}[C]\cap F^{-1}[D]\subseteq F^{-1}[C\cap D]
  }[%
    F\colon A\to B \dsep C\subseteq B \dsep D\subseteq B \vdash F^{-1}[C]\cap F^{-1}[D]\subseteq F^{-1}[C\cap D]
  ]
    \proofstepwidestar[1]{F\colon A\to B}{\rA}
    \proofstepwidestar[2]{C\subseteq B}{\rA}
    \proofstepwidestar[3]{D\subseteq B}{\rA}
    \proofstepwidestar[4]{x\in F^{-1}[C]\cap F^{-1}[D]}{\rA}
    \proofstepwidestar[4]{x\in F^{-1}[C]\land x\in F^{-1}[D]}{%
      \FormulaRefAuto{x\in A\cap B\eqvdash x\in A\land x\in B}{4}}
    \proofstepwidestar[4]{x\in F^{-1}[C]}{\rAEa{5}}
    \proofstepwidestar[4]{x\in F^{-1}[D]}{\rAEb{5}}
    \proofstepwidestar[1,2,3,6,7]{x\in F^{-1}[C\cap D]}{%
      \FormulaRefAuto{F\colon A\to B \dsep C\subseteq B \dsep D\subseteq B \dsep x\in F^{-1}[C] \dsep x\in F^{-1}[D] \vdash x\in F^{-1}[C\cap D]}{1,2,3,6,7}}
    \proofstepwidestar[1,2,3]{\forall x\,\bigl(x\in F^{-1}[C]\cap F^{-1}[D]\rightarrow x\in F^{-1}[C\cap D]\bigr)}{\rUI{\rRI{4,8}}}
    \proofstepwidestar[1,2,3]{F^{-1}[C]\cap F^{-1}[D]\subseteq F^{-1}[C\cap D]}{%
      \FormulaRefAuto{A\subseteq B \coloneq \forall x\,(x\in A \rightarrow x\in B)}{9}}
  \closeproofpartwideind

  \proofpartwide{Zusammenfassung}
    \proofstepwidestar[1]{F\colon A\to B}{\rA}
    \proofstepwidestar[2]{C\subseteq B}{\rA}
    \proofstepwidestar[3]{D\subseteq B}{\rA}
    \proofstepwidestar[1,2,3]{F^{-1}[C\cap D]\subseteq F^{-1}[C]\cap F^{-1}[D]}{%
      \FormulaRefAuto{F\colon A\to B \dsep C\subseteq B \dsep D\subseteq B \vdash F^{-1}[C\cap D]\subseteq F^{-1}[C]\cap F^{-1}[D]}{1,2,3}}
    \proofstepwidestar[1,2,3]{F^{-1}[C]\cap F^{-1}[D]\subseteq F^{-1}[C\cap D]}{%
      \FormulaRefAuto{F\colon A\to B \dsep C\subseteq B \dsep D\subseteq B \vdash F^{-1}[C]\cap F^{-1}[D]\subseteq F^{-1}[C\cap D]}{1,2,3}}
    \proofstepwidestar[1,2,3]{F^{-1}[C\cap D]=F^{-1}[C]\cap F^{-1}[D]}{%
      \FormulaRefAuto{A\subseteq B\dsep B\subseteq A \vdash A=B}{4,5}}
  \closeproofpartwide

\end{tabproofsplitwide}

\section{Graph einer Funktion}

\FormulaThmDeltaK[Graph mit eindeutigen Werten ist eine Funktion]{%
  R\subseteq A\times B\dsep
  \forall x\in A\,\exists! y\in B\,\bigl((x,y)\in R\bigr)
  \vdash R\colon A\to B
}{UniqueValuedGraphFunction}{%
  \DeltaDecl{Mengen}{A\dsep B\dsep R\dsep x\dsep y\dsep z}%
}
\begin{tabproofsplitwide}
  \proofpartwideindR[Funktionale Eindeutigkeit]{%
    \begin{aligned}[t]
      &R\subseteq A\times B\dsep
      \forall x\in A\,\exists! y\in B\,\bigl((x,y)\in R\bigr)\dsep{}\\
      &(x,y)\in R\dsep (x,z)\in R\vdash y=z
    \end{aligned}
  }{%
    R\subseteq A\times B\dsep
    \forall x\in A\,\exists! y\in B\,\bigl((x,y)\in R\bigr)\dsep
    (x,y)\in R\dsep (x,z)\in R\vdash y=z
  }
    \proofstepwidestar[1]{R\subseteq A\times B}{\rA}
    \proofstepwidestar[2]{%
      \forall x\in A\,\exists! y\in B\,\bigl((x,y)\in R\bigr)
    }{\rA}
    \proofstepwidestar[3]{(x,y)\in R}{\rA}
    \proofstepwidestar[4]{(x,z)\in R}{\rA}
    \proofstepwidestar[1,3]{(x,y)\in A\times B}{%
      \FormulaRefAuto{A\subseteq B\dsep x\in A\vdash x\in B}{1,3}}
    \proofstepwidestar[1,3]{x\in A}{%
      \FormulaRefAuto{(a,b)\in A\times B\vdash a\in A}{5}}
    \proofstepwidestar[1,3]{y\in B}{%
      \FormulaRefAuto{(a,b)\in A\times B\vdash b\in B}{5}}
    \proofstepwidestar[1,4]{(x,z)\in A\times B}{%
      \FormulaRefAuto{A\subseteq B\dsep x\in A\vdash x\in B}{1,4}}
    \proofstepwidestar[1,4]{z\in B}{%
      \FormulaRefAuto{(a,b)\in A\times B\vdash b\in B}{8}}
    \proofstepwidestar[1,2,3]{%
      \exists! y\in B\,\bigl((x,y)\in R\bigr)
    }{\rRE{\rUE{2},6}}
    \proofstepwidestar[1,2,3,4]{y=z}{%
      \UEE{10,7,9,3,4}}
  \closeproofpartwideindR

  \proofpartwideindR[Totalität]{%
    R\subseteq A\times B\dsep
    \forall x\in A\,\exists! y\in B\,\bigl((x,y)\in R\bigr)\dsep
    x\in A\vdash \exists y\,\bigl((x,y)\in R\bigr)
  }{%
    R\subseteq A\times B\dsep
    \forall x\in A\,\exists! y\in B\,\bigl((x,y)\in R\bigr)\dsep
    x\in A\vdash \exists y\,\bigl((x,y)\in R\bigr)
  }
    \proofstepwidestar[1]{R\subseteq A\times B}{\rA}
    \proofstepwidestar[2]{%
      \forall x\in A\,\exists! y\in B\,\bigl((x,y)\in R\bigr)
    }{\rA}
    \proofstepwidestar[3]{x\in A}{\rA}
    \proofstepwidestar[2,3]{%
      \exists! y\in B\,\bigl((x,y)\in R\bigr)
    }{\rRE{\rUE{2},3}}
    \proofstepwidestar[2,3]{\exists y\,\bigl((x,y)\in R\bigr)}{\UEE{4}}
  \closeproofpartwideindR



  \proofpartwide{Schluss}
    \proofstepwidestar[1]{R\subseteq A\times B}{\rA}
    \proofstepwidestar[2]{%
      \forall x\in A\,\exists! y\in B\,\bigl((x,y)\in R\bigr)
    }{\rA}
    \proofstepwidestar[1,2]{R\colon A\to B}{%
      \FormulaRefAuto{Funktion}{%
        1,
        \FormulaRefAuto{R\subseteq A\times B\dsep
          \forall x\in A\,\exists! y\in B\,\bigl((x,y)\in R\bigr)\dsep
          (x,y)\in R\dsep (x,z)\in R\vdash y=z}{1,2},
        \FormulaRefAuto{R\subseteq A\times B\dsep
          \forall x\in A\,\exists! y\in B\,\bigl((x,y)\in R\bigr)\dsep
          x\in A\vdash \exists y\,\bigl((x,y)\in R\bigr)}{1,2}%
      }}
  \closeproofpartwide
\end{tabproofsplitwide}

\FormulaDefDelta[Graph einer Funktion]{\Graph(F)\coloneq\{(x,y) \in A \times B \mid y=F(x)\}}%
{
\DeltaRow{Mengen}{x\dsep y\dsep A\dsep B}
\DeltaRow{Funktionensymbol}{F}
}

\FormulaThmDelta{\Graph(F)\subseteq A\times B}%
{
\DeltaRow{Mengen}{x\dsep y\dsep A\dsep B}
\DeltaRow{Funktionensymbol}{F}
}
\begin{tabproof}
    \proofstep{1}{\Graph(F)\coloneq\{(x,y) \in A \times B \mid y=F(x)\}}{\FormulaRefAuto{\Graph(F)\coloneq\{(x,y) \in A \times B \mid y=F(x)\}}}
    \proofstep{1}{\{(x,y) \in A \times B \mid y=F(x)\}\subseteq A\times B}{\FormulaRefAuto{\{ x \in A \mid P(x)\} \subseteq A}}
    \proofstep{1}{\Graph(F)\subseteq A\times B}{\rIE{1,2}}
\end{tabproof}

\FormulaThmDelta{(x,y)\in \Graph(F)\vdash y=F(x)}%
{
\DeltaRow{Mengen}{x\dsep y\dsep A\dsep B}
\DeltaRow{Funktionensymbol}{F}
}
\begin{tabproof}
  \proofstep{1}{(x,y)\in \Graph(F)}{\rA}
  \proofstep{}{\Graph(F)=\{(x,y) \in A \times B \mid y=F(x)\}}{\FormulaRefAuto{\Graph(F)\coloneq\{(x,y) \in A \times B \mid y=F(x)\}}}
  \proofstep{1}{y=F(x)}{\FormulaRefAuto{x \in A\dsep A=\{x \in B \mid P(x)\}\vdash P(x)}{1,2}}
\end{tabproof}

\FormulaThmDelta[Funktionale Eindeutigkeit von \(\Graph(F)\)]{(x,y)\in \Graph(F)\dsep (x,z)\in \Graph(F)\vdash y=z}%
{
\DeltaRow{Mengen}{x\dsep y\dsep z}
\DeltaRow{Funktionensymbol}{F}
}
\begin{tabproof}
  \proofstep{1}{(x,y)\in \Graph(F)}{\rA}
  \proofstep{2}{(x,z)\in \Graph(F)}{\rA}
  \proofstep{1}{y = F(x)}{\FormulaRefAuto{(x,y)\in Graph(F)\vdash y=F(x)}{1}}
  \proofstep{2}{z = F(x)}{\FormulaRefAuto{(x,y)\in Graph(F)\vdash y=F(x)}{2}}
  \proofstep{1,2}{y = z}{\FormulaRefAuto{a = b,\, c = b \vdash a = c}{3,4}}
\end{tabproof}

\FormulaThmDelta
{\forall u\in A\,F(u)\in B\dsep x\in A \vdash \exists y\,(x,y)\in\Graph(F)}
{
\DeltaRow{Mengen}{u\dsep x\dsep y}
\DeltaRow{Funktionensymbol}{F}
}
\begin{tabproof}
  \proofstep{1}{\forall u\in A\,F(u)\in B}{\rA}
  \proofstep{2}{x\in A}{\rA}
  \proofstep{1,2}{F(x)\in B}{\rRE{\rUE{1},2}}
  \proofstep{1,2}{(x,F(x))\in A\times B}{\FormulaRefAuto{a\in A\dsep b\in B\vdash (a,b)\in A\times B}{2,3}}
  \proofstep{}{F(x)=F(x)}{\rII}
  \proofstep{1,2}{(x,F(x))\in\{(x,y)\in A \times B\mid y=F(x)\}}{\FormulaRefAuto{x \in A\dsep P(x)\vdash x \in \{x \in A \mid P(x)\}}{4,5}}
  \proofstep{1,2}{(x,F(x))\in\Graph(F)}{\rIE{\FormulaRefAuto{\Graph(F)\coloneq\{(x,y) \in A \times B \mid y=F(x)\}},6}}
  \proofstep{1,2}{\exists y\,(x,y)\in\Graph(F)}{\rEI{7}}
\end{tabproof}

\FormulaThmDelta[Graph eines Funktionensymbols ist eine Funktion]
{\forall u\in A\,F(u)\in B \vdash \Graph(F)\colon A\to B}
{
  \DeltaRow{Mengen}{A\dsep B}
  \DeltaRow{Funktionensymbol}{F}
}
\begin{tabproof}
  \proofstep{1}{\forall u\in A\,F(u)\in B}{\rA}
  \proofstep{}{\Graph(F) \subseteq A\times B }%
    {\FormulaRefAuto{\Graph(F)\subseteq A\times B}}
  \proofstep{}{\forall (x,y),(x,z)\in \Graph(F)\, y=z }%
    {\FormulaRefAuto{(x,y)\in \Graph(F)\dsep (x,z)\in \Graph(F)\vdash y=z}}
  \proofstep{1}{\forall x\in A \exists y\,(x,y)\in \Graph(F)}%
    {\FormulaRefAuto{\forall u\in A\,F(u)\in B\dsep x\in A \vdash \exists y\,(x,y)\in\Graph(F)}{1}}
  \proofstep{1}{\Graph(F)\colon A\to B}%
    {\FormulaRefAuto{Funktion}{2,3,4}}
\end{tabproof}

\FormulaThmDelta[Funktionswerte von \(\Graph(F)\)]
{\forall u\in A\,F(u)\in B \vdash \forall x\in A\,\Graph(F)(x) = F(x)}
{
  \DeltaRow{Mengen}{A\dsep B\dsep x}
  \DeltaRow{Funktionensymbol}{F}
}
\begin{tabproof}
  \proofstep{1}{\forall u\in A\,F(u)\in B}{\rA}
  \proofstep{2}{x\in A}{\rA}
  \proofstep{1}{\Graph(F)\colon A\to B}{\FormulaRefAuto{\forall u\in A\,F(u)\in B \vdash \Graph(F)\colon A\to B}{1}}
  \proofstep{1,2}{(x,\Graph(F)(x))\in\Graph(F)}{%
    \FormulaRefAuto{F\colon A\to B\dsep x\in A\vdash (x,F(x))\in F}{3,2}}
  \proofstep{1,2}{\Graph(F)(x) = F(x)}{\FormulaRefAuto{(x,y)\in \Graph(F)\vdash y=F(x)}{4}}
  \proofstep{1}{\forall x\in A\,\Graph(F)(x) = F(x)}{\rUI{\rRI{2,5}}}
\end{tabproof}

\FormulaThmDelta[Graph einer Funktion ist die Funktion]{%
  F\colon A\to B \vdash \Graph(F)=F%
}{%
  \DeltaDecl{Mengen}{A\dsep B\dsep x}%
  \DeltaDecl{Funktionen}{F}%
}
\begin{tabproof}
  \proofstep{1}{F\colon A\to B}{\rA}

  \proofstep{2}{x\in A}{\rA}
  \proofstep{1,2}{F(x)\in B}{%
    \FormulaRefAuto{F\colon A\to B\dsep x\in A\vdash F(x)\in B}{1,2}}
  \proofstep{1}{\forall x\in A\,F(x)\in B}{\rUI{\rRI{2,3}}}

  \proofstep{1}{\Graph(F)\colon A\to B}{\FormulaRefAuto{\forall u\in A\,F(u)\in B \vdash \Graph(F)\colon A\to B}{4}}


  \proofstep{1}{\forall x\in A\,\bigl(\Graph(F)(x)=F(x)\bigr)}{\FormulaRefAuto{\forall u\in A\,F(u)\in B \vdash \forall x\in A\,\Graph(F)(x) = F(x)}{4}}

  \proofstep{1}{\Graph(F)=F}{%
    \FormulaRefAuto{FunctionExtensionality}[thm]{5,1,6}}
\end{tabproof}

\subsection{Direkte Funktionskonstruktion aus typisierten Termen}

Sei \(y\) ein bereits gebildeter Term, in dem \(x\) frei vorkommen darf.
In den beiden beschränkten Quantoren des folgenden Satzes bindet \(\forall x\)
auch diese Vorkommen von \(x\) in \(y\). Die Voraussetzung ist also die
universelle Abschließung der Typisierungsimplikation und nicht die bloße
offene Formel \(x\in A\rightarrow y\in B\).

\FormulaThmDeltaK[Funktion zu einem typisierten Term]{%
  \forall x\in A\,y\in B
  \vdash
  \exists!F\,\Bigl(
    F\colon A\to B\land
    \forall x\in A\,F(x)=y
  \Bigr)%
}{TypedFunctionDefinitionPrinciple}{%
  \DeltaRow{Mengen}{A\dsep B\dsep F\dsep G\dsep x\dsep z\dsep w}
  \DeltaRow{Term}{y\text{ (mit ausgezeichnetem }x\text{)}}
  \DeltaRow{Frisch f\"ur \(y\)}{F\dsep G\dsep z\dsep w}
}

Im Beweis steht \(R_y\) lediglich als Abkürzung für den Mengenterm
\[
  \bigl\{(x,z)\in A\times B\mid z=y\bigr\},
\]
wobei \(x\) und das frische \(z\) innerhalb der Aussonderung gebunden sind.
Damit wird insbesondere kein zusätzliches Funktionensymbol eingeführt.

\begin{tabproofsplitwide}
  \proofpartwideindR[Graphkonstruktion]{%
    \forall x\in A\,y\in B\vdash R_y\colon A\to B
  }{%
    \forall x\in A\,y\in B\vdash R_y\colon A\to B
  }[TypedTermGraphFunction]
    \proofstepwidestar[1]{\forall x\in A\,y\in B}{\rA}
    \proofstepwidestar[]{R_y\subseteq A\times B}{%
      \FormulaRefAuto{\{x\in A\mid P(x)\}\subseteq A}}

    \proofstepwidestar[3]{x\in A}{\rA}
    \proofstepwidestar[1,3]{y\in B}{\rRE{\rUE{1},3}}
    \proofstepwidestar[1,3]{(x,y)\in A\times B}{%
      \FormulaRefAuto{a\in A\dsep b\in B\vdash (a,b)\in A\times B}{3,4}}
    \proofstepwidestar[]{y=y}{\rII}
    \proofstepwidestar[1,3]{(x,y)\in R_y}{%
      \FormulaRefAuto{x\in A\dsep P(x)\vdash
        x\in\{x\in A\mid P(x)\}}{5,6}}
    \proofstepwidestar[1,3]{y\in B\land (x,y)\in R_y}{\rAI{4,7}}
    \proofstepwidestar[1,3]{%
      \exists z\,\bigl(z\in B\land (x,z)\in R_y\bigr)
    }{\rEI{8}}

    \proofstepwidestar[10]{z\in B\land (x,z)\in R_y}{\rA}
    \proofstepwidestar[11]{w\in B\land (x,w)\in R_y}{\rA}
    \proofstepwidestar[10]{(x,z)\in R_y}{\rAEb{10}}
    \proofstepwidestar[11]{(x,w)\in R_y}{\rAEb{11}}
    \proofstepwidestar[10]{z=y}{%
      \rAEb{\FormulaRefAuto{x\in\{u\in A\mid P(u)\}
        \eqvdash x\in A\land P(x)}{12}}}
    \proofstepwidestar[11]{w=y}{%
      \rAEb{\FormulaRefAuto{x\in\{u\in A\mid P(u)\}
        \eqvdash x\in A\land P(x)}{13}}}
    \proofstepwidestar[10,11]{z=w}{%
      \FormulaRefAuto{a=b\dsep c=b\vdash a=c}{14,15}}
    \proofstepwidestar[1,3]{%
      \exists!z\in B\,\bigl((x,z)\in R_y\bigr)
    }{\UEI{9,10,11,16}}
    \proofstepwidestarbreakkeep[1]{%
      \forall x\in A\,\exists!z\in B\,\bigl((x,z)\in R_y\bigr)
    }{\rUI{\rRI{3,17}}}
    \proofstepwidestar[1]{R_y\colon A\to B}{%
      \FormulaRefAuto{UniqueValuedGraphFunction}[thm]{2,18}}
  \closeproofpartwideindR

  \proofpartwideindR[Funktionswerte]{%
    \forall x\in A\,y\in B\vdash
    \forall x\in A\,R_y(x)=y
  }{%
    \forall x\in A\,y\in B\vdash
    \forall x\in A\,R_y(x)=y
  }[TypedTermGraphValues]
    \proofstepwidestar[1]{\forall x\in A\,y\in B}{\rA}
    \proofstepwidestar[1]{R_y\colon A\to B}{%
      \FormulaRefAuto{TypedTermGraphFunction}{1}}
    \proofstepwidestar[3]{x\in A}{\rA}
    \proofstepwidestar[1,3]{(x,R_y(x))\in R_y}{%
      \FormulaRefAuto{F\colon A\to B\dsep x\in A\vdash (x,F(x))\in F}{2,3}}
    \proofstepwidestar[1,3]{R_y(x)=y}{%
      \rAEb{\FormulaRefAuto{x\in\{u\in A\mid P(u)\}
        \eqvdash x\in A\land P(x)}{4}}}
    \proofstepwidestarbreakkeep[1]{\forall x\in A\,R_y(x)=y}{%
      \rUI{\rRI{3,5}}}
  \closeproofpartwideindR

  \proofpartwideindR[Existenz]{%
    \begin{aligned}[t]
      &\forall x\in A\,y\in B\vdash{}\\[-1mm]
      &\quad\exists F\,\Bigl(
        F\colon A\to B\land
        \forall x\in A\,F(x)=y
      \Bigr)
    \end{aligned}
  }{%
    \forall x\in A\,y\in B\vdash
    \exists F\,\Bigl(
      F\colon A\to B\land
      \forall x\in A\,F(x)=y
    \Bigr)
  }[TypedTermFunctionExistence]
    \proofstepwidestar[1]{\forall x\in A\,y\in B}{\rA}
    \proofstepwidestar[1]{R_y\colon A\to B}{%
      \FormulaRefAuto{TypedTermGraphFunction}{1}}
    \proofstepwidestar[1]{\forall x\in A\,R_y(x)=y}{%
      \FormulaRefAuto{TypedTermGraphValues}{1}}
    \proofstepwidestarbreakkeep[1]{%
      \exists F\,\Bigl(
        F\colon A\to B\land
        \forall x\in A\,F(x)=y
      \Bigr)
    }{\rEI{\rAI{2,3}}}
  \closeproofpartwideindR

  \proofpartwideindR[Eindeutigkeit]{%
    \begin{aligned}[t]
      &F\colon A\to B\land\forall x\in A\,F(x)=y\dsep{}\\[-1mm]
      &G\colon A\to B\land\forall x\in A\,G(x)=y
      \vdash F=G
    \end{aligned}
  }{%
    F\colon A\to B\land\forall x\in A\,F(x)=y\dsep
    G\colon A\to B\land\forall x\in A\,G(x)=y
    \vdash F=G
  }[TypedTermFunctionUniqueness]
    \proofstepwidestarbreakkeep[1]{%
      F\colon A\to B\land\forall x\in A\,F(x)=y}{\rA}
    \proofstepwidestarbreakkeep[2]{%
      G\colon A\to B\land\forall x\in A\,G(x)=y}{\rA}
    \proofstepwidestar[1]{F\colon A\to B}{\rAEa{1}}
    \proofstepwidestar[1]{\forall x\in A\,F(x)=y}{\rAEb{1}}
    \proofstepwidestar[2]{G\colon A\to B}{\rAEa{2}}
    \proofstepwidestar[2]{\forall x\in A\,G(x)=y}{\rAEb{2}}
    \proofstepwidestar[1,2]{\forall x\in A\,F(x)=G(x)}{%
      \FormulaRefAuto{%
        \forall x\,F(x)\rightarrow P(x)=Q(x)\dsep
        \forall x\,F(x)\rightarrow P(x)=R(x)
        \vdash \forall x\,F(x)\rightarrow Q(x)=R(x)%
      }{4,6}}
    \proofstepwidestar[1,2]{F=G}{%
      \FormulaRefAuto{FunctionExtensionality}[thm]{3,5,7}}
  \closeproofpartwideindR

  \proofpartwide{Schluss}
    \proofstepwidestar[1]{\forall x\in A\,y\in B}{\rA}
    \proofstepwidestarbreakkeep[1]{%
      \exists F\,\Bigl(
        F\colon A\to B\land
        \forall x\in A\,F(x)=y
      \Bigr)
    }{\FormulaRefAuto{TypedTermFunctionExistence}{1}}
    \proofstepwidestarbreakkeep[3]{%
      F\colon A\to B\land\forall x\in A\,F(x)=y}{\rA}
    \proofstepwidestarbreakkeep[4]{%
      G\colon A\to B\land\forall x\in A\,G(x)=y}{\rA}
    \proofstepwidestar[3,4]{F=G}{%
      \FormulaRefAuto{TypedTermFunctionUniqueness}{3,4}}
    \proofstepwidestarbreakkeep[1]{%
      \exists!F\,\Bigl(
        F\colon A\to B\land
        \forall x\in A\,F(x)=y
      \Bigr)
    }{\UEI{2,3,4,5}}
  \closeproofpartwide
\end{tabproofsplitwide}

Aus einer punktweisen Typisierungsregel \(x\in A\vdash y\in B\) folgt die
Voraussetzung des Satzes durch Implikations- und Allquantoreinführung. In
einem konkreten Definitionsblock kann daher das frische Funktionssymbol
unmittelbar mit seinem Typ und der Grundgleichung eingeführt werden. Ein
vorheriges temporäres Vorschriftssymbol \(t\) ist dafür nicht erforderlich.
Rekursive oder nur partiell bestimmte Vorschriften benötigen weiterhin einen
eigenen Existenz- und Eindeutigkeitsbeweis.




\chapter{Beispiele und Konstruktionen allgemeiner Funktionen}

F\"ur die direkten Funktionsdefinitionen dieses und der folgenden B\"ande
gilt die folgende Konvention: Der Beweis des Definitionsblocks weist nur
noch nach, dass der rechts stehende Term f\"ur jedes Argument im angegebenen
Zielbereich liegt. Theorem
\FormulaRefAuto{TypedFunctionDefinitionPrinciple}[thm] liefert daraus die
Existenz und Eindeutigkeit der angegebenen Funktion samt ihrer
Grundgleichung. Ein tempor\"ares Vorschriftssymbol und ein gesonderter
Existenzsatz werden deshalb nicht wiederholt.

\section{Elementare Beispiele}

\subsection{Konstante Funktionen}

Konstante Funktionen sind die einfachsten Abbildungen mit vorgegebenem
Wertebereich: Ist \(b\in B\), so soll jedes Argument aus \(A\) auf denselben
Wert \(b\) abgebildet werden. Der einzige inhaltliche Nachweis ist die
bereits vorausgesetzte Zugehörigkeit \(b\in B\).

\FormulaDefDeltaK[Konstante Funktion]{%
  b\in B\vdash
  \begin{aligned}[t]
    &\ConstFun{A}{B}{b}\colon A\to B,\\[-2pt]
    &\forall x\in A\,
      \bigl(\ConstFun{A}{B}{b}(x)\coloneqq b\bigr)
  \end{aligned}%
}{Konstante Funktion}{%
  \DeltaRow{Mengen}{A\dsep B\dsep b\dsep x}%
  \DeltaRow{Funktionen}{\ConstFun{A}{B}{b}}%
}
\begin{tabproof}
  \proofstep{1}{b\in B}{\rA}
  \proofstep{2}{x\in A}{\rA}
  \proofstep{1,2}{b\in B}{\rAEa{\rAI{1,2}}}
  \proofstep{1}{\forall x\in A\,b\in B}{\rUI{\rRI{2,3}}}
\end{tabproof}

\FormulaThmDelta[Konstante Funktion ist eine Abbildung]{%
  b\in B\vdash \ConstFun{A}{B}{b}\colon A\to B%
}{%
  \DeltaRow{Mengen}{A\dsep B\dsep b}%
  \DeltaRow{Funktionensymbol}{\ConstFun{A}{B}{b}}%
}
\begin{tabproof}
  \proofstep{1}{b\in B}{\rA}
  \proofstep{1}{\ConstFun{A}{B}{b}\colon A\to B}{%
    \FormulaRefAuto{Konstante Funktion}[def]{1}}
\end{tabproof}

\FormulaThmDelta[Wert der konstanten Funktion]{%
  b\in B\dsep x\in A\vdash \ConstFun{A}{B}{b}(x)=b%
}{%
  \DeltaRow{Mengen}{A\dsep B\dsep b\dsep x}%
  \DeltaRow{Funktionen}{\ConstFun{A}{B}{b}}%
}
\begin{tabproof}
  \proofstep{1}{b\in B}{\rA}
  \proofstep{2}{x\in A}{\rA}
  \proofstep{1}{\forall u\in A\,\ConstFun{A}{B}{b}(u)=b}{%
    \FormulaRefAuto{Konstante Funktion}[def]{1}}
  \proofstep{1,2}{\ConstFun{A}{B}{b}(x)=b}{%
    \rRE{\rUE{3},2}}
\end{tabproof}

\subsection{Funktionen mit einelementigem Graphen}

\FormulaThmDeltaK[Ein geordnetes Paar als Funktionsgraph]{%
  b\in B\vdash\{(a,b)\}\colon\{a\}\to B%
}{SingletonGraphFunction}{%
  \DeltaRow{Mengen}{a\dsep b\dsep B\dsep x\dsep y\dsep z}
}
\begin{tabproofwide}
  \proofstepwidestar[1]{b\in B}{\rA}
  \proofstepwidestar[]{a\in\{a\}}{\FormulaRefAuto{a\in\{a\}}[thm]}
  \proofstepwidestar[1]{(a,b)\in\{a\}\times B}{%
    \FormulaRefAuto{a\in A\dsep b\in B\vdash(a,b)\in A\times B}[thm]{2,1}}
  \proofstepwidestar[1]{\{(a,b)\}\subseteq\{a\}\times B}{%
    \FormulaRefAuto{a\in A\vdash\{a\}\subseteq A}[thm]{3}}
  \proofstepwidestar[5]{x\in\{a\}}{\rA}
  \proofstepwidestar[5]{x=a}{%
    \FormulaRefAuto{x\in\{a\}\eqvdash x=a}[thm]{5}}
  \proofstepwidestar[]{(a,b)\in\{(a,b)\}}{%
    \FormulaRefAuto{a\in\{a\}}[thm]}
  \proofstepwidestar[5]{(x,b)\in\{(a,b)\}}{\rIE{6,7}}
  \proofstepwidestar[1,5]{\exists y\in B\,((x,y)\in\{(a,b)\})}{%
    \rEI{\rAI{1,8}}}
  \proofstepwidestar[10]{y\in B\land(x,y)\in\{(a,b)\}}{\rA}
  \proofstepwidestar[11]{z\in B\land(x,z)\in\{(a,b)\}}{\rA}
  \proofstepwidestar[10]{(x,y)=(a,b)}{%
    \FormulaRefAuto{x\in\{a\}\eqvdash x=a}[thm]{\rAEb{10}}}
  \proofstepwidestar[10]{y=b}{%
    \rAEb{\FormulaRefAuto{(a,b)=(c,d)\eqvdash a=c\land b=d}[thm]{12}}}
  \proofstepwidestar[11]{(x,z)=(a,b)}{%
    \FormulaRefAuto{x\in\{a\}\eqvdash x=a}[thm]{\rAEb{11}}}
  \proofstepwidestar[11]{z=b}{%
    \rAEb{\FormulaRefAuto{(a,b)=(c,d)\eqvdash a=c\land b=d}[thm]{14}}}
  \proofstepwidestar[10,11]{y=z}{%
    \FormulaRefAuto{a=b\dsep c=b\vdash a=c}[thm]{13,15}}
  \proofstepwidestar[1,5]{\exists!y\in B\,((x,y)\in\{(a,b)\})}{%
    \UEI{9,10,11,16}}
  \proofstepwidestar[1]{\forall x\in\{a\}\,\exists!y\in B\,
    ((x,y)\in\{(a,b)\})}{\rUI{\rRI{5,17}}}
  \proofstepwidestar[1]{\{(a,b)\}\colon\{a\}\to B}{%
    \FormulaRefAuto{UniqueValuedGraphFunction}[thm]{4,18}}
\end{tabproofwide}

\FormulaThmDeltaK[Auswertung eines einelementigen Funktionsgraphen]{%
  b\in B\vdash\{(a,b)\}(a)=b%
}{SingletonGraphValue}{%
  \DeltaRow{Mengen}{a\dsep b\dsep B}
}
\begin{tabproof}
  \proofstep{1}{b\in B}{\rA}
  \proofstep{1}{\{(a,b)\}\colon\{a\}\to B}{%
    \FormulaRefAuto{SingletonGraphFunction}[thm]{1}}
  \proofstep{}{a\in\{a\}}{\FormulaRefAuto{a\in\{a\}}[thm]}
  \proofstep{1}{(a,\{(a,b)\}(a))\in\{(a,b)\}}{%
    \FormulaRefAuto{F\colon A\to B\dsep x\in A\vdash(x,F(x))\in F}[thm]{2,3}}
  \proofstep{1}{(a,\{(a,b)\}(a))=(a,b)}{%
    \FormulaRefAuto{x\in\{a\}\eqvdash x=a}[thm]{4}}
  \proofstep{1}{\{(a,b)\}(a)=b}{%
    \rAEb{\FormulaRefAuto{(a,b)=(c,d)\eqvdash a=c\land b=d}[thm]{5}}}
\end{tabproof}

\subsection{Fallunterscheidungsabbildungen}

Eine Fallabbildung wird aus zwei bereits typisierten Zweigen gebildet. Der
Fallterm aus B02 ist auf \(A\cup B\) stets ein Element von \(C\); damit kann
die Funktion unmittelbar eingeführt werden.

\FormulaDefDeltaK[Fallabbildung]{%
  f\colon A\to C\dsep g\colon B\to C\vdash
  \begin{aligned}[t]
    &\CaseFun{A}{B}{f}{g}\colon A\cup B\to C,\\[-2pt]
    &\forall x\in A\cup B\,
      \bigl(
        \CaseFun{A}{B}{f}{g}(x)\coloneqq
        \IfThenElse{x\in A}{f(x)}{g(x)}
      \bigr)
  \end{aligned}
}{CaseFunctionDef}{%
  \DeltaRow{Mengen}{A\dsep B\dsep C\dsep x}
  \DeltaRow{Funktionen}{f\dsep g\dsep \CaseFun{A}{B}{f}{g}}
}
\begin{tabproof}
  \proofstep{1}{f\colon A\to C}{\rA}
  \proofstep{2}{g\colon B\to C}{\rA}
  \proofstep{3}{x\in A\cup B}{\rA}
  \proofstep{3}{x\in A\lor x\in B}{%
    \FormulaRefAuto{z\in A\cup B\eqvdash z\in A\lor z\in B}{3}}
  \proofstep{}{x\in A\lor x\notin A}{\FormulaRefAuto{P\lor\neg P}}

  \proofstep{6}{x\in A}{\rA}
  \proofstep{1,6}{f(x)\in C}{%
    \FormulaRefAuto{F\colon A\to B\dsep x\in A\vdash F(x)\in B}{1,6}}
  \proofstep{6}{\IfThenElse{x\in A}{f(x)}{g(x)}=f(x)}{%
    \FormulaRefAuto{P\vdash\IfThenElse{P}{a}{b}=a}{6}}
  \proofstep{1,6}{\IfThenElse{x\in A}{f(x)}{g(x)}\in C}{\rIE{8,7}}

  \proofstep{10}{x\notin A}{\rA}
  \proofstep{3,10}{x\in B}{%
    \FormulaRefAuto{P\lor Q\dsep\neg P\vdash Q}{4,10}}
  \proofstep{2,3,10}{g(x)\in C}{%
    \FormulaRefAuto{F\colon A\to B\dsep x\in A\vdash F(x)\in B}{2,11}}
  \proofstep{10}{\IfThenElse{x\in A}{f(x)}{g(x)}=g(x)}{%
    \FormulaRefAuto{\neg P\vdash\IfThenElse{P}{a}{b}=b}{10}}
  \proofstep{2,3,10}{\IfThenElse{x\in A}{f(x)}{g(x)}\in C}{\rIE{13,12}}

  \proofstep{1,2,3}{\IfThenElse{x\in A}{f(x)}{g(x)}\in C}{%
    \rOE{5,6,9,10,14}}
\end{tabproof}

\FormulaThmDeltaR[Fallabbildung ist eine Funktion]{%
  \begin{aligned}[t]
    &f\colon A\to C\dsep g\colon B\to C\vdash{}\\[-2pt]
    &\CaseFun{A}{B}{f}{g}\colon A\cup B\to C
  \end{aligned}
}{%
  f\colon A\to C\dsep g\colon B\to C
  \vdash \CaseFun{A}{B}{f}{g}\colon A\cup B\to C
}{%
  \DeltaRow{Mengen}{A\dsep B\dsep C}
  \DeltaRow{Funktionen}{f\dsep g}
  \DeltaRow{Funktionensymbol}{\CaseFun{A}{B}{f}{g}}
}
\begin{tabproof}
  \proofstep{1}{f\colon A\to C}{\rA}
  \proofstep{2}{g\colon B\to C}{\rA}
  \proofstep{1,2}{\CaseFun{A}{B}{f}{g}\colon A\cup B\to C}{%
    \FormulaRefAuto{CaseFunctionDef}[def]{1,2}}
\end{tabproof}

\FormulaThmDeltaR[Grundgleichung der Fallabbildung]{%
  \begin{aligned}[t]
    &f\colon A\to C\dsep g\colon B\to C\dsep
      x\in A\cup B\vdash{}\\[-2pt]
    &\CaseFun{A}{B}{f}{g}(x)=\IfThenElse{x\in A}{f(x)}{g(x)}
  \end{aligned}
}{%
  f\colon A\to C\dsep g\colon B\to C\dsep x\in A\cup B
  \vdash
  \CaseFun{A}{B}{f}{g}(x)=\IfThenElse{x\in A}{f(x)}{g(x)}
}{%
  \DeltaRow{Mengen}{A\dsep B\dsep C\dsep x}
  \DeltaRow{Funktionen}{f\dsep g}
  \DeltaRow{Funktionensymbol}{\CaseFun{A}{B}{f}{g}}
}
\begin{tabproof}
  \proofstep{1}{f\colon A\to C}{\rA}
  \proofstep{2}{g\colon B\to C}{\rA}
  \proofstep{3}{x\in A\cup B}{\rA}
  \proofstep{1,2}{\forall u\in A\cup B\,
    \CaseFun{A}{B}{f}{g}(u)=\IfThenElse{u\in A}{f(u)}{g(u)}}{%
    \FormulaRefAuto{CaseFunctionDef}[def]{1,2}}
  \proofstep{1,2,3}{\CaseFun{A}{B}{f}{g}(x)=
    \IfThenElse{x\in A}{f(x)}{g(x)}}{%
    \rRE{\rUE{4},3}}
\end{tabproof}

\FormulaThmDeltaR[Fallabbildung im linken Fall]{%
  \begin{aligned}[t]
    &f\colon A\to C\dsep g\colon B\to C\dsep x\in A\vdash{}\\[-2pt]
    &\CaseFun{A}{B}{f}{g}(x)=f(x)
  \end{aligned}
}{%
  f\colon A\to C\dsep g\colon B\to C\dsep x\in A
  \vdash \CaseFun{A}{B}{f}{g}(x)=f(x)
}{%
  \DeltaRow{Mengen}{A\dsep B\dsep C\dsep x}
  \DeltaRow{Funktionen}{f\dsep g}
  \DeltaRow{Funktionensymbol}{\CaseFun{A}{B}{f}{g}}
}
\begin{tabproof}
  \proofstep{1}{f\colon A\to C}{\rA}
  \proofstep{2}{g\colon B\to C}{\rA}
  \proofstep{3}{x\in A}{\rA}
  \proofstep{3}{x\in A\cup B}{%
    \FormulaRefAuto{z \in A \vdash z \in A \cup B}{3}}
  \proofstep{1,2,3}{\CaseFun{A}{B}{f}{g}(x)=
    \IfThenElse{x\in A}{f(x)}{g(x)}}{%
    \FormulaRefAuto{f\colon A\to C\dsep g\colon B\to C\dsep x\in A\cup B
    \vdash
    \CaseFun{A}{B}{f}{g}(x)=
    \IfThenElse{x\in A}{f(x)}{g(x)}}{1,2,4}}
  \proofstep{3}{\IfThenElse{x\in A}{f(x)}{g(x)}=f(x)}{%
    \FormulaRefAuto{P \vdash \IfThenElse{P}{a}{b}=a}{3}}
  \proofstep{1,2,3}{\CaseFun{A}{B}{f}{g}(x)=f(x)}{%
    \FormulaRefAuto{a=b,\,b=c\vdash a=c}{5,6}}
\end{tabproof}

\FormulaThmDeltaR[Fallabbildung im rechten Fall]{%
  \begin{aligned}[t]
    &A\cap B=\varnothing\dsep f\colon A\to C\dsep
      g\colon B\to C\dsep x\in B\vdash{}\\[-2pt]
    &\CaseFun{A}{B}{f}{g}(x)=g(x)
  \end{aligned}
}{%
  A\cap B=\varnothing\dsep f\colon A\to C\dsep g\colon B\to C\dsep x\in B
  \vdash \CaseFun{A}{B}{f}{g}(x)=g(x)
}{%
  \DeltaRow{Mengen}{A\dsep B\dsep C\dsep x}
  \DeltaRow{Funktionen}{f\dsep g}
  \DeltaRow{Funktionensymbol}{\CaseFun{A}{B}{f}{g}}
}
\begin{tabproof}
  \proofstep{1}{A\cap B=\varnothing}{\rA}
  \proofstep{2}{f\colon A\to C}{\rA}
  \proofstep{3}{g\colon B\to C}{\rA}
  \proofstep{4}{x\in B}{\rA}
  \proofstep{4}{x\in A\cup B}{%
    \FormulaRefAuto{z \in B \vdash z \in A \cup B}{4}}
  \proofstep{1,4}{x\notin A}{%
    \FormulaRefAuto{A \cap B = \varnothing,\ x \in B \vdash x \notin A}{1,4}}
  \proofstep{2,3,4}{\CaseFun{A}{B}{f}{g}(x)=
    \IfThenElse{x\in A}{f(x)}{g(x)}}{%
    \FormulaRefAuto{f\colon A\to C\dsep g\colon B\to C\dsep x\in A\cup B
    \vdash
    \CaseFun{A}{B}{f}{g}(x)=
    \IfThenElse{x\in A}{f(x)}{g(x)}}{2,3,5}}
  \proofstep{1,4}{\IfThenElse{x\in A}{f(x)}{g(x)}=g(x)}{%
    \FormulaRefAuto{\neg P \vdash \IfThenElse{P}{a}{b}=b}{6}}
  \proofstep{1,2,3,4}{\CaseFun{A}{B}{f}{g}(x)=g(x)}{%
    \FormulaRefAuto{a=b,\,b=c\vdash a=c}{7,8}}
\end{tabproof}


\section{Fasern und Einschränkungen}

\subsection{Die Faserabbildung}

\FormulaDefDeltaK[Faser zu $F$]{%
  F\colon A\to B\vdash
  \begin{aligned}[t]
    &\Fib_F\colon B\to\powerset(A),\\[-2pt]
    &\forall y\in B\,
      \bigl(\Fib_F(y)\coloneqq F^{-1}[\{y\}]\bigr)
  \end{aligned}%
}{FiberMapDef}{%
  \DeltaRow{Mengen}{A\dsep B\dsep y}
  \DeltaRow{Funktionen}{F\dsep \Fib_F}
}
\begin{tabproof}
  \proofstep{1}{F\colon A\to B}{\rA}
  \proofstep{2}{y\in B}{\rA}
  \proofstep{2}{\{y\}\subseteq B}{%
    \FormulaRefAuto{x\in A\vdash\{x\}\subseteq A}{2}}
  \proofstep{1,2}{F^{-1}[\{y\}]
    =\{u\in A\mid F(u)\in\{y\}\}}{%
    \FormulaRefAuto{F\colon A\to B\dsep C\subseteq B\vdash
      F^{-1}[C]\coloneqq\{x\in A\mid F(x)\in C\}}{1,3}}
  \proofstep{}{\{u\in A\mid F(u)\in\{y\}\}\subseteq A}{%
    \FormulaRefAuto{\{x\in A\mid P(x)\}\subseteq A}}
  \proofstep{1,2}{F^{-1}[\{y\}]\subseteq A}{\rIE{4,5}}
  \proofstep{1,2}{F^{-1}[\{y\}]\in\powerset(A)}{%
    \FormulaRefAuto{x\in\powerset(A)\eqvdash x\subseteq A}{6}}
\end{tabproof}

\FormulaThmDelta[Faser ist eine Funktion]{%
  \Fib_F\colon B\to \powerset(A)%
}{
  \DeltaRow{Mengen}{A\dsep B}
  \DeltaPrem{Funktionen}{F\colon A\to B}
}
\begin{tabproof}
  \proofstep{}{\Fib_F\colon B\to \powerset(A)}{%
    \FormulaRefAuto{FiberMapDef}[def]}
\end{tabproof}

\FormulaThmDelta[Grundgleichung der Faser]{%
  y\in B \vdash \Fib_F(y)=F^{-1}[\{y\}]%
}{
  \DeltaRow{Mengen}{A\dsep B\dsep y}
  \DeltaPrem{Funktionen}{F\colon A\to B}
  \DeltaRow{Funktionensymbol}{\Fib_F}
}
\begin{tabproof}
  \proofstep{1}{y\in B}{\rA}
  \proofstep{}{\forall u\in B\,\Fib_F(u)=F^{-1}[\{u\}]}{%
    \FormulaRefAuto{FiberMapDef}[def]}
  \proofstep{1}{\Fib_F(y)=F^{-1}[\{y\}]}{%
    \rRE{\rUE{2},1}}
\end{tabproof}

\subsubsection{Die Faser als Menge}

\FormulaThmDelta[Faser als Urbild eines Singletons]{%
  y\in B \vdash \Fib_F(y) = \{\,x\in A \mid F(x)=y\,\}%
}{
  \DeltaRow{Mengen}{A\dsep B\dsep x\dsep y}
  \DeltaPrem{Funktionen}{F\colon A\to B}
}
\begin{tabproofwide}
  \proofstepwidestar[1]{y\in B}{\rA}
  \proofstepwidestar[1]{\{y\}\subseteq B}{%
    \FormulaRefAuto{x\in A\vdash \{x\}\subseteq A}{1}}

  \proofstepwide[1]{\Fib_F(y)}{=}{F^{-1}[\{y\}]}{%
    \FormulaRefAuto{y\in B \vdash \Fib_F(y)=F^{-1}[\{y\}]}{1}}

  \proofstepwide[1]{}{=}{\{\,u\in A \mid F(u)\in \{y\}\,\}}{%
    \FormulaRefAuto{F\colon A\to B\dsep C\subseteq B
    \vdash F^{-1}[C] \coloneqq \{\,x\in A \mid F(x)\in C\,\}}{%
    \FormulaRefAuto{Funktion},2}}

  \proofstepwide[1]{}{=}{\{\,u\in A \mid F(u)=y\,\}}{%
    \FormulaRefAuto{\{x\in A\mid F(x)\in \{a\}\}=\{x\in A\mid F(x)=a\}}}

  \proofstepwide[1]{\Fib_F(y)}{=}{\{\,u\in A \mid F(u)=y\,\}}{%
    \rChain{3,4,5}}
\end{tabproofwide}


\FormulaThmDelta{y\in B\dsep x\in \Fib_F(y) \vdash x\in A\land F(x)=y}
{
  \DeltaRow{Mengen}{A\dsep B\dsep x\dsep y}
  \DeltaPrem{Funktionen}{F\colon A\to B}
}
\begin{tabproof}
  \proofstep{1}{y\in B}{\rA}
  \proofstep{2}{x\in \Fib_F(y)}{\rA}
  \proofstep{1}{\Fib_F(y)=\{\,t\in A \mid F(t)=y\,\}}{%
    \FormulaRefAuto{y\in B \vdash \Fib_F(y) = \{\,x\in A \mid F(x)=y\,\}}{1}}
  \proofstep{1,2}{x\in \{\,t\in A \mid F(t)=y\,\}}{\rIE{3,2}}
  \proofstep{1,2}{x\in A \land F(x)=y}{%
    \FormulaRefAuto{x\in\{u\in A\mid P(u)\}\eqvdash x\in A\land P(x)}{4}}
\end{tabproof}

\FormulaThmDelta{y\in B\dsep x\in \Fib_F(y) \vdash x\in A}
{
  \DeltaRow{Mengen}{A\dsep B\dsep x\dsep y}
  \DeltaPrem{Funktionen}{F\colon A\to B}
}
\begin{tabproof}
  \proofstep{1}{y\in B}{\rA}
  \proofstep{2}{x\in \Fib_F(y)}{\rA}
  \proofstep{1,2}{x\in A\land F(x)=y}{\FormulaRefAuto{y\in B\dsep x\in \Fib_F(y) \vdash x\in A\land F(x)=y}{1,2}}
  \proofstep{1,2}{x\in A}{\rAEa{3}}
\end{tabproof}

\FormulaThmDelta{y\in B\dsep x\in \Fib_F(y) \vdash F(x)=y}
{
  \DeltaRow{Mengen}{A\dsep B\dsep x\dsep y}
  \DeltaPrem{Funktionen}{F\colon A\to B}
}
\begin{tabproof}
  \proofstep{1}{y\in B}{\rA}
  \proofstep{2}{x\in \Fib_F(y)}{\rA}
  \proofstep{1,2}{x\in A\land F(x)=y}{\FormulaRefAuto{y\in B\dsep x\in \Fib_F(y) \vdash x\in A\land F(x)=y}{1,2}}
  \proofstep{1,2}{F(x)=y}{\rAEb{3}}
\end{tabproof}

\FormulaThmDelta
{y\in B\dsep x\in A\dsep F(x)=y\vdash x\in \Fib_F(y)}
{
  \DeltaRow{Mengen}{A\dsep B\dsep x\dsep y}
  \DeltaPrem{Funktionen}{F\colon A\to B}
}
\begin{tabproof}
  \proofstep{1}{y\in B}{\rA}
  \proofstep{2}{x\in A}{\rA}
  \proofstep{3}{F(x)=y}{\rA}
  \proofstep{2,3}{x\in A\land F(x)=y}{\rAI{2,3}}
  \proofstep{2,3}{x\in \{\,t\in A \mid F(t)=y\,\}}{%
    \FormulaRefAuto{x\in\{u\in A\mid P(u)\}\eqvdash x\in A\land P(x)}{4}}
  \proofstep{1}{\Fib_F(y)=\{\,t\in A \mid F(t)=y\,\}}{%
    \FormulaRefAuto{y\in B \vdash \Fib_F(y) = \{\,x\in A \mid F(x)=y\,\}}{1}}
  \proofstep{1,2,3}{x\in \Fib_F(y)}{\rIE{6,5}}
\end{tabproof}


\FormulaThmDelta
{y\in B\vdash x\in\Fib_F(y)\leftrightarrow x\in A\land F(x)=y}
{
  \DeltaRow{Mengen}{A\dsep B\dsep x\dsep y}
  \DeltaPrem{Funktionen}{F\colon A\to B}
}
\begin{tabproof}
  \proofstep{1}{y\in B}{\rA}

  \proofstep{2}{x\in \Fib_F(y)}{\rA}
  \proofstep{1,2}{x\in A\land F(x)=y}{%
    \FormulaRefAuto{y\in B\dsep x\in \Fib_F(y) \vdash x\in A\land F(x)=y}{1,2}}
  \proofstep{1}{x\in \Fib_F(y)\rightarrow (x\in A\land F(x)=y)}{\rRI{2,3}}

  \proofstep{5}{x\in A\land F(x)=y}{\rA}
  \proofstep{5}{x\in A}{\rAEa{5}}
  \proofstep{5}{F(x)=y}{\rAEb{5}}
  \proofstep{1,5}{x\in \Fib_F(y)}{%
    \FormulaRefAuto{y\in B\dsep x\in A\dsep F(x)=y\vdash x\in \Fib_F(y)}{1,6,7}}
  \proofstep{1}{(x\in A\land F(x)=y)\rightarrow x\in \Fib_F(y)}{\rRI{5,8}}

  \proofstep{1}{x\in\Fib_F(y)\leftrightarrow (x\in A\land F(x)=y)}{\rLRI{4,9}}
\end{tabproof}

\FormulaThmDelta[Verschiedene Fasern sind disjunkt]{%
  y\in B \dsep z\in B \dsep y\neq z \vdash \Fib_F(y)\cap \Fib_F(z)=\varnothing%
}{
  \DeltaRow{Mengen}{A\dsep B\dsep y\dsep z\dsep a}
  \DeltaPrem{Funktionen}{F\colon A\to B}
}
\begin{tabproof}
  \proofstep{1}{y\in B}{\rA}
  \proofstep{2}{z\in B}{\rA}
  \proofstep{3}{y\neq z}{\rA}

  \proofstep{4}{a\in \Fib_F(y)\cap \Fib_F(z)}{\rA}
  \proofstep{4}{a\in \Fib_F(y)}{%
    \FormulaRefAuto{x \in A \cap B \vdash x \in A}{4}}
  \proofstep{4}{a\in \Fib_F(z)}{%
    \FormulaRefAuto{x \in A \cap B \vdash x \in B}{4}}

  \proofstep{1,4}{F(a)=y}{%
    \FormulaRefAuto{y\in B\dsep x\in \Fib_F(y) \vdash F(x)=y}{1,5}}
  \proofstep{2,4}{F(a)=z}{%
    \FormulaRefAuto{y\in B\dsep x\in \Fib_F(y) \vdash F(x)=y}{2,6}}

  \proofstep{1,2,4}{y=z}{%
    \FormulaRefAuto{a=b \dsep a=c \vdash b=c}{7,8}}
  \proofstep{1,2,3,4}{\bot}{\rBI{3,9}}

  \proofstep{1,2,3}{a\notin \Fib_F(y)\cap \Fib_F(z)}{\rCE{4,10}}
  \proofstep{1,2,3}{\forall a\,\bigl(a\notin \Fib_F(y)\cap \Fib_F(z)\bigr)}{\rUI{11}}
  \proofstep{1,2,3}{\Fib_F(y)\cap \Fib_F(z)=\varnothing}{%
    \FormulaRefAuto{\forall b\,b\notin X\vdash X=\varnothing}{12}}
\end{tabproof}

\FormulaThmDelta{%
  y\in B \dsep z\in B \dsep X=\Fib_F(y) \dsep Y=\Fib_F(z) \dsep X\neq Y
  \vdash X\cap Y=\varnothing%
}{
  \DeltaRow{Mengen}{A\dsep B\dsep X\dsep Y\dsep y\dsep z}
  \DeltaPrem{Funktionen}{F\colon A\to B}
}
\begin{tabproof}
  \proofstep{1}{y\in B}{\rA}
  \proofstep{2}{z\in B}{\rA}
  \proofstep{3}{X=\Fib_F(y)}{\rA}
  \proofstep{4}{Y=\Fib_F(z)}{\rA}
  \proofstep{5}{X\neq Y}{\rA}

  \proofstep{6}{y=z}{\rA}
  \proofstep{4,6}{Y=\Fib_F(y)}{%
    \FormulaRefAuto{x = y \dsep a = F(y) \vdash a = F(x)}{6,4}}

  \proofstep{3,4,6}{X=Y}{%
    \FormulaRefAuto{a = b,\, c = b \vdash a = c}{3,7}}

  \proofstep{3,4,5,6}{\bot}{\rBI{5,8}}
  \proofstep{1,2,3,4,5}{y\neq z}{\rCI{6,9}}

  \proofstep{1,2,3,4,5}{\Fib_F(y)\cap \Fib_F(z)=\varnothing}{%
    \FormulaRefAuto{y\in B \dsep z\in B \dsep y\neq z \vdash \Fib_F(y)\cap \Fib_F(z)=\varnothing}{1,2,10}}

  \proofstep{1,2,3,4,5}{X\cap Y=\varnothing}{%
    \FormulaRefAuto{X=U \dsep Y=V \dsep U\cap V=W \vdash X\cap Y=W}{3,4,11}}
\end{tabproof}



\FormulaThmDelta{%
  X\in \Fib_F[B] \dsep Y\in \Fib_F[B] \dsep X\neq Y \vdash X\cap Y=\varnothing%
}{
  \DeltaRow{Mengen}{A\dsep B\dsep X\dsep Y\dsep y\dsep z}
  \DeltaPrem{Funktionen}{F\colon A\to B}
}
\begin{tabproof}
  \proofstep{1}{X\in \Fib_F[B]}{\rA}
  \proofstep{2}{Y\in \Fib_F[B]}{\rA}
  \proofstep{3}{X\neq Y}{\rA}
  \proofstep{}{\Fib_F\colon B\to \powerset(A)}{%
    \FormulaRefAuto{\Fib_F\colon B\to \powerset(A)}}

  \proofstep{1}{\exists y\in B\,X=\Fib_F(y)}{%
    \FormulaRefAuto{F\colon A\to B\dsep y\in F[A] \vdash \exists x\in A\,y=F(x)}{4,1}}

  \proofstep{2}{\exists z\in B\,Y=\Fib_F(z)}{%
    \FormulaRefAuto{F\colon A\to B\dsep y\in F[A] \vdash \exists x\in A\,y=F(x)}{4,2}}

  \proofstep{7}{y\in B \land X=\Fib_F(y)}{\rA}
  \proofstep{7}{y\in B}{\rAEa{7}}
  \proofstep{7}{X=\Fib_F(y)}{\rAEb{7}}

  \proofstep{10}{z\in B \land Y=\Fib_F(z)}{\rA}
  \proofstep{10}{z\in B}{\rAEa{10}}
  \proofstep{10}{Y=\Fib_F(z)}{\rAEb{10}}

  \proofstep{3,7,10}{X\cap Y=\varnothing}{%
    \FormulaRefAuto{y\in B \dsep z\in B \dsep X=\Fib_F(y) \dsep Y=\Fib_F(z) \dsep X\neq Y \vdash X\cap Y=\varnothing}{8,11,9,12,3}}

  \proofstep{2,3,7}{X\cap Y=\varnothing}{\rEE{6,10,13}}
  \proofstep{1,2,3}{X\cap Y=\varnothing}{\rEE{5,7,14}}
\end{tabproof}


\subsection{Die Einschränkung einer Funktion}

\FormulaDefDeltaK[Einschränkung von \(F\) auf \(C\)]{%
  C\subseteq A\vdash
  \begin{aligned}[t]
    &F\restriction_C\colon C\to B,\\[-2pt]
    &\forall x\in C\,
      \bigl((F\restriction_C)(x)\coloneqq F(x)\bigr)
  \end{aligned}%
}{RestrictionDef}{%
  \DeltaRow{Mengen}{A\dsep B\dsep C\dsep x}
  \DeltaPrem{Funktionen}{F\colon A\to B}
  \DeltaRow{Funktionensymbol}{F\restriction_C}
}
\begin{tabproof}
  \proofstep{1}{C\subseteq A}{\rA}
  \proofstep{2}{x\in C}{\rA}
  \proofstep{1,2}{x\in A}{%
    \FormulaRefAuto{A\subseteq B\dsep x\in A\vdash x\in B}{1,2}}
  \proofstep{1,2}{F(x)\in B}{%
    \FormulaRefAuto{F\colon A\to B\dsep x\in A\vdash F(x)\in B}{3}}
\end{tabproof}

\FormulaThmDelta[Einschränkung von \(F\) auf \(C\)]
{C\subseteq A \vdash F\restriction_C\colon C\to B}
{
  \DeltaRow{Mengen}{A\dsep B\dsep C}
  \DeltaPrem{Funktionen}{F\colon A\to B}
}
\begin{tabproof}
  \proofstep{1}{C\subseteq A}{\rA}
  \proofstep{1}{F\restriction_C\colon C\to B}{%
    \FormulaRefAuto{RestrictionDef}[def]{1}}
\end{tabproof}

\FormulaThmDelta
{C\subseteq A \vdash \TotRel{F\restriction_C,C,B}}
{
  \DeltaRow{Mengen}{A\dsep B}
}
\begin{tabproof}
    \proofstep{1}{C\subseteq A }{\rA}
    \proofstep{1}{F\restriction_C\colon C\to B}{\FormulaRefAuto{C\subseteq A \vdash F\restriction_C\colon C\to B}{1}}
    \proofstep{1}{\TotRel{F\restriction_C,C,B}}{\FormulaRefAuto{F\colon A\to B\vdash\TotRel{F,A,B}}{2}}
\end{tabproof}

\FormulaThmDelta[Grundgleichung der Einschränkung]
{C\subseteq A \dsep x\in C \vdash F\restriction_C(x)=F(x)}
{
  \DeltaRow{Mengen}{A\dsep B\dsep C\dsep x}
  \DeltaPrem{Funktionen}{F\colon A\to B}
}
\begin{tabproof}
  \proofstep{1}{C\subseteq A}{\rA}
  \proofstep{2}{x\in C}{\rA}
  \proofstep{1}{\forall u\in C\,(F\restriction_C)(u)=F(u)}{%
    \FormulaRefAuto{RestrictionDef}[def]{1}}
  \proofstep{1,2}{F\restriction_C(x)=F(x)}{%
    \rRE{\rUE{3},2}}
\end{tabproof}

\FormulaThmDeltaK[Gleichheitskriterium für Funktionseinschränkungen]{%
  F\colon A\to B\dsep G\colon C\to B\dsep C\subseteq A\vdash
  \Bigl(
    F\restriction_C=G
    \leftrightarrow
    \forall x\in C\,(F(x)=G(x))
  \Bigr)%
}{FunctionRestrictionEqualityCriterion}{%
  \DeltaRow{Mengen}{A\dsep B\dsep C\dsep x}
  \DeltaRow{Funktionen}{F\dsep G}
}
\begin{tabproofwide}
  \proofstepwidestar[1]{F\colon A\to B}{\rA}
  \proofstepwidestar[2]{G\colon C\to B}{\rA}
  \proofstepwidestar[3]{C\subseteq A}{\rA}
  \proofstepwidestar[1,3]{F\restriction_C\colon C\to B}{%
    \FormulaRefAuto{RestrictionDef}[def]{3}}

  \proofstepwidestar[5]{F\restriction_C=G}{\rA}
  \proofstepwidestar[6]{x\in C}{\rA}
  \proofstepwidestar[1,3,6]{(F\restriction_C)(x)=F(x)}{%
    \FormulaRefAuto{C\subseteq A\dsep x\in C
      \vdash F\restriction_C(x)=F(x)}[thm]{3,6}}
  \proofstepwidestar[1,2,3,5,6]{(F\restriction_C)(x)=G(x)}{%
    \FormulaRefAuto{FunctionEqualityEvaluation}[thm]{4,2,6,5}}
  \proofstepwidestar[1,3,6]{F(x)=(F\restriction_C)(x)}{%
    \FormulaRefAuto{a=b\vdash b=a}[thm]{7}}
  \proofstepwidestar[1,2,3,5,6]{F(x)=G(x)}{%
    \FormulaRefAuto{a=b\dsep b=c\vdash a=c}[thm]{9,8}}
  \proofstepwidestar[1,2,3,5]{%
    \forall x\in C\,(F(x)=G(x))}{%
    \rUI{\rRI{6,10}}}
  \proofstepwidestar[1,2,3]{%
    F\restriction_C=G\rightarrow
    \forall x\in C\,(F(x)=G(x))}{%
    \rRI{5,11}}

  \proofstepwidestar[13]{\forall x\in C\,(F(x)=G(x))}{\rA}
  \proofstepwidestar[14]{x\in C}{\rA}
  \proofstepwidestar[1,3,14]{(F\restriction_C)(x)=F(x)}{%
    \FormulaRefAuto{C\subseteq A\dsep x\in C
      \vdash F\restriction_C(x)=F(x)}[thm]{3,14}}
  \proofstepwidestar[13,14]{F(x)=G(x)}{%
    \rRE{\rUE{13},14}}
  \proofstepwidestar[1,3,13,14]{(F\restriction_C)(x)=G(x)}{%
    \FormulaRefAuto{a=b\dsep b=c\vdash a=c}[thm]{15,16}}
  \proofstepwidestar[1,3,13]{%
    \forall x\in C\,((F\restriction_C)(x)=G(x))}{%
    \rUI{\rRI{14,17}}}
  \proofstepwidestar[1,2,3,13]{F\restriction_C=G}{%
    \FormulaRefAuto{FunctionExtensionality}[thm]{4,2,18}}
  \proofstepwidestar[1,2,3]{%
    \forall x\in C\,(F(x)=G(x))\rightarrow F\restriction_C=G}{%
    \rRI{13,19}}
  \proofstepwidestar[1,2,3]{%
    F\restriction_C=G
    \leftrightarrow
    \forall x\in C\,(F(x)=G(x))}{%
    \rLRI{12,20}}
\end{tabproofwide}

\subsubsection{Bild der Einschränkung}

\FormulaThmDelta[Bild der Einschränkung]{%
  C\subseteq A \dsep F\colon A\to B \vdash (F\restriction_C)[C]=F[C]
}{%
  \DeltaRow{Mengen}{A\dsep B\dsep C\dsep x\dsep y}
  \DeltaRow{Funktionen}{F}
}
\begin{tabproofsplitwide}
  \proofpartwideind[Teilmengenrichtung \((F\restriction_C)[C]\subseteq F[C]\)]{%
    C\subseteq A \dsep F\colon A\to B \vdash (F\restriction_C)[C]\subseteq F[C]
  }[%
    C\subseteq A \dsep F\colon A\to B \vdash (F\restriction_C)[C]\subseteq F[C]
  ]
    \proofstepwidestar[1]{C\subseteq A}{\rA}
    \proofstepwidestar[2]{F\colon A\to B}{\rA}
    \proofstepwidestar[1,2]{F\restriction_C\colon C\to B}{%
      \FormulaRefAuto{C\subseteq A \vdash F\restriction_C\colon C\to B}{1}}
    \proofstepwidestar[4]{y\in (F\restriction_C)[C]}{\rA}
    \proofstepwidestar[1,2,4]{\exists x\in C\,y=F\restriction_C(x)}{%
      \FormulaRefAuto{F\colon A\to B\dsep y\in F[A] \vdash \exists x\in A\,y=F(x)}{3,4}}
    \proofstepwidestar[6]{x\in C\land y=F\restriction_C(x)}{\rA}
    \proofstepwidestar[6]{x\in C}{\rAEa{6}}
    \proofstepwidestar[6]{y=F\restriction_C(x)}{\rAEb{6}}
    \proofstepwidestar[1,6]{F\restriction_C(x)=F(x)}{%
      \FormulaRefAuto{C\subseteq A \dsep x\in C \vdash F\restriction_C(x)=F(x)}{1,7}}
    \proofstepwidestar[1,6]{y=F(x)}{\FormulaRefAuto{a=b,\, b=c \vdash a=c}{8,9}}
    \proofstepwidestar[1,6]{\exists x\in C\,y=F(x)}{\rEI{\rAI{7,10}}}
    \proofstepwidestar[1,2,6]{y\in F[C]}{%
      \FormulaRefAuto{C\subseteq A\dsep F\colon A\to B\dsep \exists x\in C\,y=F(x)\vdash y\in F[C]}{1,2,11}}
    \proofstepwidestar[1,2]{\forall y\,(y\in (F\restriction_C)[C]\rightarrow y\in F[C])}{\rUI{\rRI{4,12}}}
    \proofstepwidestar[1,2]{(F\restriction_C)[C]\subseteq F[C]}{%
      \FormulaRefAuto{A\subseteq B \coloneq \forall x\,(x\in A \rightarrow x\in B)}{13}}
  \closeproofpartwideind

  \proofpartwideind[Teilmengenrichtung \(F[C]\subseteq (F\restriction_C)[C]\)]{%
    C\subseteq A \dsep F\colon A\to B \vdash F[C]\subseteq (F\restriction_C)[C]
  }[%
    C\subseteq A \dsep F\colon A\to B \vdash F[C]\subseteq (F\restriction_C)[C]
  ]
    \proofstepwidestar[1]{C\subseteq A}{\rA}
    \proofstepwidestar[2]{F\colon A\to B}{\rA}
    \proofstepwidestar[1,2]{F\restriction_C\colon C\to B}{%
      \FormulaRefAuto{C\subseteq A \vdash F\restriction_C\colon C\to B}{1}}
    \proofstepwidestar[4]{y\in F[C]}{\rA}
    \proofstepwidestar[1,2,4]{\exists x\in C\,y=F(x)}{%
      \FormulaRefAuto{C\subseteq A\dsep F\colon A\to B\dsep y\in F[C]\vdash \exists x\in C\,y=F(x)}{1,2,4}}
    \proofstepwidestar[6]{x\in C\land y=F(x)}{\rA}
    \proofstepwidestar[6]{x\in C}{\rAEa{6}}
    \proofstepwidestar[6]{y=F(x)}{\rAEb{6}}
    \proofstepwidestar[1,6]{F\restriction_C(x)=F(x)}{%
      \FormulaRefAuto{C\subseteq A \dsep x\in C \vdash F\restriction_C(x)=F(x)}{1,7}}
    \proofstepwidestar[1,6]{F(x)=F\restriction_C(x)}{\FormulaRefAuto{a=b\vdash b=a}{9}}
    \proofstepwidestar[1,6]{y=F\restriction_C(x)}{\FormulaRefAuto{a=b,\, b=c \vdash a=c}{8,10}}
    \proofstepwidestar[1,6]{\exists x\in C\,y=F\restriction_C(x)}{\rEI{\rAI{7,11}}}
    \proofstepwidestar[1,2,6]{y\in (F\restriction_C)[C]}{%
      \FormulaRefAuto{F\colon A\to B\dsep \exists x\in A\,y=F(x)\vdash y\in F[A]}{3,12}}
    \proofstepwidestar[1,2]{\forall y\,(y\in F[C]\rightarrow y\in (F\restriction_C)[C])}{\rUI{\rRI{4,13}}}
    \proofstepwidestar[1,2]{F[C]\subseteq (F\restriction_C)[C]}{%
      \FormulaRefAuto{A\subseteq B \coloneq \forall x\,(x\in A \rightarrow x\in B)}{14}}
  \closeproofpartwideind

  \proofpartwide{Zusammenfassung}
    \proofstepwidestar[1]{C\subseteq A}{\rA}
    \proofstepwidestar[2]{F\colon A\to B}{\rA}
    \proofstepwidestar[1,2]{(F\restriction_C)[C]\subseteq F[C]}{%
      \FormulaRefAuto{C\subseteq A \dsep F\colon A\to B \vdash (F\restriction_C)[C]\subseteq F[C]}{1,2}}
    \proofstepwidestar[1,2]{F[C]\subseteq (F\restriction_C)[C]}{%
      \FormulaRefAuto{C\subseteq A \dsep F\colon A\to B \vdash F[C]\subseteq (F\restriction_C)[C]}{1,2}}
    \proofstepwidestar[1,2]{(F\restriction_C)[C]=F[C]}{%
      \FormulaRefAuto{A\subseteq B\dsep B\subseteq A \vdash A=B}{3,4}}
  \closeproofpartwide
\end{tabproofsplitwide}


\subsubsection{Die Auswertungseigenschaft}

\FormulaThmDelta{%
C\subseteq A \dsep x\in C \dsep y\in C \dsep F(x)=F(y)\vdash F\restriction_C(x)=F\restriction_C(y)
}{
  \DeltaRow{Mengen}{A \dsep B \dsep C \dsep x \dsep y}
  \DeltaPrem{Funktionen}{F\colon A\to B}
}
\begin{tabproof}
  \proofstep{1}{C\subseteq A}{\rA}
  \proofstep{2}{x\in C}{\rA}
  \proofstep{3}{y\in C}{\rA}
  \proofstep{4}{F(x)=F(y)}{\rA}
  \proofstep{1,2}{F\restriction_C(x)=F(x)}{%
    \FormulaRefAuto{C\subseteq A \dsep x\in C \vdash F\restriction_C(x)=F(x)}{1,2}}
  \proofstep{1,3}{F\restriction_C(y)=F(y)}{%
    \FormulaRefAuto{C\subseteq A \dsep x\in C \vdash F\restriction_C(x)=F(x)}{1,3}}
  \proofstep{1,3}{F(y)=F\restriction_C(y)}{%
    \FormulaRefAuto{a=b\vdash b=a}{6}}
  \proofstep{1,3,4}{F(x)=F\restriction_C(y)}{%
    \FormulaRefAuto{a = b,\, b = c \vdash a = c}{4,7}}
  \proofstep{1,2,3,4}{F\restriction_C(x)=F\restriction_C(y)}{%
    \FormulaRefAuto{a = b,\, b = c \vdash a = c}{5,8}}
\end{tabproof}

\FormulaThmDelta{%
C\subseteq A \dsep x\in C \dsep y\in C \dsep F\restriction_C(x)=F\restriction_C(y)\vdash F(x)=F(y)
}{
  \DeltaRow{Mengen}{A \dsep B \dsep C \dsep x \dsep y}
  \DeltaPrem{Funktionen}{F\colon A\to B}
}
\begin{tabproof}
  \proofstep{1}{C\subseteq A}{\rA}
  \proofstep{2}{x\in C}{\rA}
  \proofstep{3}{y\in C}{\rA}
  \proofstep{4}{F\restriction_C(x)=F\restriction_C(y)}{\rA}
  \proofstep{1,2}{F\restriction_C(x)=F(x)}{%
    \FormulaRefAuto{C\subseteq A \dsep x\in C \vdash F\restriction_C(x)=F(x)}{1,2}}
  \proofstep{1,3}{F\restriction_C(y)=F(y)}{%
    \FormulaRefAuto{C\subseteq A \dsep x\in C \vdash F\restriction_C(x)=F(x)}{1,3}}
  \proofstep{1,2}{F(x)=F\restriction_C(x)}{%
    \FormulaRefAuto{a=b\vdash b=a}{5}}
  \proofstep{1,2,4}{F(x)=F\restriction_C(y)}{%
    \FormulaRefAuto{a = b,\, b = c \vdash a = c}{7,4}}
  \proofstep{1,2,3,4}{F(x)=F(y)}{%
    \FormulaRefAuto{a = b,\, b = c \vdash a = c}{8,6}}
\end{tabproof}


\subsubsection{Verkleinerung der Zielmenge}

\FormulaThmDelta[Verkleinerung der Zielmenge]{%
  F\colon A\to B\dsep \forall x\in A\,F(x)\in C
  \vdash
  F\colon A\to C
}{%
  \DeltaRow{Mengen}{A\dsep B\dsep C\dsep x\dsep y\dsep z}
  \DeltaRow{Funktionen}{F}
}
\begin{tabproof}
  \proofstep{1}{F\colon A\to B}{\rA}
  \proofstep{2}{\forall x\in A\,F(x)\in C}{\rA}

  \proofstep{3}{(x,y)\in F}{\rA}
  \proofstep{1,3}{x\in A}{%
    \FormulaRefAuto{F\colon A\to B\dsep (x,y)\in F\vdash x\in A}{1,3}}
  \proofstep{1,2,3}{F(x)\in C}{\rRE{2,4}}
  \proofstep{1,3}{y=F(x)}{%
    \FormulaRefAuto{F\colon A\to B\dsep (x,y)\in F\vdash y=F(x)}{1,3}}
  \proofstep{1,2,3}{y\in C}{\rIE{6,5}}
  \proofstep{1,2,3}{(x,y)\in A\times C}{%
    \FormulaRefAuto{x\in A\dsep y\in C\vdash (x,y)\in A\times C}{4,7}}
  \proofstep{1,2}{(x,y)\in F\rightarrow (x,y)\in A\times C}{\rRI{3,8}}
  \proofstep{1,2}{F\subseteq A\times C}{%
    \FormulaRefAuto{A \subseteq B \coloneq \forall x\,(x\in A \rightarrow x\in B)}{\rUI{9}}}

  \proofstep{1}{(x,y)\in F\land (x,z)\in F\rightarrow y=z}{%
    \FormulaRefAuto{(x,y)\in F\dsep (x,z)\in F \vdash y=z}{1}}
  \proofstep{1}{x\in A \rightarrow \exists y\,(x,y)\in F}{%
    \FormulaRefAuto{x\in A \vdash \exists y\,(x,y)\in F}{1}}

  \proofstep{1,2}{F\colon A\to C}{\FormulaRefAuto{Funktion}{10,11,12}}
\end{tabproof}

\FormulaThmDelta[Erweiterung der Zielmenge]{%
  F\colon A\to B\dsep B\subseteq C\vdash F\colon A\to C
}{%
  \DeltaRow{Mengen}{A\dsep B\dsep C\dsep x}
  \DeltaRow{Funktionen}{F}
}
\begin{tabproof}
  \proofstep{1}{F\colon A\to B}{\rA}
  \proofstep{2}{B\subseteq C}{\rA}
  \proofstep{3}{x\in A}{\rA}
  \proofstep{1,3}{F(x)\in B}{%
    \FormulaRefAuto{F\colon A\to B\dsep x\in A\vdash F(x)\in B}{1,3}}
  \proofstep{2,3}{F(x)\in C}{%
    \FormulaRefAuto{A\subseteq B,\, x\in A \vdash x\in B}{2,4}}
  \proofstep{1,2}{\forall x\in A\,F(x)\in C}{\rUI{\rRI{3,5}}}
  \proofstep{1,2}{F\colon A\to C}{%
    \FormulaRefAuto{F\colon A\to B\dsep \forall x\in A\,F(x)\in C\vdash F\colon A\to C}{1,6}}
\end{tabproof}

\subsubsection{Einschränkung auf das Urbild einer Teilmenge}

\FormulaThmDelta[Einschränkung auf ein Urbild als Funktion]{%
  F\colon A\to B\dsep T\subseteq B \vdash F\restriction_{F^{-1}[T]}\colon F^{-1}[T]\to T
}{
  \DeltaRow{Mengen}{A\dsep B\dsep T\dsep x}
  \DeltaRow{Funktionen}{F}
}
\begin{tabproof}
  \proofstep{1}{F\colon A\to B}{\rA}
  \proofstep{2}{T\subseteq B}{\rA}
  \proofstep{1,2}{F^{-1}[T]\subseteq A}{%
    \FormulaRefAuto{F\colon A\to B\dsep C\subseteq B \vdash F^{-1}[C]\subseteq A}{1,2}}
  \proofstep{1,2}{F\restriction_{F^{-1}[T]}\colon F^{-1}[T]\to B}{%
    \FormulaRefAuto{C\subseteq A \vdash F\restriction_C\colon C\to B}{3}}
  \proofstep{5}{x\in F^{-1}[T]}{\rA}
  \proofstep{1,2,5}{x\in F^{-1}[T]\leftrightarrow (x\in A \land F(x)\in T)}{%
    \FormulaRefAuto{F\colon A\to B\dsep C\subseteq B
    \vdash \bigl(x\in F^{-1}[C] \leftrightarrow
    (x\in A \land F(x)\in C)\bigr)}{1,2}}
  \proofstep{1,2,5}{x\in A \land F(x)\in T}{%
    \FormulaRefAuto{P \leftrightarrow Q, P \vdash Q}{6,5}}
  \proofstep{1,2,5}{F(x)\in T}{\rAEb{7}}
  \proofstep{1,2,5}{F\restriction_{F^{-1}[T]}(x)=F(x)}{%
    \FormulaRefAuto{C\subseteq A \dsep x\in C \vdash F\restriction_C(x)=F(x)}{3,5}}
  \proofstep{1,2,5}{F\restriction_{F^{-1}[T]}(x)\in T}{\rIE{9,8}}
  \proofstep{1,2}{\forall x\in F^{-1}[T]\,F\restriction_{F^{-1}[T]}(x)\in T}{%
    \rUI{\rRI{5,10}}}
  \proofstep{1,2}{F\restriction_{F^{-1}[T]}\colon F^{-1}[T]\to T}{%
    \FormulaRefAuto{F\colon A\to B\dsep \forall x\in A\,F(x)\in C \vdash F\colon A\to C}{4,11}}
\end{tabproof}


\subsection{Vereinigung kompatibler Funktionsgraphen}

\FormulaThmDeltaK[Vereinigung einer kompatiblen Familie von
Funktionsgraphen]{%
  \begin{aligned}[t]
    &H\colon I\to\powerset(U\times B)\dsep
      D\colon I\to\powerset(U)\dsep{}\\[-2pt]
    &\forall i\in I\,\bigl(H(i)\colon D(i)\to B\bigr)\dsep{}\\[-2pt]
    &\forall i\in I\,\forall j\in I\,
      \forall x\in D(i)\cap D(j)\,
        H(i)(x)=H(j)(x)\\[-2pt]
    &\vdash
      \bigcup H[I]\colon\bigcup D[I]\to B\dsep{}\\[-2pt]
    &\phantom{\vdash{}}
      \forall i\in I\,\forall x\in D(i)\,
        \bigl(\bigcup H[I]\bigr)(x)=H(i)(x)
  \end{aligned}%
}{CompatibleFunctionGraphFamilyUnion}{%
  \DeltaRow{Mengen}{I\dsep U\dsep B\dsep
    G\dsep E\dsep i\dsep j\dsep x\dsep y\dsep z\dsep p}
  \DeltaRow{Funktionen}{D\dsep H}
}
\begingroup
\let\proofstepwidestar\proofstepwidestarbreak
\begin{tabproofsplitwide}
  \proofpartwideindR[Lokalisierung eines Paars]{%
    \begin{aligned}[t]
      &H\colon I\to\powerset(U\times B)\dsep
        D\colon I\to\powerset(U)\dsep{}\\[-2pt]
      &\forall i\in I\,\bigl(H(i)\colon D(i)\to B\bigr)\dsep
        (x,y)\in\bigcup H[I]\\[-2pt]
      &\vdash\exists i\in I\,\Bigl(
        (x,y)\in H(i)\land\bigl(x\in D(i)\land y\in B\bigr)\Bigr)
    \end{aligned}
  }{%
    \begin{aligned}[t]
      &H\colon I\to\powerset(U\times B)\dsep
        D\colon I\to\powerset(U)\dsep{}\\[-2pt]
      &\forall i\in I\,\bigl(H(i)\colon D(i)\to B\bigr)\dsep
        (x,y)\in\bigcup H[I]\\[-2pt]
      &\vdash\exists i\in I\,\Bigl(
        (x,y)\in H(i)\land\bigl(x\in D(i)\land y\in B\bigr)\Bigr)
    \end{aligned}
  }[CompatibleFunctionGraphFamilyPairLocalization]
    \proofstepwidestar[1]{H\colon I\to\powerset(U\times B)}{\rA}
    \proofstepwidestar[2]{D\colon I\to\powerset(U)}{\rA}
    \proofstepwidestar[3]{%
      \forall i\in I\,\bigl(H(i)\colon D(i)\to B\bigr)}{\rA}
    \proofstepwidestar[4]{(x,y)\in\bigcup H[I]}{\rA}
    \proofstepwidestar[4]{%
      \exists G\in H[I],\bigl((x,y)\in G\bigr)}{%
      \FormulaRefAuto{x \in \bigcup A \eqvdash
        \exists B\in A\,(x\in B)}[thm]{4}}
    \proofstepwidestar[6]{G\in H[I]\land(x,y)\in G}{\rA}
    \proofstepwidestar[6]{G\in H[I]}{\rAEa{6}}
    \proofstepwidestar[6]{(x,y)\in G}{\rAEb{6}}
    \proofstepwidestar[1,6]{\exists i\in I\,G=H(i)}{%
      \FormulaRefAuto{F\colon A\to B\dsep y\in F[A]
        \vdash\exists x\in A\,y=F(x)}[thm]{1,7}}
    \proofstepwidestar[10]{i\in I\land G=H(i)}{\rA}
    \proofstepwidestar[10]{i\in I}{\rAEa{10}}
    \proofstepwidestar[10]{G=H(i)}{\rAEb{10}}
    \proofstepwidestar[2,3,10]{H(i)\colon D(i)\to B}{%
      \rRE{\rUE{3},11}}
    \proofstepwidestar[6,10]{(x,y)\in H(i)}{\rIE{12,8}}
    \proofstepwidestar[2,3,6,10]{x\in D(i)}{%
      \FormulaRefAuto{F\colon A\to B\dsep (x,y)\in F
        \vdash x\in A}[thm]{13,14}}
    \proofstepwidestar[2,3,6,10]{y\in B}{%
      \FormulaRefAuto{F\colon A\to B\dsep (x,y)\in F
        \vdash y\in B}[thm]{13,14}}
    \proofstepwidestar[2,3,6,10]{%
      (x,y)\in H(i)\land\bigl(x\in D(i)\land y\in B\bigr)}{%
      \rAI{14,\rAI{15,16}}}
    \proofstepwidestar[2,3,6,10]{%
      \exists i\in I\,\Bigl(
        (x,y)\in H(i)\land\bigl(x\in D(i)\land y\in B\bigr)\Bigr)}{%
      \rEI{\rAI{11,17}}}
    \proofstepwidestar[1,2,3,6]{%
      \exists i\in I\,\Bigl(
        (x,y)\in H(i)\land\bigl(x\in D(i)\land y\in B\bigr)\Bigr)}{%
      \rEE{9,10,18}}
    \proofstepwidestar[1,2,3,4]{%
      \exists i\in I\,\Bigl(
        (x,y)\in H(i)\land\bigl(x\in D(i)\land y\in B\bigr)\Bigr)}{%
      \rEE{5,6,19}}
  \closeproofpartwideindR

  \proofpartwideindR[Graphteilmenge]{%
    \begin{aligned}[t]
      &H\colon I\to\powerset(U\times B)\dsep
        D\colon I\to\powerset(U)\dsep{}\\[-2pt]
      &\forall i\in I\,\bigl(H(i)\colon D(i)\to B\bigr)
        \vdash\bigcup H[I]\subseteq\bigl(\bigcup D[I]\bigr)\times B
    \end{aligned}
  }{%
    \begin{aligned}[t]
      &H\colon I\to\powerset(U\times B)\dsep
        D\colon I\to\powerset(U)\dsep{}\\[-2pt]
      &\forall i\in I\,\bigl(H(i)\colon D(i)\to B\bigr)
        \vdash\bigcup H[I]\subseteq\bigl(\bigcup D[I]\bigr)\times B
    \end{aligned}
  }[CompatibleFunctionGraphFamilyUnionSubset]
    \proofstepwidestar[1]{H\colon I\to\powerset(U\times B)}{\rA}
    \proofstepwidestar[2]{D\colon I\to\powerset(U)}{\rA}
    \proofstepwidestar[3]{%
      \forall i\in I\,\bigl(H(i)\colon D(i)\to B\bigr)}{\rA}
    \proofstepwidestar[4]{p\in\bigcup H[I]}{\rA}
    \proofstepwidestar[4]{\exists G\in H[I]\,(p\in G)}{%
      \FormulaRefAuto{x \in \bigcup A \eqvdash
        \exists B\in A\,(x\in B)}[thm]{4}}
    \proofstepwidestar[6]{G\in H[I]\land p\in G}{\rA}
    \proofstepwidestar[6]{G\in H[I]}{\rAEa{6}}
    \proofstepwidestar[6]{p\in G}{\rAEb{6}}
    \proofstepwidestar[1,6]{\exists i\in I\,G=H(i)}{%
      \FormulaRefAuto{F\colon A\to B\dsep y\in F[A]
        \vdash\exists x\in A\,y=F(x)}[thm]{1,7}}
    \proofstepwidestar[10]{i\in I\land G=H(i)}{\rA}
    \proofstepwidestar[10]{i\in I}{\rAEa{10}}
    \proofstepwidestar[10]{G=H(i)}{\rAEb{10}}
    \proofstepwidestar[2,3,10]{H(i)\colon D(i)\to B}{%
      \rRE{\rUE{3},11}}
    \proofstepwidestar[6,10]{p\in H(i)}{\rIE{12,8}}
    \proofstepwidestar[2,3,10]{H(i)\subseteq D(i)\times B}{%
      \FormulaRefAuto{F \subseteq A \times B}{13}}
    \proofstepwidestar[2,3,6,10]{p\in D(i)\times B}{%
      \FormulaRefAuto{A\subseteq B\dsep x\in A\vdash x\in B}[thm]{15,14}}
    \proofstepwidestar[2,3,6,10]{%
      \exists x\in D(i)\,\exists y\in B\,p=(x,y)}{%
      \FormulaRefAuto{x\in A\times B\eqvdash
        \exists a\in A\,\exists b\in B\,x=(a,b)}[thm]{16}}
    \proofstepwidestar[18]{%
      x\in D(i)\land y\in B\land p=(x,y)}{\rA}
    \proofstepwidestar[18]{x\in D(i)}{\rAEa{18}}
    \proofstepwidestar[18]{y\in B}{\rAEa{\rAEb{18}}}
    \proofstepwidestar[18]{p=(x,y)}{\rAEb{\rAEb{18}}}
    \proofstepwidestar[2,10]{D(i)\in D[I]}{%
      \FormulaRefAuto{C\subseteq A\dsep F\colon A\to B\dsep
        x\in C\vdash F(x)\in F[C]}[thm]{%
          \FormulaRefAuto{A\subseteq A}[thm],2,11}}
    \proofstepwidestar[2,10,18]{x\in\bigcup D[I]}{%
      \FormulaRefAuto{B\in A\dsep x\in B\vdash x\in\bigcup A}[thm]{22,19}}
    \proofstepwidestar[2,10,18]{(x,y)\in(\bigcup D[I])\times B}{%
      \FormulaRefAuto{a\in A\dsep b\in B\vdash
        (a,b)\in A\times B}[thm]{23,20}}
    \proofstepwidestar[2,10,18]{p\in(\bigcup D[I])\times B}{%
      \rIE{21,24}}
    \proofstepwidestar[2,3,6,10]{p\in(\bigcup D[I])\times B}{%
      \rEE{17,18,25}}
    \proofstepwidestar[1,2,3,6]{p\in(\bigcup D[I])\times B}{%
      \rEE{9,10,26}}
    \proofstepwidestar[1,2,3,4]{p\in(\bigcup D[I])\times B}{%
      \rEE{5,6,27}}
    \proofstepwidestar[1,2,3]{%
      \forall p\,\bigl(p\in\bigcup H[I]\rightarrow
        p\in(\bigcup D[I])\times B\bigr)}{%
      \rUI{\rRI{4,28}}}
    \proofstepwidestar[1,2,3]{%
      \bigcup H[I]\subseteq(\bigcup D[I])\times B}{%
      \FormulaRefAuto{A\subseteq B \coloneq
        \forall x\,(x\in A\rightarrow x\in B)}[def]{29}}
  \closeproofpartwideindR

  \proofpartwideindR[Funktionale Eindeutigkeit]{%
    \begin{aligned}[t]
      &H\colon I\to\powerset(U\times B)\dsep
        D\colon I\to\powerset(U)\dsep{}\\[-2pt]
      &\forall i\in I\,\bigl(H(i)\colon D(i)\to B\bigr)\dsep{}\\[-2pt]
      &\forall i\in I\,\forall j\in I\,
        \forall u\in D(i)\cap D(j)\,H(i)(u)=H(j)(u)\dsep{}\\[-2pt]
      &(x,y)\in\bigcup H[I]\dsep(x,z)\in\bigcup H[I]\vdash y=z
    \end{aligned}
  }{%
    \begin{aligned}[t]
      &H\colon I\to\powerset(U\times B)\dsep
        D\colon I\to\powerset(U)\dsep{}\\[-2pt]
      &\forall i\in I\,\bigl(H(i)\colon D(i)\to B\bigr)\dsep{}\\[-2pt]
      &\forall i\in I\,\forall j\in I\,
        \forall u\in D(i)\cap D(j)\,H(i)(u)=H(j)(u)\dsep{}\\[-2pt]
      &(x,y)\in\bigcup H[I]\dsep(x,z)\in\bigcup H[I]\vdash y=z
    \end{aligned}
  }[CompatibleFunctionGraphFamilyUnionFunctional]
    \proofstepwidestar[1]{H\colon I\to\powerset(U\times B)}{\rA}
    \proofstepwidestar[2]{D\colon I\to\powerset(U)}{\rA}
    \proofstepwidestar[3]{%
      \forall i\in I\,\bigl(H(i)\colon D(i)\to B\bigr)}{\rA}
    \proofstepwidestar[4]{%
      \forall i\in I\,\forall j\in I\,
      \forall u\in D(i)\cap D(j)\,H(i)(u)=H(j)(u)}{\rA}
    \proofstepwidestar[5]{(x,y)\in\bigcup H[I]}{\rA}
    \proofstepwidestar[6]{(x,z)\in\bigcup H[I]}{\rA}
    \proofstepwidestar[1,2,3,5]{%
      \exists i\in I\,\Bigl(
        (x,y)\in H(i)\land\bigl(x\in D(i)\land y\in B\bigr)\Bigr)}{%
      \FormulaRefAuto{CompatibleFunctionGraphFamilyPairLocalization}[thm]{1,2,3,5}}
    \proofstepwidestar[1,2,3,6]{%
      \exists j\in I\,\Bigl(
        (x,z)\in H(j)\land\bigl(x\in D(j)\land z\in B\bigr)\Bigr)}{%
      \FormulaRefAuto{CompatibleFunctionGraphFamilyPairLocalization}[thm]{1,2,3,6}}
    \proofstepwidestar[9]{%
      i\in I\land\Bigl((x,y)\in H(i)\land
        \bigl(x\in D(i)\land y\in B\bigr)\Bigr)}{\rA}
    \proofstepwidestar[10]{%
      j\in I\land\Bigl((x,z)\in H(j)\land
        \bigl(x\in D(j)\land z\in B\bigr)\Bigr)}{\rA}
    \proofstepwidestar[9]{i\in I}{\rAEa{9}}
    \proofstepwidestar[9]{(x,y)\in H(i)}{\rAEa{\rAEb{9}}}
    \proofstepwidestar[9]{x\in D(i)}{\rAEa{\rAEb{\rAEb{9}}}}
    \proofstepwidestar[10]{j\in I}{\rAEa{10}}
    \proofstepwidestar[10]{(x,z)\in H(j)}{\rAEa{\rAEb{10}}}
    \proofstepwidestar[10]{x\in D(j)}{\rAEa{\rAEb{\rAEb{10}}}}
    \proofstepwidestar[9,10]{x\in D(i)\cap D(j)}{%
      \FormulaRefAuto{x\in A\dsep x\in B\vdash x\in A\cap B}[thm]{13,16}}
    \proofstepwidestar[4,9,10]{H(i)(x)=H(j)(x)}{%
      \rRE{17,\rUE{\rRE{14,\rUE{\rRE{11,\rUE{4}}}}}}}
    \proofstepwidestar[2,3,9]{H(i)\colon D(i)\to B}{%
      \rRE{\rUE{3},11}}
    \proofstepwidestar[2,3,10]{H(j)\colon D(j)\to B}{%
      \rRE{\rUE{3},14}}
    \proofstepwidestar[2,3,9]{H(i)(x)=y}{%
      \FormulaRefAuto{F\colon A\to B\dsep (x,y)\in F
        \vdash F(x)=y}[thm]{19,12}}
    \proofstepwidestar[2,3,10]{H(j)(x)=z}{%
      \FormulaRefAuto{F\colon A\to B\dsep (x,y)\in F
        \vdash F(x)=y}[thm]{20,15}}
    \proofstepwidestar[2,3,9]{y=H(i)(x)}{%
      \FormulaRefAuto{a=b\vdash b=a}[thm]{21}}
    \proofstepwidestar[2,3,4,9,10]{y=H(j)(x)}{%
      \FormulaRefAuto{a=b\dsep b=c\vdash a=c}[thm]{23,18}}
    \proofstepwidestar[2,3,4,9,10]{y=z}{%
      \FormulaRefAuto{a=b\dsep b=c\vdash a=c}[thm]{24,22}}
    \proofstepwidestar[1,2,3,4,5,9]{y=z}{\rEE{8,10,25}}
    \proofstepwidestar[1,2,3,4,5,6]{y=z}{\rEE{7,9,26}}
  \closeproofpartwideindR

  \proofpartwideindR[Totalität]{%
    \begin{aligned}[t]
      &H\colon I\to\powerset(U\times B)\dsep
        D\colon I\to\powerset(U)\dsep{}\\[-2pt]
      &\forall i\in I\,\bigl(H(i)\colon D(i)\to B\bigr)\dsep
        x\in\bigcup D[I]\\[-2pt]
      &\vdash\exists y\,\bigl((x,y)\in\bigcup H[I]\bigr)
    \end{aligned}
  }{%
    \begin{aligned}[t]
      &H\colon I\to\powerset(U\times B)\dsep
        D\colon I\to\powerset(U)\dsep{}\\[-2pt]
      &\forall i\in I\,\bigl(H(i)\colon D(i)\to B\bigr)\dsep
        x\in\bigcup D[I]\\[-2pt]
      &\vdash\exists y\,\bigl((x,y)\in\bigcup H[I]\bigr)
    \end{aligned}
  }[CompatibleFunctionGraphFamilyUnionTotal]
    \proofstepwidestar[1]{H\colon I\to\powerset(U\times B)}{\rA}
    \proofstepwidestar[2]{D\colon I\to\powerset(U)}{\rA}
    \proofstepwidestar[3]{%
      \forall i\in I\,\bigl(H(i)\colon D(i)\to B\bigr)}{\rA}
    \proofstepwidestar[4]{x\in\bigcup D[I]}{\rA}
    \proofstepwidestar[4]{\exists E\in D[I]\,(x\in E)}{%
      \FormulaRefAuto{x \in \bigcup A \eqvdash
        \exists B\in A\,(x\in B)}[thm]{4}}
    \proofstepwidestar[6]{E\in D[I]\land x\in E}{\rA}
    \proofstepwidestar[6]{E\in D[I]}{\rAEa{6}}
    \proofstepwidestar[6]{x\in E}{\rAEb{6}}
    \proofstepwidestar[2,6]{\exists i\in I\,E=D(i)}{%
      \FormulaRefAuto{F\colon A\to B\dsep y\in F[A]
        \vdash\exists x\in A\,y=F(x)}[thm]{2,7}}
    \proofstepwidestar[10]{i\in I\land E=D(i)}{\rA}
    \proofstepwidestar[10]{i\in I}{\rAEa{10}}
    \proofstepwidestar[10]{E=D(i)}{\rAEb{10}}
    \proofstepwidestar[6,10]{x\in D(i)}{\rIE{12,8}}
    \proofstepwidestar[2,3,10]{H(i)\colon D(i)\to B}{%
      \rRE{\rUE{3},11}}
    \proofstepwidestar[2,3,6,10]{\exists y\,((x,y)\in H(i))}{%
      \FormulaRefAuto{x\in A\vdash\exists y\,(x,y)\in F}{14,13}}
    \proofstepwidestar[1,10]{H(i)\in H[I]}{%
      \FormulaRefAuto{C\subseteq A\dsep F\colon A\to B\dsep
        x\in C\vdash F(x)\in F[C]}[thm]{%
          \FormulaRefAuto{A\subseteq A}[thm],1,11}}
    \proofstepwidestar[1,10]{H(i)\subseteq\bigcup H[I]}{%
      \FormulaRefAuto{B\in A\vdash B\subseteq\bigcup A}[thm]{16}}
    \proofstepwidestar[18]{(x,y)\in H(i)}{\rA}
    \proofstepwidestar[1,10,18]{(x,y)\in\bigcup H[I]}{%
      \FormulaRefAuto{A\subseteq B\dsep x\in A\vdash x\in B}[thm]{17,18}}
    \proofstepwidestar[1,10,18]{%
      \exists y\,\bigl((x,y)\in\bigcup H[I]\bigr)}{\rEI{19}}
    \proofstepwidestar[1,2,3,6,10]{%
      \exists y\,\bigl((x,y)\in\bigcup H[I]\bigr)}{\rEE{15,18,20}}
    \proofstepwidestar[1,2,3,6]{%
      \exists y\,\bigl((x,y)\in\bigcup H[I]\bigr)}{\rEE{9,10,21}}
    \proofstepwidestar[1,2,3,4]{%
      \exists y\,\bigl((x,y)\in\bigcup H[I]\bigr)}{\rEE{5,6,22}}
  \closeproofpartwideindR

  \proofpartwide{Schluss}
    \proofstepwidestar[1]{H\colon I\to\powerset(U\times B)}{\rA}
    \proofstepwidestar[2]{D\colon I\to\powerset(U)}{\rA}
    \proofstepwidestar[3]{%
      \forall i\in I\,\bigl(H(i)\colon D(i)\to B\bigr)}{\rA}
    \proofstepwidestar[4]{%
      \forall i\in I\,\forall j\in I\,
      \forall x\in D(i)\cap D(j)\,H(i)(x)=H(j)(x)}{\rA}
    \proofstepwidestar[1,2,3,4]{%
      \bigcup H[I]\colon\bigcup D[I]\to B}{%
      \FormulaRefAuto{Funktion}[def]{%
        \FormulaRefAuto{CompatibleFunctionGraphFamilyUnionSubset}[thm]{1,2,3},
        \FormulaRefAuto{CompatibleFunctionGraphFamilyUnionFunctional}[thm]{1,2,3,4},
        \FormulaRefAuto{CompatibleFunctionGraphFamilyUnionTotal}[thm]{1,2,3}}}
    \proofstepwidestar[6]{i\in I}{\rA}
    \proofstepwidestar[7]{x\in D(i)}{\rA}
    \proofstepwidestar[2,3,6]{H(i)\colon D(i)\to B}{%
      \rRE{\rUE{3},6}}
    \proofstepwidestar[2,3,6,7]{(x,H(i)(x))\in H(i)}{%
      \FormulaRefAuto{F\colon A\to B\dsep x\in A
        \vdash(x,F(x))\in F}[thm]{8,7}}
    \proofstepwidestar[1,6]{H(i)\in H[I]}{%
      \FormulaRefAuto{C\subseteq A\dsep F\colon A\to B\dsep
        x\in C\vdash F(x)\in F[C]}[thm]{%
          \FormulaRefAuto{A\subseteq A}[thm],1,6}}
    \proofstepwidestar[1,6]{H(i)\subseteq\bigcup H[I]}{%
      \FormulaRefAuto{B\in A\vdash B\subseteq\bigcup A}[thm]{10}}
    \proofstepwidestar[1,2,3,6,7]{%
      (x,H(i)(x))\in\bigcup H[I]}{%
      \FormulaRefAuto{A\subseteq B\dsep x\in A\vdash x\in B}[thm]{11,9}}
    \proofstepwidestar[1,2,3,4,6,7]{%
      (\bigcup H[I])(x)=H(i)(x)}{%
      \FormulaRefAuto{F\colon A\to B\dsep (x,y)\in F
        \vdash F(x)=y}[thm]{5,12}}
    \proofstepwidestar[1,2,3,4,6]{%
      \forall x\in D(i)\,(\bigcup H[I])(x)=H(i)(x)}{%
      \rUI{\rRI{7,13}}}
    \proofstepwidestar[1,2,3,4]{%
      \forall i\in I\,\forall x\in D(i)\,
        (\bigcup H[I])(x)=H(i)(x)}{%
      \rUI{\rRI{6,14}}}
  \closeproofpartwide
\end{tabproofsplitwide}
\endgroup


\section{Komposition und induzierte Abbildungen}

\subsection{Die Komposition von Funktionen}

\FormulaDefDeltaK[Komposition $G\circ F$]{%
  F\colon A\to B\dsep G\colon B\to C\vdash
  \begin{aligned}[t]
    &G\circ F\colon A\to C,\\[-2pt]
    &\forall x\in A\,
      \bigl((G\circ F)(x)\coloneqq G(F(x))\bigr)
  \end{aligned}%
}{CompositionDef}{%
  \DeltaRow{Mengen}{A\dsep B\dsep C\dsep x}
  \DeltaRow{Funktionen}{F\dsep G\dsep G\circ F}
}
\begin{tabproof}
  \proofstep{1}{F\colon A\to B}{\rA}
  \proofstep{2}{G\colon B\to C}{\rA}
  \proofstep{3}{x\in A}{\rA}
  \proofstep{1,3}{F(x)\in B}{%
    \FormulaRefAuto{F\colon A\to B\dsep x\in A\vdash F(x)\in B}{1,3}}
  \proofstep{1,2,3}{G(F(x))\in C}{%
    \FormulaRefAuto{F\colon A\to B\dsep x\in A\vdash F(x)\in B}{2,4}}
\end{tabproof}

\FormulaThmDelta[Komposition als Abbildung]{%
  G\circ F\colon A\to C%
}{
  \DeltaRow{Mengen}{A\dsep B\dsep C}
  \DeltaPrem{Funktionen}{F\colon A\to B\dsep G\colon B\to C}
}
\begin{tabproof}
  \proofstep{}{G\circ F\colon A\to C}{%
    \FormulaRefAuto{CompositionDef}[def]}
\end{tabproof}

\FormulaThmDelta[Grundgleichung der Komposition (für $F\circ G$)]%
{x\in A \vdash (F\circ G)(x)=F(G(x))}%
{
  \DeltaRow{Mengen}{A\dsep B\dsep C\dsep x}
  \DeltaPrem{Funktionen}{G\colon A\to B\dsep F\colon B\to C}
}
\begin{tabproof}
  \proofstep{1}{x\in A}{\rA}
  \proofstep{}{\forall u\in A\,(F\circ G)(u)=F(G(u))}{%
    \FormulaRefAuto{CompositionDef}[def]}
  \proofstep{1}{(F\circ G)(x)=F(G(x))}{%
    \rRE{\rUE{2},1}}
\end{tabproof}


% ------------------------------------------------------------
% (2) Rückwärts-Variante (saubere Gleichungskette / Symmetrie)
% ------------------------------------------------------------
\FormulaThmDelta%
{x\in A \vdash F(G(x))=(F\circ G)(x)}%
{
  \DeltaRow{Mengen}{A\dsep B\dsep C\dsep x}
  \DeltaPrem{Funktionen}{G\colon A\to B\dsep F\colon B\to C}
}
\begin{tabproof}
  \proofstep{1}{x\in A}{\rA}
  \proofstep{1}{(F\circ G)(x)=F(G(x))}{%
    \FormulaRefAuto{x\in A \vdash (F\circ G)(x)=F(G(x))}{1}}
  \proofstep{1}{F(G(x))=(F\circ G)(x)}{%
    \FormulaRefAuto{a=b\vdash b=a}{2}}
\end{tabproof}

% ------------------------------------------------------------
% (3) Grundgleichung der Komposition für drei Komponenten
% ------------------------------------------------------------
\FormulaThmDelta[Grundgleichung der Komposition für drei Komponenten]%
{x\in A \vdash ((F\circ G)\circ H)(x)=F(G(H(x)))}%
{
  \DeltaRow{Mengen}{A\dsep B\dsep C\dsep D\dsep x}
  \DeltaPrem{Funktionen}{H\colon A\to B\dsep G\colon B\to C\dsep F\colon C\to D}
}
\begin{tabproofwide}
  \proofstepwidestar[1]{x\in A}{\rA}

  \proofstepwidestar[1]{H(x)\in B}{%
    \FormulaRefAuto{H\colon A\to B\dsep x\in A\vdash H(x)\in B}{1}}

  \proofstepwide[1]{((F\circ G)\circ H)(x)}{=}{(F\circ G)(H(x))}{%
    \FormulaRefAuto{x\in A \vdash (F\circ G)(x)=F(G(x))}{1}}

  \proofstepwide[1]{}{=}{F(G(H(x)))}{%
    \FormulaRefAuto{x\in A \vdash (F\circ G)(x)=F(G(x))}{2}}

  \proofstepwide[1]{((F\circ G)\circ H)(x)}{=}{F(G(H(x)))}{\rChain{3,4}}
\end{tabproofwide}


% ------------------------------------------------------------
% (4) Rückwärts-Variante für drei Komponenten
% ------------------------------------------------------------
\FormulaThmDelta%
{x\in A \vdash F(G(H(x)))=((F\circ G)\circ H)(x)}%
{
  \DeltaRow{Mengen}{A\dsep B\dsep C\dsep D\dsep x}
  \DeltaPrem{Funktionen}{H\colon A\to B\dsep G\colon B\to C\dsep F\colon C\to D}
}
\begin{tabproof}
  \proofstep{1}{x\in A}{\rA}
  \proofstep{1}{((F\circ G)\circ H)(x)=F(G(H(x)))}{%
    \FormulaRefAuto{x\in A \vdash ((F\circ G)\circ H)(x)=F(G(H(x)))}{1}}
  \proofstep{1}{F(G(H(x)))=((F\circ G)\circ H)(x)}{%
    \FormulaRefAuto{a=b\vdash b=a}{2}}
\end{tabproof}


\subsubsection{Die Auswertungseigenschaft}

\FormulaThmDelta{%
x\in A\dsep y\in A\dsep G(F(x))=G(F(y))\vdash (G\circ F)(x)=(G\circ F)(y)
}{
  \DeltaRow{Mengen}{x \dsep A \dsep B \dsep C}
  \DeltaPrem{Funktionen}{F\colon A\to B \dsep G\colon B\to C}
}
\begin{tabproof}
\proofstep{1}{x\in A}{\rA}
\proofstep{2}{y\in A}{\rA}
\proofstep{3}{G(F(x))=G(F(y))}{\rA}
\proofstep{1}{(G\circ F)(x)=G(F(x))}{\FormulaRefAuto{x\in A \vdash (G\circ F)(x)=G(F(x))}{1}}
\proofstep{2}{(G\circ F)(y)=G(F(y))}{\FormulaRefAuto{x\in A \vdash (G\circ F)(x)=G(F(x))}{2}}
\proofstep{3}{(G\circ F)(x)=G(F(y))}{\rIE{4,3}}
\proofstep{3}{(G\circ F)(x)=(G\circ F)(y)}{\rIE{5,6}}
\end{tabproof}

\FormulaThmDelta{%
x\in A\dsep y\in A\dsep (G\circ F)(x)=(G\circ F)(y) \vdash G(F(x))=G(F(y))
}{
  \DeltaRow{Mengen}{x \dsep A \dsep B \dsep C}
  \DeltaPrem{Funktionen}{F\colon A\to B \dsep G\colon B\to C}
}
\begin{tabproof}
  \proofstep{1}{x\in A}{\rA}
  \proofstep{2}{y\in A}{\rA}
  \proofstep{3}{(G\circ F)(x)=(G\circ F)(y)}{\rA}

  \proofstep{1}{(G\circ F)(x)=G(F(x))}{\FormulaRefAuto{x\in A \vdash (G\circ F)(x)=G(F(x))}{1}}
  \proofstep{2}{(G\circ F)(y)=G(F(y))}{\FormulaRefAuto{x\in A \vdash (G\circ F)(x)=G(F(x))}{2}}

  \proofstep{3}{G(F(x))=(G\circ F)(y)}{\rIE{4,3}}
  \proofstep{3}{G(F(x))=G(F(y))}{\rIE{5,6}}
\end{tabproof}

\FormulaThmDeltaR[Ungleichheit auf Kompositionswerte heben]{%
  \begin{aligned}[t]
    &x\in A\dsep y\in A\dsep G(F(x))\neq G(F(y))
    \vdash{}\\
    &(G\circ F)(x)\neq (G\circ F)(y)
  \end{aligned}
}{%
  x\in A\dsep y\in A\dsep G(F(x))\neq G(F(y))
  \vdash (G\circ F)(x)\neq (G\circ F)(y)
}{%
  \DeltaDecl{Mengen}{A\dsep x\dsep y}%
  \DeltaDecl{Funktionen}{F\dsep G}%
}
\begin{tabproof}
  \proofstep{1}{x\in A}{\rA}
  \proofstep{2}{y\in A}{\rA}
  \proofstep{3}{G(F(x))\neq G(F(y))}{\rA}
  \proofstep{4}{(G\circ F)(x)=(G\circ F)(y)}{\rA}
  \proofstep{1,2,4}{G(F(x))=G(F(y))}{%
    \FormulaRefAuto{x\in A\dsep y\in A\dsep (G\circ F)(x)=(G\circ F)(y) \vdash G(F(x))=G(F(y))}{1,2,4}}
  \proofstep{1,2,3,4}{\bot}{\rBI{3,5}}
  \proofstep{1,2,3}{(G\circ F)(x)\neq (G\circ F)(y)}{\rCI{4,6}}
\end{tabproof}


\subsubsection{Rechenregeln}

\FormulaThmDelta{%
G = H \vdash F\circ G = F\circ H
}{
  \DeltaRow{Mengen}{A \dsep B \dsep C \dsep D}
  \DeltaPrem{Funktionen}{G,H\colon A\to B \dsep F\colon B\to C}
}
\begin{tabproof}
  \proofstep{1}{G=H}{\rA}
  \proofstep{}{F\circ G=F\circ G}{\rII}
  \proofstep{}{F\circ G=F\circ H}{\rIE{1,2}}
\end{tabproof}

\FormulaThmDelta{%
G = H \vdash G\circ F = H\circ F
}{
  \DeltaRow{Mengen}{x \dsep A \dsep B \dsep C}
  \DeltaPrem{Funktionen}{F\colon A\to B \dsep G,H\colon B\to C}
}
\begin{tabproof}
  \proofstep{1}{G=H}{\rA}
  \proofstep{}{G\circ F=G\circ F}{\rII}
  \proofstep{}{G\circ F=H\circ F}{\rIE{1,2}}
\end{tabproof}

% — Assoziativität der Komposition —
\FormulaThmDelta[Assoziativität]{%
F\colon A\to B\dsep G\colon B\to C\dsep H\colon C\to D
\vdash H\circ (G\circ F) = (H\circ G)\circ F
}{
  \DeltaRow{Mengen}{A \dsep B \dsep C \dsep D}
  \DeltaRow{Funktionen}{F\dsep G\dsep H}
}
\begin{tabproofwide}
  \proofstepwidestar[1]{F\colon A\to B}{\rA}
  \proofstepwidestar[2]{G\colon B\to C}{\rA}
  \proofstepwidestar[3]{H\colon C\to D}{\rA}
  \proofstepwidestar[1,2]{G\circ F\colon A\to C}{%
    \FormulaRefAuto{CompositionDef}[def]{1,2}}
  \proofstepwidestar[1,2,3]{H\circ(G\circ F)\colon A\to D}{%
    \FormulaRefAuto{CompositionDef}[def]{4,3}}
  \proofstepwidestar[2,3]{H\circ G\colon B\to D}{%
    \FormulaRefAuto{CompositionDef}[def]{2,3}}
  \proofstepwidestar[1,2,3]{(H\circ G)\circ F\colon A\to D}{%
    \FormulaRefAuto{CompositionDef}[def]{1,6}}
  \proofstepwidestar[8]{x\in A}{\rA}
  \proofstepwidestar[1,8]{F(x)\in B}{\FormulaRefAuto{F\colon A\to B\dsep x\in A\vdash F(x)\in B}{1,8}}
  \proofstepwidestar[8]{(G\circ F)(x)=G(F(x))}{\FormulaRefAuto{x\in A \vdash (G\circ F)(x)=G(F(x))}{8}}
  \proofstepwide[8]{(H\circ(G\circ F))(x)}{=}{H((G\circ F)(x))}{\FormulaRefAuto{x\in A \vdash (G\circ F)(x)=G(F(x))}{8}}
  \proofstepwide[8]{}{=}{H(G(F(x)))}{\rIE{10,11}}
  \proofstepwide[8]{}{=}{(H\circ G)(F(x))}{\FormulaRefAuto{x\in A \vdash (G\circ F)(x)=G(F(x))}{9}}
  \proofstepwide[8]{}{=}{((H\circ G)\circ F)(x)}{\FormulaRefAuto{x\in A \vdash (G\circ F)(x)=G(F(x))}{8}}
  \proofstepwide[8]{(H\circ(G\circ F))(x)}{=}{((H\circ G)\circ F)(x)}{\rChain{11,14}}
  \proofstepwide[1,2,3]{\forall x\in A\,((H\circ(G\circ F))(x)}{=}{((H\circ G)\circ F)(x))}{\rUI{\rRI{8,15}}}
  \proofstepwide[1,2,3]{H\circ (G\circ F)}{=}{(H\circ G)\circ F}{%
    \FormulaRefAuto{FunctionExtensionality}[thm]{5,7,16}}
\end{tabproofwide}

\subsubsection{Bildmengen unter einer Linksinversen}

\FormulaThmDeltaK[Bildmengen unter einer Linksinversen]{%
  \begin{aligned}[t]
    &F\colon A\to B\dsep G\colon B\to A\dsep
      \forall x\in A\,G(F(x))=x\dsep K\subseteq A\vdash{}\\[-2pt]
    &G[F[K]]=K
  \end{aligned}%
}{LeftInverseImages}{%
  \DeltaRow{Mengen}{A\dsep B\dsep K\dsep u\dsep v\dsep x}%
  \DeltaRow{Funktionen}{F\dsep G}%
}
\begin{tabproofsplitwide}
  \proofpartwideindR[Erste Inklusion]{%
    F\colon A\to B\dsep G\colon B\to A\dsep
    \forall x\in A\,G(F(x))=x\dsep K\subseteq A
    \vdash G[F[K]]\subseteq K%
  }{%
    F\colon A\to B\dsep G\colon B\to A\dsep
    \forall x\in A\,G(F(x))=x\dsep K\subseteq A
    \vdash G[F[K]]\subseteq K%
  }[LeftInverseImagesForward]
    \proofstepwidestar[1]{F\colon A\to B}{\rA}
    \proofstepwidestar[2]{G\colon B\to A}{\rA}
    \proofstepwidestar[3]{\forall x\in A\,G(F(x))=x}{\rA}
    \proofstepwidestar[4]{K\subseteq A}{\rA}
    \proofstepwidestar[5]{u\in G[F[K]]}{\rA}
    \proofstepwidestar[1,4]{F[K]\subseteq B}{%
      \FormulaRefAuto{C\subseteq A\dsep F\colon A\to B
        \vdash F[C]\subseteq B}{4,1}}
    \proofstepwidestar[1,2,4,5]{\exists v\in F[K]\,(u=G(v))}{%
      \FormulaRefAuto{C\subseteq A\dsep F\colon A\to B\dsep
        y\in F[C]\vdash\exists x\in C\,y=F(x)}{6,2,5}}
    \proofstepwidestar[8]{v\in F[K]\land u=G(v)}{\rA}
    \proofstepwidestar[1,4,8]{\exists x\in K\,(v=F(x))}{%
      \FormulaRefAuto{C\subseteq A\dsep F\colon A\to B\dsep
        y\in F[C]\vdash\exists x\in C\,y=F(x)}{4,1,\rAEa{8}}}
    \proofstepwidestar[10]{x\in K\land v=F(x)}{\rA}
    \proofstepwidestar[4,10]{x\in A}{%
      \FormulaRefAuto{A\subseteq B\dsep x\in A\vdash x\in B}{4,\rAEa{10}}}
    \proofstepwidestar[3,4,10]{G(F(x))=x}{%
      \rRE{\rUE{3},11}}
    \proofstepwidestar[1,2,4,8,10]{u=G(F(x))}{%
      \rChain{\rAEb{8},\rLRS{\rAEb{10}}}}
    \proofstepwidestar[1,2,3,4,8,10]{u=x}{\rChain{13,12}}
    \proofstepwidestar[1,2,3,4,8,10]{x=u}{%
      \FormulaRefAuto{a=b\vdash b=a}{14}}
    \proofstepwidestar[1,2,3,4,8,10]{u\in K}{\rIE{15,\rAEa{10}}}
    \proofstepwidestar[1,2,3,4,8]{u\in K}{\rEE{9,10,16}}
    \proofstepwidestar[1,2,3,4,5]{u\in K}{\rEE{7,8,17}}
    \proofstepwidestar[1,2,3,4]{G[F[K]]\subseteq K}{%
      \FormulaRefAuto{A\subseteq B\coloneq
        \forall x\,(x\in A\rightarrow x\in B)}{\rUI{\rRI{5,18}}}}
  \closeproofpartwideindR

  \proofpartwideindR[Zweite Inklusion]{%
    F\colon A\to B\dsep G\colon B\to A\dsep
    \forall x\in A\,G(F(x))=x\dsep K\subseteq A
    \vdash K\subseteq G[F[K]]%
  }{%
    F\colon A\to B\dsep G\colon B\to A\dsep
    \forall x\in A\,G(F(x))=x\dsep K\subseteq A
    \vdash K\subseteq G[F[K]]%
  }[LeftInverseImagesBackward]
    \proofstepwidestar[1]{F\colon A\to B}{\rA}
    \proofstepwidestar[2]{G\colon B\to A}{\rA}
    \proofstepwidestar[3]{\forall x\in A\,G(F(x))=x}{\rA}
    \proofstepwidestar[4]{K\subseteq A}{\rA}
    \proofstepwidestar[5]{u\in K}{\rA}
    \proofstepwidestar[4,5]{u\in A}{%
      \FormulaRefAuto{A\subseteq B\dsep x\in A\vdash x\in B}{4,5}}
    \proofstepwidestar[1,4,5]{F(u)\in F[K]}{%
      \FormulaRefAuto{C\subseteq A\dsep F\colon A\to B\dsep
        x\in C\vdash F(x)\in F[C]}{4,1,5}}
    \proofstepwidestar[1,4]{F[K]\subseteq B}{%
      \FormulaRefAuto{C\subseteq A\dsep F\colon A\to B
        \vdash F[C]\subseteq B}{4,1}}
    \proofstepwidestar[1,2,4,5]{G(F(u))\in G[F[K]]}{%
      \FormulaRefAuto{C\subseteq A\dsep F\colon A\to B\dsep
        x\in C\vdash F(x)\in F[C]}{8,2,7}}
    \proofstepwidestar[3,4,5]{G(F(u))=u}{%
      \rRE{\rUE{3},6}}
    \proofstepwidestar[1,2,3,4,5]{u\in G[F[K]]}{\rIE{10,9}}
    \proofstepwidestar[1,2,3,4]{K\subseteq G[F[K]]}{%
      \FormulaRefAuto{A\subseteq B\coloneq
        \forall x\,(x\in A\rightarrow x\in B)}{\rUI{\rRI{5,11}}}}
  \closeproofpartwideindR

  \proofpartwide{Schluss}
    \proofstepwidestar[1]{F\colon A\to B}{\rA}
    \proofstepwidestar[2]{G\colon B\to A}{\rA}
    \proofstepwidestar[3]{\forall x\in A\,G(F(x))=x}{\rA}
    \proofstepwidestar[4]{K\subseteq A}{\rA}
    \proofstepwidestar[1,2,3,4]{G[F[K]]\subseteq K}{%
      \FormulaRefAuto{LeftInverseImagesForward}{1,2,3,4}}
    \proofstepwidestar[1,2,3,4]{K\subseteq G[F[K]]}{%
      \FormulaRefAuto{LeftInverseImagesBackward}{1,2,3,4}}
    \proofstepwidestar[1,2,3,4]{G[F[K]]=K}{%
      \FormulaRefAuto{A\subseteq B\dsep B\subseteq A\vdash A=B}{5,6}}
  \closeproofpartwide
\end{tabproofsplitwide}


\subsection{Die induzierte Potenzmengenabbildung}

\FormulaDefDeltaK[Induzierte Potenzmengenabbildung]{%
  F\colon A\to B\vdash
  \begin{aligned}[t]
    &\PowMap{F}\colon\powerset(A)\to\powerset(B),\\[-2pt]
    &\forall X\in\powerset(A)\,
      \bigl(\PowMap{F}(X)\coloneqq F[X]\bigr)
  \end{aligned}
}{PowerMapDef}{%
  \DeltaRow{Mengen}{A\dsep B\dsep X}
  \DeltaRow{Funktionen}{F\dsep \PowMap{F}}
}
\begin{tabproof}
  \proofstep{1}{F\colon A\to B}{\rA}
  \proofstep{2}{X\in\powerset(A)}{\rA}
  \proofstep{2}{X\subseteq A}{%
    \FormulaRefAuto{x\in\powerset(A)\eqvdash x\subseteq A}{2}}
  \proofstep{1,2}{F[X]\subseteq B}{%
    \FormulaRefAuto{C\subseteq A\dsep F\colon A\to B\vdash F[C]\subseteq B}{3,1}}
  \proofstep{1,2}{F[X]\in\powerset(B)}{%
    \FormulaRefAuto{x\in\powerset(A)\eqvdash x\subseteq A}{4}}
\end{tabproof}

\FormulaThmDelta[Typisierung der induzierten Potenzmengenabbildung]{%
  F\colon A\to B \vdash \PowMap{F}\colon \powerset(A)\to \powerset(B)
}{%
  \DeltaRow{Mengen}{A\dsep B\dsep X}
  \DeltaRow{Funktionen}{F\dsep \PowMap{F}}
}
\begin{tabproof}
  \proofstep{1}{F\colon A\to B}{\rA}
  \proofstep{1}{\PowMap{F}\colon \powerset(A)\to \powerset(B)}{%
    \FormulaRefAuto{PowerMapDef}[def]{1}}
\end{tabproof}

\FormulaThmDelta[Explizite Werte der induzierten Potenzmengenabbildung]{%
  F\colon A\to B \dsep X\in\powerset(A)\vdash \PowMap{F}(X)=F[X]
}{%
  \DeltaRow{Mengen}{A\dsep B\dsep X}
  \DeltaRow{Funktionen}{F\dsep \PowMap{F}}
}
\begin{tabproof}
  \proofstep{1}{F\colon A\to B}{\rA}
  \proofstep{2}{X\in\powerset(A)}{\rA}
  \proofstep{1}{\forall Y\in\powerset(A)\,\PowMap{F}(Y)=F[Y]}{%
    \FormulaRefAuto{PowerMapDef}[def]{1}}
  \proofstep{1,2}{\PowMap{F}(X)=F[X]}{%
    \rRE{\rUE{3},2}}
\end{tabproof}


\FormulaDefDeltaK[Eingeschränkte Potenzmengenabbildung]{%
  f\colon A\to B\vdash
  \PowMapNE{f}\coloneqq
  \PowMap{f}\restriction_{\PowNE{A}}%
}{NonemptyPowerMapDef}{%
  \DeltaRow{Mengen}{A\dsep B}
  \DeltaRow{Funktionen}{f\dsep\PowMap{f}}
  \DeltaRow{\textbf{Neues Symbol}}{}
  \DeltaPrem{Nichtleere Potenzmengenabbildung}{\PowMapNE{f}}
}

\FormulaThmDeltaK[Werte der eingeschränkten Potenzmengenabbildung]{%
  \begin{aligned}[t]
    &f\colon A\to B\dsep X\in\PowNE{A}\vdash{}\\[-2pt]
    &\PowMapNE{f}(X)=f[X]
  \end{aligned}%
}{NonemptyPowerMapValue}{%
  \DeltaRow{Mengen}{A\dsep B\dsep X}
  \DeltaRow{Funktionen}{f\dsep\PowMap{f}\dsep\PowMapNE{f}}
}
\begin{tabproofwide}
  \proofstepwidestar[1]{f\colon A\to B}{\rA}
  \proofstepwidestar[2]{X\in\PowNE{A}}{\rA}
  \proofstepwidestar[]{\PowNE{A}\subseteq\powerset(A)}{%
    \FormulaRefAuto{A\setminus B\subseteq A}[thm]}
  \proofstepwidestar[2]{X\in\powerset(A)}{%
    \FormulaRefAuto{A\subseteq B,\,x\in A\vdash x\in B}[thm]{3,2}}
  \proofstepwidestar[1]{\PowMap{f}\colon\powerset(A)\to\powerset(B)}{%
    \FormulaRefAuto{F\colon A\to B\vdash
      \PowMap{F}\colon\powerset(A)\to\powerset(B)}[thm]{1}}
  \proofstepwidestar[1]{%
    \PowMapNE{f}=\PowMap{f}\restriction_{\PowNE{A}}}{%
    \FormulaRefAuto{NonemptyPowerMapDef}[def]{1}}
  \proofstepwidestar[]{%
    \bigl(\PowMap{f}\restriction_{\PowNE{A}}\bigr)(X)
    =\bigl(\PowMap{f}\restriction_{\PowNE{A}}\bigr)(X)}{\rII}
  \proofstepwide[1,2]{%
    \PowMapNE{f}(X)}{=}{%
    \bigl(\PowMap{f}\restriction_{\PowNE{A}}\bigr)(X)}{\rIE{6,7}}
  \proofstepwide[1,2]{}{=}{\PowMap{f}(X)}{%
    \FormulaRefAuto{C\subseteq A\dsep x\in C\vdash
      F\restriction_C(x)=F(x)}[thm]{3,2}}
  \proofstepwide[1,2]{}{=}{f[X]}{%
    \FormulaRefAuto{F\colon A\to B\dsep X\in\powerset(A)
      \vdash\PowMap{F}(X)=F[X]}[thm]{1,4}}
  \proofstepwide[1,2]{\PowMapNE{f}(X)}{=}{f[X]}{\rChain{8,10}}
\end{tabproofwide}

\FormulaThmDeltaK[Typisierung der eingeschränkten Potenzmengenabbildung]{%
  f\colon A\to B\vdash
  \PowMapNE{f}\colon\PowNE{A}\to\PowNE{B}%
}{NonemptyPowerMapTyped}{%
  \DeltaRow{Mengen}{A\dsep B\dsep X}
  \DeltaRow{Funktionen}{f\dsep\PowMap{f}\dsep\PowMapNE{f}}
}
\begin{tabproofwide}
  \proofstepwidestar[1]{f\colon A\to B}{\rA}
  \proofstepwidestar[]{\PowNE{A}\subseteq\powerset(A)}{%
    \FormulaRefAuto{A\setminus B\subseteq A}[thm]}
  \proofstepwidestar[1]{\PowMap{f}\colon\powerset(A)\to\powerset(B)}{%
    \FormulaRefAuto{F\colon A\to B\vdash
      \PowMap{F}\colon\powerset(A)\to\powerset(B)}[thm]{1}}
  \proofstepwidestar[1]{%
    \PowMap{f}\restriction_{\PowNE{A}}
      \colon\PowNE{A}\to\powerset(B)}{%
    \FormulaRefAuto{C\subseteq A\vdash
      F\restriction_C\colon C\to B}[thm]{2}}
  \proofstepwidestar[1]{%
    \PowMapNE{f}=\PowMap{f}\restriction_{\PowNE{A}}}{%
    \FormulaRefAuto{NonemptyPowerMapDef}[def]{1}}
  \proofstepwidestar[1]{%
    \PowMapNE{f}\colon\PowNE{A}\to\powerset(B)}{\rIE{5,4}}
  \proofstepwidestar[7]{X\in\PowNE{A}}{\rA}
  \proofstepwidestar[1,7]{f[X]\in\PowNE{B}}{%
    \FormulaRefAuto{NonemptyDirectImage}[thm]{1,7}}
  \proofstepwidestar[1,7]{\PowMapNE{f}(X)=f[X]}{%
    \FormulaRefAuto{NonemptyPowerMapValue}[thm]{1,7}}
  \proofstepwidestar[1,7]{\PowMapNE{f}(X)\in\PowNE{B}}{\rIE{9,8}}
  \proofstepwidestar[1]{%
    \forall X\in\PowNE{A}\,\PowMapNE{f}(X)\in\PowNE{B}}{%
    \rUI{\rRI{7,10}}}
  \proofstepwidestar[1]{%
    \PowMapNE{f}\colon\PowNE{A}\to\PowNE{B}}{%
    \FormulaRefAuto{F\colon A\to B\dsep
      \forall x\in A\,F(x)\in C\vdash F\colon A\to C}[thm]{6,11}}
\end{tabproofwide}

\FormulaThmDeltaK[Urbild der nichtleeren Potenzmenge]{%
  f\colon A\to B\vdash
  (\PowMap{f})^{-1}[\PowNE{B}]=\PowNE{A}%
}{NonemptyPowerMapPreimage}{%
  \DeltaRow{Mengen}{A\dsep B\dsep X}
  \DeltaRow{Funktionen}{f\dsep\PowMap{f}}
}
\begin{tabproofsplitwide}
  \proofpartwideindR[Urbildpunkte sind nichtleer]{%
    \begin{aligned}[t]
      &f\colon A\to B\dsep
        X\in(\PowMap{f})^{-1}[\PowNE{B}]\vdash{}\\[-2pt]
      &X\neq\varnothing
    \end{aligned}
  }{%
    \begin{aligned}[t]
      &f\colon A\to B\dsep
        X\in(\PowMap{f})^{-1}[\PowNE{B}]\vdash{}\\[-2pt]
      &X\neq\varnothing
    \end{aligned}
  }[NonemptyPowerMapPreimageNonempty]
    \proofstepwidestar[1]{f\colon A\to B}{\rA}
    \proofstepwidestar[2]{X\in(\PowMap{f})^{-1}[\PowNE{B}]}{\rA}
    \proofstepwidestar[1]{\PowMap{f}\colon\powerset(A)\to\powerset(B)}{%
      \FormulaRefAuto{F\colon A\to B\vdash
        \PowMap{F}\colon\powerset(A)\to\powerset(B)}[thm]{1}}
    \proofstepwidestar[]{\PowNE{B}\subseteq\powerset(B)}{%
      \FormulaRefAuto{A\setminus B\subseteq A}[thm]}
    \proofstepwidestar[1]{%
      \begin{aligned}[t]
        &X\in(\PowMap{f})^{-1}[\PowNE{B}]\\[-2pt]
        &\quad\leftrightarrow
          \bigl(X\in\powerset(A)\\[-2pt]
        &\hspace{30mm}\land\PowMap{f}(X)\in\PowNE{B}\bigr)
      \end{aligned}}{%
      \FormulaRefAuto{F\colon A\to B\dsep C\subseteq B
        \vdash\bigl(x\in F^{-1}[C]\leftrightarrow
        (x\in A\land F(x)\in C)\bigr)}[thm]{3,4}}
    \proofstepwidestar[1,2]{%
      X\in\powerset(A)\land\PowMap{f}(X)\in\PowNE{B}}{%
      \FormulaRefAuto{P\leftrightarrow Q, P\vdash Q}[thm]{5,2}}
    \proofstepwidestar[1,2]{X\in\powerset(A)}{\rAEa{6}}
    \proofstepwidestar[1,2]{\PowMap{f}(X)\in\PowNE{B}}{\rAEb{6}}
    \proofstepwidestar[1,2]{\PowMap{f}(X)\neq\varnothing}{%
      \rAEb{\FormulaRefAuto{NonemptyPowersetMembership}[thm]{8}}}
    \proofstepwidestar[10]{X=\varnothing}{\rA}
    \proofstepwide[1,2]{\PowMap{f}(X)}{=}{f[X]}{%
      \FormulaRefAuto{F\colon A\to B\dsep X\in\powerset(A)
        \vdash\PowMap{F}(X)=F[X]}[thm]{1,7}}
    \proofstepwide[1,2,10]{}{=}{f[\varnothing]}{\rLRS{10}}
    \proofstepwide[1,2,10]{}{=}{\varnothing}{%
      \FormulaRefAuto{F\colon A\to B\vdash
        F[\varnothing]=\varnothing}[thm]{1}}
    \proofstepwide[1,2,10]{\PowMap{f}(X)}{=}{\varnothing}{%
      \rChain{11,13}}
    \proofstepwidestar[1,2,10]{\bot}{\rBI{9,14}}
    \proofstepwidestar[1,2]{X\neq\varnothing}{\rCI{10,15}}
  \closeproofpartwideindR

  \proofpartwideindR[Urbildpunkte gehören zur nichtleeren Potenzmenge]{%
    \begin{aligned}[t]
      &f\colon A\to B\dsep
        X\in(\PowMap{f})^{-1}[\PowNE{B}]\vdash{}\\[-2pt]
      &X\in\PowNE{A}
    \end{aligned}
  }{%
    \begin{aligned}[t]
      &f\colon A\to B\dsep
        X\in(\PowMap{f})^{-1}[\PowNE{B}]\vdash{}\\[-2pt]
      &X\in\PowNE{A}
    \end{aligned}
  }[NonemptyPowerMapPreimageElement]
    \proofstepwidestar[1]{f\colon A\to B}{\rA}
    \proofstepwidestar[2]{X\in(\PowMap{f})^{-1}[\PowNE{B}]}{\rA}
    \proofstepwidestar[1]{\PowMap{f}\colon\powerset(A)\to\powerset(B)}{%
      \FormulaRefAuto{F\colon A\to B\vdash
        \PowMap{F}\colon\powerset(A)\to\powerset(B)}[thm]{1}}
    \proofstepwidestar[]{\PowNE{B}\subseteq\powerset(B)}{%
      \FormulaRefAuto{A\setminus B\subseteq A}[thm]}
    \proofstepwidestar[1]{%
      \begin{aligned}[t]
        &X\in(\PowMap{f})^{-1}[\PowNE{B}]\\[-2pt]
        &\quad\leftrightarrow
          \bigl(X\in\powerset(A)\\[-2pt]
        &\hspace{30mm}\land\PowMap{f}(X)\in\PowNE{B}\bigr)
      \end{aligned}}{%
      \FormulaRefAuto{F\colon A\to B\dsep C\subseteq B
        \vdash\bigl(x\in F^{-1}[C]\leftrightarrow
        (x\in A\land F(x)\in C)\bigr)}[thm]{3,4}}
    \proofstepwidestar[1,2]{%
      X\in\powerset(A)\land\PowMap{f}(X)\in\PowNE{B}}{%
      \FormulaRefAuto{P\leftrightarrow Q, P\vdash Q}[thm]{5,2}}
    \proofstepwidestar[1,2]{X\in\powerset(A)}{\rAEa{6}}
    \proofstepwidestar[1,2]{X\neq\varnothing}{%
      \FormulaRefAuto{NonemptyPowerMapPreimageNonempty}[thm]{1,2}}
    \proofstepwidestar[1,2]{X\subseteq A}{%
      \FormulaRefAuto{x\in\powerset(A)\eqvdash x\subseteq A}[thm]{7}}
    \proofstepwidestar[1,2]{X\subseteq A\land X\neq\varnothing}{%
      \rAI{9,8}}
    \proofstepwidestar[1,2]{X\in\PowNE{A}}{%
      \FormulaRefAuto{NonemptyPowersetMembership}[thm]{10}}
  \closeproofpartwideindR

  \proofpartwideindR[Nichtleere Teilmengen liegen im Urbild]{%
    f\colon A\to B\dsep X\in\PowNE{A}
    \vdash X\in(\PowMap{f})^{-1}[\PowNE{B}]
  }{%
    f\colon A\to B\dsep X\in\PowNE{A}
    \vdash X\in(\PowMap{f})^{-1}[\PowNE{B}]
  }[NonemptyPowerMapPreimageReverseElement]
    \proofstepwidestar[1]{f\colon A\to B}{\rA}
    \proofstepwidestar[2]{X\in\PowNE{A}}{\rA}
    \proofstepwidestar[1]{\PowMap{f}\colon\powerset(A)\to\powerset(B)}{%
      \FormulaRefAuto{F\colon A\to B\vdash
        \PowMap{F}\colon\powerset(A)\to\powerset(B)}[thm]{1}}
    \proofstepwidestar[]{\PowNE{B}\subseteq\powerset(B)}{%
      \FormulaRefAuto{A\setminus B\subseteq A}[thm]}
    \proofstepwidestar[2]{X\subseteq A}{%
      \rAEa{\FormulaRefAuto{NonemptyPowersetMembership}[thm]{2}}}
    \proofstepwidestar[2]{X\in\powerset(A)}{%
      \FormulaRefAuto{x\in\powerset(A)\eqvdash x\subseteq A}[thm]{5}}
    \proofstepwidestar[1,2]{f[X]\in\PowNE{B}}{%
      \FormulaRefAuto{NonemptyDirectImage}[thm]{1,2}}
    \proofstepwidestar[1,2]{\PowMap{f}(X)=f[X]}{%
      \FormulaRefAuto{F\colon A\to B\dsep X\in\powerset(A)
        \vdash\PowMap{F}(X)=F[X]}[thm]{1,6}}
    \proofstepwidestar[1,2]{\PowMap{f}(X)\in\PowNE{B}}{\rIE{8,7}}
    \proofstepwidestar[1,2]{%
      X\in\powerset(A)\land\PowMap{f}(X)\in\PowNE{B}}{\rAI{6,9}}
    \proofstepwidestar[1]{%
      \begin{aligned}[t]
        &X\in(\PowMap{f})^{-1}[\PowNE{B}]\\[-2pt]
        &\quad\leftrightarrow
          \bigl(X\in\powerset(A)\\[-2pt]
        &\hspace{30mm}\land\PowMap{f}(X)\in\PowNE{B}\bigr)
      \end{aligned}}{%
      \FormulaRefAuto{F\colon A\to B\dsep C\subseteq B
        \vdash\bigl(x\in F^{-1}[C]\leftrightarrow
        (x\in A\land F(x)\in C)\bigr)}[thm]{3,4}}
    \proofstepwidestar[1,2]{X\in(\PowMap{f})^{-1}[\PowNE{B}]}{%
      \FormulaRefAuto{P\leftrightarrow Q, Q\vdash P}[thm]{11,10}}
  \closeproofpartwideindR

  \proofpartwideindR[Erste Teilmengenrichtung]{%
    f\colon A\to B\vdash
    (\PowMap{f})^{-1}[\PowNE{B}]\subseteq\PowNE{A}
  }{%
    f\colon A\to B\vdash
    (\PowMap{f})^{-1}[\PowNE{B}]\subseteq\PowNE{A}
  }[NonemptyPowerMapPreimageForwardSubset]
    \proofstepwidestar[1]{f\colon A\to B}{\rA}
    \proofstepwidestar[2]{X\in(\PowMap{f})^{-1}[\PowNE{B}]}{\rA}
    \proofstepwidestar[1,2]{X\in\PowNE{A}}{%
      \FormulaRefAuto{NonemptyPowerMapPreimageElement}[thm]{1,2}}
    \proofstepwidestar[1]{%
      \forall X\bigl(X\in(\PowMap{f})^{-1}[\PowNE{B}]
      \rightarrow X\in\PowNE{A}\bigr)}{\rUI{\rRI{2,3}}}
    \proofstepwidestar[1]{%
      (\PowMap{f})^{-1}[\PowNE{B}]\subseteq\PowNE{A}}{%
      \FormulaRefAuto{A\subseteq B\coloneq
        \forall x\,(x\in A\rightarrow x\in B)}[def]{4}}
  \closeproofpartwideindR

  \proofpartwideindR[Zweite Teilmengenrichtung]{%
    f\colon A\to B\vdash
    \PowNE{A}\subseteq(\PowMap{f})^{-1}[\PowNE{B}]
  }{%
    f\colon A\to B\vdash
    \PowNE{A}\subseteq(\PowMap{f})^{-1}[\PowNE{B}]
  }[NonemptyPowerMapPreimageReverseSubset]
    \proofstepwidestar[1]{f\colon A\to B}{\rA}
    \proofstepwidestar[2]{X\in\PowNE{A}}{\rA}
    \proofstepwidestar[1,2]{X\in(\PowMap{f})^{-1}[\PowNE{B}]}{%
      \FormulaRefAuto{NonemptyPowerMapPreimageReverseElement}[thm]{1,2}}
    \proofstepwidestar[1]{%
      \forall X\bigl(X\in\PowNE{A}\rightarrow
      X\in(\PowMap{f})^{-1}[\PowNE{B}]\bigr)}{\rUI{\rRI{2,3}}}
    \proofstepwidestar[1]{%
      \PowNE{A}\subseteq(\PowMap{f})^{-1}[\PowNE{B}]}{%
      \FormulaRefAuto{A\subseteq B\coloneq
        \forall x\,(x\in A\rightarrow x\in B)}[def]{4}}
  \closeproofpartwideindR

  \proofpartwide{Schluss}
    \proofstepwidestar[1]{f\colon A\to B}{\rA}
    \proofstepwidestar[1]{%
      (\PowMap{f})^{-1}[\PowNE{B}]\subseteq\PowNE{A}}{%
      \FormulaRefAuto{NonemptyPowerMapPreimageForwardSubset}[thm]{1}}
    \proofstepwidestar[1]{%
      \PowNE{A}\subseteq(\PowMap{f})^{-1}[\PowNE{B}]}{%
      \FormulaRefAuto{NonemptyPowerMapPreimageReverseSubset}[thm]{1}}
    \proofstepwidestar[1]{%
      (\PowMap{f})^{-1}[\PowNE{B}]=\PowNE{A}}{%
      \FormulaRefAuto{A\subseteq B, B\subseteq A\vdash A=B}[thm]{2,3}}
  \closeproofpartwide
\end{tabproofsplitwide}

\subsection{Spurabbildungen mengenwertiger Funktionen}

Ist \(F\colon A\to\powerset(B)\) eine mengenwertige Funktion, so
beschreibt die Spurabbildung zu einem Wert \(b\in B\) die Menge aller
Argumente \(x\in A\), deren Funktionswert \(b\) enthält.

\FormulaThmDeltaK[Spurmengen liegen in der Potenzmenge des Definitionsbereichs]{%
  F\colon A\to\powerset(B)\dsep b\in B
  \vdash
  \{x\in A\mid b\in F(x)\}\in\powerset(A)
}{TraceSetInPower}{%
  \DeltaRow{Mengen}{A\dsep B\dsep b\dsep x}
  \DeltaRow{Funktionen}{F}
}
\begin{tabproof}
  \proofstep{1}{F\colon A\to\powerset(B)}{\rA}
  \proofstep{2}{b\in B}{\rA}
  \proofstep{3}{x\in
    \{x\in A\mid b\in F(x)\}}{\rA}
  \proofstep{3}{x\in A}{%
    \FormulaRefAuto{x \in \{x \in A \mid P(x)\}\vdash x\in A}{3}}
  \proofstep{}{\{x\in A\mid b\in F(x)\}\subseteq A}{%
    \FormulaRefAuto{A \subseteq B \coloneq
    \forall x\,(x\in A \rightarrow x\in B)}{\rUI{\rRI{3,4}}}}
  \proofstep{}{\{x\in A\mid b\in F(x)\}\in\powerset(A)}{%
    \FormulaRefAuto{x \in \powerset(A)\eqvdash x \subseteq A}{5}}
\end{tabproof}

\FormulaDefDeltaK[Spurabbildung einer mengenwertigen Funktion]{%
  F\colon A\to\powerset(B)\vdash
  \begin{aligned}[t]
    &\TraceMap{F}\colon B\to\powerset(A),\\[-2pt]
    &\forall b\in B\,
      \bigl(
        \TraceMap{F}(b)\coloneqq\{x\in A\mid b\in F(x)\}
      \bigr)
  \end{aligned}
}{TraceMapDef}{%
  \DeltaRow{Mengen}{A\dsep B\dsep b}
  \DeltaRow{Funktionen}{F\dsep \TraceMap{F}}
  \DeltaRow{Neue Symbole}{\TraceMap{F}}
}
\begin{tabproof}
  \proofstep{1}{F\colon A\to\powerset(B)}{\rA}
  \proofstep{2}{b\in B}{\rA}
  \proofstep{1,2}{\{x\in A\mid b\in F(x)\}\in\powerset(A)}{%
    \FormulaRefAuto{TraceSetInPower}[thm]{1,2}}
\end{tabproof}

\FormulaThmDeltaK[Typisierung der Spurabbildung]{%
  F\colon A\to\powerset(B)\vdash
  \TraceMap{F}\colon B\to\powerset(A)
}{TraceMapFunction}{%
  \DeltaRow{Mengen}{A\dsep B}
  \DeltaRow{Funktionen}{F\dsep \TraceMap{F}}
}
\begin{tabproof}
  \proofstep{1}{F\colon A\to\powerset(B)}{\rA}
  \proofstep{1}{\TraceMap{F}\colon B\to\powerset(A)}{%
    \FormulaRefAuto{TraceMapDef}[def]{1}}
\end{tabproof}

\FormulaThmDeltaK[Wert der Spurabbildung]{%
  F\colon A\to\powerset(B)\dsep b\in B
  \vdash
  \TraceMap{F}(b)=\{x\in A\mid b\in F(x)\}
}{TraceMapValue}{%
  \DeltaRow{Mengen}{A\dsep B\dsep b\dsep x}
  \DeltaRow{Funktionen}{F\dsep \TraceMap{F}}
}
\begin{tabproof}
  \proofstep{1}{F\colon A\to\powerset(B)}{\rA}
  \proofstep{2}{b\in B}{\rA}
  \proofstep{1}{\forall c\in B\,
    \TraceMap{F}(c)=\{x\in A\mid c\in F(x)\}}{%
    \FormulaRefAuto{TraceMapDef}[def]{1}}
  \proofstep{1,2}{\TraceMap{F}(b)=
    \{x\in A\mid b\in F(x)\}}{%
    \rRE{\rUE{3},2}}
\end{tabproof}

\FormulaThmDeltaK[Gleiche Spuren liefern gleiche Bildmitgliedschaften]{%
  \begin{aligned}[t]
    &F\colon A\to\powerset(B)\dsep b\in B\dsep c\in B\\
    &\dsep \TraceMap{F}(b)=\TraceMap{F}(c)\dsep x\in A\\
    &\vdash b\in F(x)\leftrightarrow c\in F(x)
  \end{aligned}
}{TraceMapEqualityMembership}{%
  \DeltaRow{Mengen}{A\dsep B\dsep b\dsep c\dsep x}
  \DeltaRow{Funktionen}{F\dsep \TraceMap{F}}
}
\begin{tabproofwide}
  \proofstepwidestar[1]{F\colon A\to\powerset(B)}{\rA}
  \proofstepwidestar[2]{b\in B}{\rA}
  \proofstepwidestar[3]{c\in B}{\rA}
  \proofstepwidestar[4]{\TraceMap{F}(b)=\TraceMap{F}(c)}{\rA}
  \proofstepwidestar[5]{x\in A}{\rA}
  \proofstepwidestar[1,2]{\TraceMap{F}(b)=
    \{x\in A\mid b\in F(x)\}}{%
    \FormulaRefAuto{TraceMapValue}{1,2}}
  \proofstepwidestar[1,3]{\TraceMap{F}(c)=
    \{x\in A\mid c\in F(x)\}}{%
    \FormulaRefAuto{TraceMapValue}{1,3}}
  \proofstepwidestar[1,2]{%
    \{x\in A\mid b\in F(x)\}=\TraceMap{F}(b)}{%
    \FormulaRefAuto{a=b\vdash b=a}{6}}
  \proofstepwidestar[1,2,4]{%
    \{x\in A\mid b\in F(x)\}=\TraceMap{F}(c)}{%
    \FormulaRefAuto{a = b,\, b = c \vdash a = c}{8,4}}
  \proofstepwidestar[1,2,3,4]{%
    \{x\in A\mid b\in F(x)\}=
    \{x\in A\mid c\in F(x)\}}{%
    \FormulaRefAuto{a = b,\, b = c \vdash a = c}{9,7}}
  \proofstepwidestar[1,2,3,4,5]{%
    b\in F(x)\leftrightarrow c\in F(x)}{%
    \FormulaRefAuto{SeparationEqualityPredicateEquivalence}{10,5}}
\end{tabproofwide}



\section{Erweiterungen und leere Bereiche}

\subsection{Erweiterung einer Abbildung um einen Punkt}

\FormulaDefDeltaK[Erweiterungsabbildung]{%
  a\notin A\dsep f\colon A\to M\dsep b\in M\vdash
  \begin{aligned}[t]
    &\ExtMap{f}{a}{b}\colon(A\cup\{a\})\to M,\\[-2pt]
    &\forall x\in A\cup\{a\}\,
      \bigl(
        \ExtMap{f}{a}{b}(x)\coloneqq
        \IfThenElse{x=a}{b}{f(x)}
      \bigr)
  \end{aligned}%
}{ExtensionMapDef}{%
  \DeltaDecl{Mengen}{A\dsep a\dsep M\dsep b\dsep x}%
  \DeltaDecl{Funktionen}{f\dsep \ExtMap{f}{a}{b}}%
}
\begin{tabproof}
  \proofstep{1}{a\notin A}{\rA}
  \proofstep{2}{f\colon A\to M}{\rA}
  \proofstep{3}{b\in M}{\rA}
  \proofstep{4}{x\in A\cup\{a\}}{\rA}
  \proofstep{4}{x\in A\lor x=a}{%
    \FormulaRefAuto{x\in A\cup\{a\}\vdash x\in A\lor x=a}{4}}

  \proofstep{6}{x\in A}{\rA}
  \proofstep{1,6}{x\neq a}{%
    \FormulaRefAuto{x\in A\dsep a\notin A\vdash x\neq a}{6,1}}
  \proofstep{7}{\IfThenElse{x=a}{b}{f(x)}=f(x)}{%
    \FormulaRefAuto{\neg P\vdash\IfThenElse{P}{a}{b}=b}{7}}
  \proofstep{2,6}{f(x)\in M}{%
    \FormulaRefAuto{F\colon A\to B\dsep x\in A\vdash F(x)\in B}{2,6}}
  \proofstep{1,2,6}{\IfThenElse{x=a}{b}{f(x)}\in M}{\rIE{8,9}}

  \proofstep{11}{x=a}{\rA}
  \proofstep{11}{\IfThenElse{x=a}{b}{f(x)}=b}{%
    \FormulaRefAuto{P\vdash\IfThenElse{P}{a}{b}=a}{11}}
  \proofstep{3,11}{\IfThenElse{x=a}{b}{f(x)}\in M}{\rIE{12,3}}

  \proofstep{1,2,3,4}{\IfThenElse{x=a}{b}{f(x)}\in M}{%
    \rOE{5,6,10,11,13}}
\end{tabproof}

\FormulaThmDelta[Erweiterungsabbildung ist eine Funktion]{%
  a\notin A \dsep f\colon A\to M \dsep b\in M \vdash
  \ExtMap{f}{a}{b}\colon (A\cup\{a\})\to M%
}{%
  \DeltaDecl{Mengen}{A\dsep a\dsep M\dsep b}%
  \DeltaDecl{Funktionen}{f}%
  \DeltaDecl{Funktionensymbol}{\ExtMap{f}{a}{b}}%
}
\begin{tabproof}
  \proofstep{1}{a\notin A}{\rA}
  \proofstep{2}{f\colon A\to M}{\rA}
  \proofstep{3}{b\in M}{\rA}

  \proofstep{1,2,3}{\ExtMap{f}{a}{b}\colon (A\cup\{a\})\to M}{%
    \FormulaRefAuto{ExtensionMapDef}[def]{1,2,3}}
\end{tabproof}

\FormulaThmDelta[Erweiterung trifft den neuen Punkt]{%
  a\notin A \dsep f\colon A\to M \dsep b\in M \vdash \ExtMap{f}{a}{b}(a)=b%
}{%
  \DeltaDecl{Mengen}{A\dsep a\dsep M\dsep b}%
  \DeltaDecl{Funktionen}{f}%
  \DeltaDecl{Funktionensymbol}{\ExtMap{f}{a}{b}}%
}
\begin{tabproof}
  \proofstep{1}{a\notin A}{\rA}
  \proofstep{2}{f\colon A\to M}{\rA}
  \proofstep{3}{b\in M}{\rA}

  \proofstep{4}{a\in A\cup\{a\}}{%
    \FormulaRefAuto{a\in A\cup\{a\}}}
  \proofstep{1,2,3}{\forall x\in A\cup\{a\}\,
    \ExtMap{f}{a}{b}(x)=\IfThenElse{x=a}{b}{f(x)}}{%
    \FormulaRefAuto{ExtensionMapDef}[def]{1,2,3}}
  \proofstep{1,2,3,4}{\ExtMap{f}{a}{b}(a)=
    \IfThenElse{a=a}{b}{f(a)}}{%
    \rRE{\rUE{5},4}}
  \proofstep{}{ \IfThenElse{a=a}{b}{f(a)}=b }{%
    \FormulaRefAuto{P \vdash \IfThenElse{P}{a}{b}=a}{\rII}}
  \proofstep{1,2,3}{\ExtMap{f}{a}{b}(a)=b}{%
    \FormulaRefAuto{a=b,\,b=c\vdash a=c}{6,7}}
\end{tabproof}

\FormulaThmDelta[Erweiterung stimmt auf $A$ mit $f$ überein]{%
  a\notin A \dsep f\colon A\to M \dsep b\in M \vdash
  \forall x\in A\,\ExtMap{f}{a}{b}(x)=f(x)%
}{%
  \DeltaDecl{Mengen}{A\dsep a\dsep M\dsep b\dsep x}%
  \DeltaDecl{Funktionen}{f}%
  \DeltaDecl{Funktionensymbol}{\ExtMap{f}{a}{b}}%
}
\begin{tabproof}
  \proofstep{1}{a\notin A}{\rA}
  \proofstep{2}{f\colon A\to M}{\rA}
  \proofstep{3}{b\in M}{\rA}

  \proofstep{4}{x\in A}{\rA}
  \proofstep{4}{x\in A\cup\{a\}}{%
    \FormulaRefAuto{z\in A\vdash z\in A\cup B}{4}}
  \proofstep{1,2,3}{\forall y\in A\cup\{a\}\,
    \ExtMap{f}{a}{b}(y)=\IfThenElse{y=a}{b}{f(y)}}{%
    \FormulaRefAuto{ExtensionMapDef}[def]{1,2,3}}
  \proofstep{1,2,3,4}{\ExtMap{f}{a}{b}(x)=
    \IfThenElse{x=a}{b}{f(x)}}{%
    \rRE{\rUE{6},5}}
  \proofstep{1,4}{x\neq a}{%
    \FormulaRefAuto{x\in A \dsep a\notin A \vdash x\neq a}{4,1}}
  \proofstep{1,2,4}{\IfThenElse{x=a}{b}{f(x)}=f(x)}{%
    \FormulaRefAuto{\neg P \vdash \IfThenElse{P}{a}{b}=b}{8}}
  \proofstep{1,2,3,4}{\ExtMap{f}{a}{b}(x)=f(x)}{%
    \FormulaRefAuto{a=b,\,b=c\vdash a=c}{7,9}}
  \proofstep{1,2,3}{\forall x\in A\,\ExtMap{f}{a}{b}(x)=f(x)}{\rUI{\rRI{4,10}}}
\end{tabproof}

\FormulaThmDelta[Auswertungsformel der Erweiterung]{%
  a\notin A \dsep f\colon A\to M \dsep b\in M \dsep x\in A\cup\{a\}
  \vdash \ExtMap{f}{a}{b}(x)=\IfThenElse{x=a}{b}{f(x)}%
}{%
  \DeltaDecl{Mengen}{A\dsep a\dsep M\dsep b\dsep x}%
  \DeltaDecl{Funktionen}{f}%
  \DeltaDecl{Funktionensymbol}{\ExtMap{f}{a}{b}}%
}
\begin{tabproof}
  \proofstep{1}{a\notin A}{\rA}
  \proofstep{2}{f\colon A\to M}{\rA}
  \proofstep{3}{b\in M}{\rA}
  \proofstep{4}{x\in A\cup\{a\}}{\rA}

  \proofstep{1,2,3}{\forall y\in A\cup\{a\}\,
    \ExtMap{f}{a}{b}(y)=\IfThenElse{y=a}{b}{f(y)}}{%
    \FormulaRefAuto{ExtensionMapDef}[def]{1,2,3}}
  \proofstep{1,2,3,4}{\ExtMap{f}{a}{b}(x)=\IfThenElse{x=a}{b}{f(x)}}{%
    \rRE{\rUE{5},4}}
\end{tabproof}

\FormulaThmDelta[Punkt-Erweiterung]{%
  \begin{aligned}[t]
    &a\notin A \dsep f\colon A\to M \dsep b\in M\\
    &\vdash\ \exists g\colon (A\cup\{a\})\to M\,\Bigl(
      g(a)=b \land \forall x\in A\,g(x)=f(x)
    \Bigr)
  \end{aligned}%
}{%
  \DeltaDecl{Mengen}{A\dsep a\dsep M\dsep b\dsep x}%
  \DeltaDecl{Funktionen}{f\dsep g}%
  \DeltaDecl{Funktionensymbol}{\ExtMap{f}{a}{b}}%
}
\begin{tabproof}
  \proofstep{1}{a\notin A}{\rA}
  \proofstep{2}{f\colon A\to M}{\rA}
  \proofstep{3}{b\in M}{\rA}

  \proofstep{1,2,3}{\ExtMap{f}{a}{b}\colon (A\cup\{a\})\to M}{%
    \FormulaRefAuto{a\notin A \dsep f\colon A\to M \dsep b\in M \vdash
      \ExtMap{f}{a}{b}\colon (A\cup\{a\})\to M}{1,2,3}}

  \proofstep{1,2,3}{\ExtMap{f}{a}{b}(a)=b}{%
    \FormulaRefAuto{a\notin A \dsep f\colon A\to M \dsep b\in M \vdash
      \ExtMap{f}{a}{b}(a)=b}{1,2,3}}

  \proofstep{1,2,3}{\forall x\in A\,\ExtMap{f}{a}{b}(x)=f(x)}{%
    \FormulaRefAuto{a\notin A \dsep f\colon A\to M \dsep b\in M \vdash
      \forall x\in A\,\ExtMap{f}{a}{b}(x)=f(x)}{1,2,3}}

  \proofstep{1,2,3}{%
    \ExtMap{f}{a}{b}(a)=b \land \forall x\in A\,\ExtMap{f}{a}{b}(x)=f(x)}{%
    \rAI{5,6}}

\proofstep{1,2,3}{%
  \begin{aligned}[t]
    \exists g\colon (A\cup\{a\})\to M\,\Bigl(
      &g(a)=b \\
      &\land \forall x\in A\,g(x)=f(x)
    \Bigr)
  \end{aligned}%
}{\rEI{\rAI{4,7}}}
\end{tabproof}


\subsection{Abbildungen mit leerem Definitions- oder Zielbereich}

\FormulaThmDelta[Jede Funktion aus $\varnothing$ ist leer]{%
  f\colon \varnothing\to M \vdash f=\varnothing%
}{%
  \DeltaDecl{Mengen}{M\dsep x\dsep y\dsep u}%
  \DeltaDecl{Funktionen}{f}%
}
\begin{tabproof}
  \proofstep{1}{f\colon \varnothing\to M}{\rA}
  \proofstep{1}{f\subseteq \varnothing\times M}{\FormulaRefAuto{Funktion}{1}}
  \proofstep{3}{u\in f}{\rA}
  \proofstep{1,3}{u\in \varnothing\times M}{\FormulaRefAuto{A \subseteq B,\, x \in A \vdash x \in B}{2,3}}
  \proofstep{}{\varnothing\times M=\varnothing}{\FormulaRefAuto{\varnothing\times A=\varnothing}}
  \proofstep{1,3}{u\in\varnothing}{\rIE{5,4}}
  \proofstep{}{u\notin\varnothing}{\FormulaRefAuto{x \not\in \varnothing}}
  \proofstep{1,3}{\bot}{\rBI{6,7}}
  \proofstep{1}{u\notin f}{\rCI{2,8}}
  \proofstep{1}{\forall u\,u\notin f}{\rUI{9}}
  \proofstep{1}{f=\varnothing}{\FormulaRefAuto{\forall x\, x\notin A\vdash A=\varnothing}{10}}
\end{tabproof}

% =========================================================
% Abbildungen aus der leeren Menge
% =========================================================

\FormulaThmDelta[Leere Relation ist eine Funktion]{%
  \varnothing\colon \varnothing\to M%
}{%
  \DeltaDecl{Mengen}{M\dsep x\dsep y\dsep z}%
}
\begin{tabproof}
  \proofstep{}{\varnothing\subseteq\varnothing\times M}{%
    \FormulaRefAuto{\varnothing\subseteq A}}

  \proofstep{2}{%
    (x,y)\in\varnothing\land(x,z)\in\varnothing}{\rA}
  \proofstep{2}{(x,y)\in\varnothing}{\rAEa{2}}
  \proofstep{}{(x,y)\notin\varnothing}{%
    \FormulaRefAuto{x \not\in \varnothing}}
  \proofstep{2}{\bot}{\rBI{3,4}}
  \proofstep{2}{y=z}{%
    \FormulaRefAuto{P\dsep \neg P \vdash Q}{3,4}}
  \proofstep{}{%
    \bigl((x,y)\in\varnothing\land(x,z)\in\varnothing\bigr)\rightarrow y=z}{%
    \rRI{2,6}}

  \proofstep{8}{x\in\varnothing}{\rA}
  \proofstep{}{x\notin\varnothing}{%
    \FormulaRefAuto{x \not\in \varnothing}}
  \proofstep{8}{\bot}{\rBI{8,9}}
  \proofstep{8}{\exists y\,(x,y)\in\varnothing}{%
    \FormulaRefAuto{P\dsep \neg P \vdash Q}{8,9}}
  \proofstep{}{%
    x\in\varnothing\rightarrow\exists y\,(x,y)\in\varnothing}{%
    \rRI{8,11}}

  \proofstep{}{\varnothing\colon\varnothing\to M}{%
    \FormulaRefAuto{Funktion}{1,7,12}}
\end{tabproof}

\FormulaThmDelta[Eindeutige Existenz einer Funktion aus $\varnothing$]{%
  \exists! f\,\bigl(f\colon \varnothing\to M\bigr)%
}{%
  \DeltaDecl{Mengen}{M}%
  \DeltaDecl{Funktionen}{f\dsep g}%
}
\begin{tabproof}
  \proofstep{}{ \varnothing\colon \varnothing\to M }{%
    \FormulaRefAuto{\varnothing\colon \varnothing\to M}}

  \proofstep{}{ \exists f\,\bigl(f\colon \varnothing\to M\bigr) }{%
    \rEI{1}}

  \proofstep{3}{ f\colon \varnothing\to M }{\rA}
  \proofstep{4}{ g\colon \varnothing\to M }{\rA}

  \proofstep{3}{ f=\varnothing }{%
    \FormulaRefAuto{f\colon \varnothing\to M \vdash f=\varnothing}{3}}

  \proofstep{4}{ g=\varnothing }{%
    \FormulaRefAuto{f\colon \varnothing\to M \vdash f=\varnothing}{4}}

  \proofstep{3,4}{ f=g }{%
    \FormulaRefAuto{a = b,\, c = b \vdash a = c}{5,6}}

  \proofstep{}{ \exists! f\,\bigl(f\colon \varnothing\to M\bigr) }{%
    \rUEI{2,3,4,7}}
\end{tabproof}

\FormulaThmDelta{%
  A\neq \varnothing \vdash \neg \exists f\colon A\to \varnothing%
}{%
  \DeltaDecl{Mengen}{A\dsep x}%
  \DeltaDecl{Funktionen}{f}%
}
\begin{tabproof}
  \proofstep{1}{A\neq\varnothing}{\rA}
  \proofstep{1}{\exists x\,(x\in A)}{%
    \FormulaRefAuto{A\neq\varnothing \eqvdash \exists x\,(x \in A)}{1}}

  \proofstep{3}{\exists f\colon A\to\varnothing}{\rA}

  \proofstep{4}{x\in A}{\rA}
  \proofstep{5}{f\colon A\to\varnothing}{\rA}

  \proofstep{4,5}{(x,f(x))\in f}{%
    \FormulaRefAuto{F\colon A\to B\dsep x\in A\vdash (x,F(x))\in F}{5,4}}

  \proofstep{4,5}{f(x)\in\varnothing}{%
    \FormulaRefAuto{F\colon A\to B\dsep (x,y)\in F\vdash y\in B}{5,6}}

  \proofstep{}{f(x)\notin\varnothing}{%
    \FormulaRefAuto{x \not\in \varnothing}}

  \proofstep{4,5}{\bot}{\rBI{7,8}}

  \proofstep{1,5}{\bot}{\rEE{2,4,9}}
  \proofstep{1,3}{\bot}{\rEE{3,5,10}}

  \proofstep{1}{\neg \exists f\colon A\to\varnothing}{\rCI{3,11}}
\end{tabproof}



\section{Vereinigung einer überdeckenden Teilmengenfamilie}

\FormulaThmDeltaK[Vereinigung einer überdeckenden Teilmengenfamilie]{%
  D\colon I\to\powerset(W)\dsep
  \forall w\in W\,\exists i\in I\,(w\in D(i))
  \vdash\bigcup D[I]=W%
}{IndexedSubsetFamilyCoverUnion}{%
  \DeltaRow{Mengen}{I\dsep W\dsep E\dsep i\dsep w\dsep x}
  \DeltaRow{Funktionen}{D}
}
\begin{tabproofwide}
  \proofstepwidestar[1]{D\colon I\to\powerset(W)}{\rA}
  \proofstepwidestar[2]{%
    \forall w\in W\,\exists i\in I\,(w\in D(i))}{\rA}
  \proofstepwidestar[3]{x\in\bigcup D[I]}{\rA}
  \proofstepwidestar[3]{\exists E\in D[I]\,(x\in E)}{%
    \FormulaRefAuto{x\in\bigcup A\eqvdash
      \exists B\in A\,(x\in B)}[thm]{3}}
  \proofstepwidestar[5]{E\in D[I]\land x\in E}{\rA}
  \proofstepwidestar[1,5]{\exists i\in I\,E=D(i)}{%
    \FormulaRefAuto{F\colon A\to B\dsep y\in F[A]
      \vdash\exists x\in A\,y=F(x)}[thm]{1,\rAEa{5}}}
  \proofstepwidestar[7]{i\in I\land E=D(i)}{\rA}
  \proofstepwidestar[1,7]{D(i)\in\powerset(W)}{%
    \FormulaRefAuto{F\colon A\to B\dsep x\in A\vdash
      F(x)\in B}[thm]{1,\rAEa{7}}}
  \proofstepwidestar[1,7]{D(i)\subseteq W}{%
    \FormulaRefAuto{x\in\powerset(A)\eqvdash x\subseteq A}[thm]{8}}
  \proofstepwidestar[5,7]{x\in D(i)}{\rIE{\rAEb{7},\rAEb{5}}}
  \proofstepwidestar[1,5,7]{x\in W}{%
    \FormulaRefAuto{A\subseteq B\dsep x\in A\vdash x\in B}[thm]{9,10}}
  \proofstepwidestar[1,5]{x\in W}{\rEE{6,7,11}}
  \proofstepwidestar[1,3]{x\in W}{\rEE{4,5,12}}
  \proofstepwidestar[1]{\bigcup D[I]\subseteq W}{%
    \FormulaRefAuto{A\subseteq B\coloneq
      \forall x\,(x\in A\rightarrow x\in B)}[def]{\rUI{\rRI{3,13}}}}
  \proofstepwidestar[15]{x\in W}{\rA}
  \proofstepwidestar[2,15]{\exists i\in I\,(x\in D(i))}{%
    \rRE{\rUE{2},15}}
  \proofstepwidestar[17]{i\in I\land x\in D(i)}{\rA}
  \proofstepwidestar[]{I\subseteq I}{\FormulaRefAuto{A\subseteq A}[thm]}
  \proofstepwidestar[1,17]{D(i)\in D[I]}{%
    \FormulaRefAuto{C\subseteq A\dsep F\colon A\to B\dsep
      x\in C\vdash F(x)\in F[C]}[thm]{18,1,\rAEa{17}}}
  \proofstepwidestar[1,17]{D(i)\subseteq\bigcup D[I]}{%
    \FormulaRefAuto{B\in A\vdash B\subseteq\bigcup A}[thm]{19}}
  \proofstepwidestar[1,17]{x\in\bigcup D[I]}{%
    \FormulaRefAuto{A\subseteq B\dsep x\in A\vdash x\in B}[thm]{%
      20,\rAEb{17}}}
  \proofstepwidestar[1,2,15]{x\in\bigcup D[I]}{\rEE{16,17,21}}
  \proofstepwidestar[1,2]{W\subseteq\bigcup D[I]}{%
    \FormulaRefAuto{A\subseteq B\coloneq
      \forall x\,(x\in A\rightarrow x\in B)}[def]{\rUI{\rRI{15,22}}}}
  \proofstepwidestar[1,2]{\bigcup D[I]=W}{%
    \FormulaRefAuto{A\subseteq B\dsep B\subseteq A\vdash A=B}[thm]{14,23}}
\end{tabproofwide}

\chapter{Wegweiser zu späteren Funktionsbeispielen}

Die lineare Bandfolge bildet die logischen Abhängigkeiten ab, nicht die
ausschließliche thematische Zugehörigkeit. Die Eigenschaftsbände B06--B08
führen die allgemeinen Standardkonstruktionen mit kennzeichnender
Abbildungseigenschaft ein: Band~06 die Inklusionsabbildung und die injektive
Adjunktionsabbildung einer Mengenfamilie, Band~07 die Projektionen und
Band~08 Identitäts- und Umkehrfunktionen. Band~09 behandelt das
Auswahlprinzip.

In Band~10 werden die Nachfolger- und Vorgängerfunktion, rekursiv definierte
Abbildungen, Translationen, Folgenverschiebungen sowie
Stellenwertauswertungen untersucht. Band~17 ergänzt die Ganzzahleinbettung
sowie ganzzahlige Nachfolger- und Vorgängerabbildungen. Die Bände~18 und~19
konstruieren die Einbettungen der ganzen in die rationalen und der rationalen
in die reellen Zahlen. Band~20 enthält endliche Bijektionen und
Kardinalitätsvergleiche; Band~21 führt indizierte Familien und Folgen als
spezielle Funktionen ein. Die Bände~28--40 behandeln algebraische
Strukturen, Band~41 die Halbverbandsoperationen und Band~42 die für Frankls
Vermutung konstruierten
Normalisierungs-, Profil-, Komplement- und Quotientenabbildungen.

\end{document}
